% Options for packages loaded elsewhere
\PassOptionsToPackage{unicode}{hyperref}
\PassOptionsToPackage{hyphens}{url}
\PassOptionsToPackage{dvipsnames,svgnames,x11names}{xcolor}
%
\documentclass[
  10pt,
]{article}
\usepackage{amsmath,amssymb}
\usepackage{iftex}
\ifPDFTeX
  \usepackage[T1]{fontenc}
  \usepackage[utf8]{inputenc}
  \usepackage{textcomp} % provide euro and other symbols
\else % if luatex or xetex
  \usepackage{unicode-math} % this also loads fontspec
  \defaultfontfeatures{Scale=MatchLowercase}
  \defaultfontfeatures[\rmfamily]{Ligatures=TeX,Scale=1}
\fi
\usepackage{lmodern}
\ifPDFTeX\else
  % xetex/luatex font selection
  \setmainfont[]{DejaVu Serif}
  \setsansfont[]{DejaVu Sans}
  \setmonofont[]{DejaVu Sans Mono}
  \setmathfont[]{Latin Modern Math}
\fi
% Use upquote if available, for straight quotes in verbatim environments
\IfFileExists{upquote.sty}{\usepackage{upquote}}{}
\IfFileExists{microtype.sty}{% use microtype if available
  \usepackage[]{microtype}
  \UseMicrotypeSet[protrusion]{basicmath} % disable protrusion for tt fonts
}{}
\makeatletter
\@ifundefined{KOMAClassName}{% if non-KOMA class
  \IfFileExists{parskip.sty}{%
    \usepackage{parskip}
  }{% else
    \setlength{\parindent}{0pt}
    \setlength{\parskip}{6pt plus 2pt minus 1pt}}
}{% if KOMA class
  \KOMAoptions{parskip=half}}
\makeatother
\usepackage{xcolor}
\usepackage[margin=2.0cm]{geometry}
\usepackage{longtable,booktabs,array}
\usepackage{calc} % for calculating minipage widths
% Correct order of tables after \paragraph or \subparagraph
\usepackage{etoolbox}
\makeatletter
\patchcmd\longtable{\par}{\if@noskipsec\mbox{}\fi\par}{}{}
\makeatother
% Allow footnotes in longtable head/foot
\IfFileExists{footnotehyper.sty}{\usepackage{footnotehyper}}{\usepackage{footnote}}
\makesavenoteenv{longtable}
\setlength{\emergencystretch}{3em} % prevent overfull lines
\providecommand{\tightlist}{%
  \setlength{\itemsep}{0pt}\setlength{\parskip}{0pt}}
\setcounter{secnumdepth}{-\maxdimen} % remove section numbering
% doc_header.tex -- LaTeX preamble additions used by build_docs.sh for every
% Markdown -> LaTeX -> PDF conversion (README and provenance documents of Revision_SuperLoop).
\usepackage{etoolbox}
% longtable may not be loaded by pandoc when a document has no table
\usepackage{longtable}
% bold vectors: \bm -> \symbf under unicode-math (XeLaTeX), bm package otherwise
\makeatletter
\@ifpackageloaded{unicode-math}{\providecommand{\bm}[1]{\symbf{#1}}}{\usepackage{bm}}
\makeatother
% smaller type in long tables so that wide audit tables fit the text width;
% inside a table a line break is allowed after an underscore of plain text
% (pandoc writes it as \_), so that long identifiers wrap inside their column
% rows are set a little apart, so that symbols with indices do not touch
\AtBeginEnvironment{longtable}{\footnotesize\renewcommand{\arraystretch}{1.45}%
  \let\sptunderscore\_\def\_{\sptunderscore\allowbreak}}
\setlength{\LTcapwidth}{\textwidth}
% allow line breaks after slashes and hyphens inside table cells
\usepackage{xurl}
% code blocks: break over-long lines instead of running into the margin.
% Plain ``` blocks are typeset by pandoc as {verbatim}, highlighted ones
% (```bash) as {Highlighting}; both are re-declared as breakable fancyvrb
% environments.  The continuation line is marked with a small hooked arrow.
\usepackage{fvextra}
\RecustomVerbatimEnvironment{verbatim}{Verbatim}{breaklines,breakanywhere,fontsize=\small}
\makeatletter
\@ifundefined{Highlighting}{}{%
  \RecustomVerbatimEnvironment{Highlighting}{Verbatim}{breaklines,breakanywhere,fontsize=\small,commandchars=\\\{\}}}
\makeatother
% inline code (pandoc writes it as \texttt{...}): never hyphenate -- an inserted
% hyphen could be read as part of a file name -- but allow a line break after
% an underscore, slash, comma, colon or equals sign, so that long paths and
% column lists wrap inside the text width and inside table columns.
% Inside a formula (e.g. the name of a notebook symbol in an equation) the same
% command must first leave math mode, which \text of amsmath does.
\ExplSyntaxOn
\tl_new:N \l__spt_code_tl
\cs_new_protected:Npn \__spt_texttt:n #1
  {
    \tl_set:Nn \l__spt_code_tl {#1}
    \regex_replace_all:nnN { ([\/\,\:\=]) } { \1 \c{allowbreak} } \l__spt_code_tl
    \regex_replace_all:nnN { \c{_} } { \c{_} \c{allowbreak} } \l__spt_code_tl
    \group_begin: \ttfamily \hyphenchar \font = -1 \scan_stop: \l__spt_code_tl \group_end:
  }
\RenewDocumentCommand \texttt { m }
  {
    \mode_if_math:TF { \text { \__spt_texttt:n {#1} } } { \__spt_texttt:n {#1} }
  }
\ExplSyntaxOff
% A figure is written in the Markdown source as an image paragraph followed by
% a caption paragraph ("**Figure S1.** ..."), so that the caption is also
% visible on GitHub.  An image that starts a paragraph is centered.
\usepackage{graphicx}% pandoc loads it only for documents that contain an image
\NewCommandCopy\sptincludegraphics\includegraphics
\RenewDocumentCommand\includegraphics{O{} m}{%
  \ifvmode\begin{center}\sptincludegraphics[#1]{#2}\end{center}%
  \else\sptincludegraphics[#1]{#2}\fi}
% keep figures where they are referenced
\usepackage{float}
\floatplacement{figure}{H}
% running header with the document title
\usepackage{fancyhdr}
\pagestyle{fancy}
\fancyhf{}
\fancyhead[L]{\small PROJECT spacetimes, Revision\_SuperLoop}
\fancyhead[R]{\small\thepage}
\renewcommand{\headrulewidth}{0.4pt}
% 16x16 matrices: tighter columns so that a full matrix fits the text width
\setlength{\arraycolsep}{3.2pt}
% The documents of this project contain long tagged symbols, file paths and
% identifiers that cannot be hyphenated.  A paragraph is therefore allowed to
% have looser lines before TeX gives up and lets a line stick out of the margin.
\AtBeginDocument{\tolerance=2000 \emergencystretch=5em \hbadness=2000}
\ifLuaTeX
  \usepackage{selnolig}  % disable illegal ligatures
\fi
\IfFileExists{bookmark.sty}{\usepackage{bookmark}}{\usepackage{hyperref}}
\IfFileExists{xurl.sty}{\usepackage{xurl}}{} % add URL line breaks if available
\urlstyle{same}
\hypersetup{
  pdftitle={The task, the rules, what was done, what was not, and the decisions that are the author's},
  pdfauthor={Prepared for Patrick Lee Nash, Ph.D.},
  colorlinks=true,
  linkcolor={blue},
  filecolor={Maroon},
  citecolor={blue},
  urlcolor={blue},
  pdfcreator={LaTeX via pandoc}}

\title{The task, the rules, what was done, what was not, and the
decisions that are the author's}
\usepackage{etoolbox}
\makeatletter
\providecommand{\subtitle}[1]{% add subtitle to \maketitle
  \apptocmd{\@title}{\par {\large #1 \par}}{}{}
}
\makeatother
\subtitle{Provenance document \texttt{P0} of PROJECT spacetimes,
Revision\_SuperLoop}
\author{Prepared for Patrick Lee Nash, Ph.D.}
\date{November 16, 2026}

\begin{document}
\maketitle

\hypertarget{what-this-document-is}{%
\section{1. What this document is}\label{what-this-document-is}}

This document is the account of the work in the folder
\texttt{Revision\_SuperLoop}. It states

\begin{itemize}
\tightlist
\item
  the task, verbatim, and what it asks for; the later request for an
  explanation of one sentence of the manuscript; and the test of the
  sentence that follows it (Section 2);
\item
  each of the author's hard rules and how it is met (Section 3);
\item
  what the folder contains (Section 4) and what was done, in order
  (Section 5);
\item
  the results in brief (Section 6);
\item
  every \textbf{deviation} from the wording of the task, with its reason
  (Section 7) --- first of all why a counterexample to a proved theorem
  whose hypotheses hold cannot exist and what is provided instead, and
  why the requested ``super-group'' is a super-loop;
\item
  the \textbf{decisions that are the author's} to take or to confirm
  (Section 8);
\item
  what has \textbf{not} been done (Section 9);
\item
  the inputs (Section 10) and the \textbf{log of errors} made during the
  work (Section 11).
\end{itemize}

The three provenance documents of the folder are called \texttt{P0}
(this one), \texttt{P1} (the symbols) and \texttt{P2} (the
calculations), after the beginning of their file names. Every symbol
that occurs here is defined in Section 12, the complete table of the
symbols of the folder, and every formula that is used is written out
here. Nothing in this document has to be looked up elsewhere.

\hypertarget{the-task}{%
\section{2. The task}\label{the-task}}

\hypertarget{the-request-of-november-16-2026-1611-verbatim}{%
\subsection{2.1 The request of November 16, 2026, 16:11
(verbatim)}\label{the-request-of-november-16-2026-1611-verbatim}}

``Again Explore the so-called `no-go' theorems that forbid combining
Lorentz degrees of freedom with internal degrees of freedom, and look
for hypotheses that were used to prove these theorems such as
associativity that do not hold for the octonion algebra and the split
octonion algebra. Create a new `super-group' (that you either define or
invent) that combines support for SO(3,1) Lorentz degrees of freedom and
support for spinors with internal degrees of freedom and operators that
might possess eigenvalues of, say, electric charge, chirality, or
invented new entities. So Create exotic counter examples to these
`no-go' theorems that exist because hypotheses were used to prove these
`no-go' theorems such as associativity do not hold . Create counter
examples on any manifold that you can either find or invent, and that
supports your `super-group'. Look very hard and do not give up. Create
new folders for this effort, and files that contain a thorough and
correct analysis for any solution for a spinor superfluid boson singlet
(S=0) states, but employing your new `super-group'. Determine if these
theories are renormalizable. Document this solution, with provenance and
symbol files, and include in your manuscript. Push results to repo.''

\hypertarget{what-the-task-asks-for-item-by-item}{%
\subsection{2.2 What the task asks for, item by
item}\label{what-the-task-asks-for-item-by-item}}

\begin{enumerate}
\def\labelenumi{\arabic{enumi}.}
\tightlist
\item
  \textbf{The hypotheses of the no-go theorems} (McGlinn 1964,
  O'Raifeartaigh 1965, Coleman and Mandula 1967, with the extension of
  Pelc and Horwitz to any number of space dimensions, and the graded
  extension of Haag, Łopuszański and Sohnius) that do not hold for the
  octonions and the split octonions --- associativity first of all.
\item
  \textbf{A new ``super-group''} that contains the Lorentz group
  \(SO(3,1)\) (in the form of its double cover \(SL(2,\mathbb C)\)),
  spinors, and internal degrees of freedom, with operators whose
  eigenvalues are an electric charge, a chirality, or new quantities.
\item
  \textbf{Counterexamples to the no-go theorems} that exist because
  associativity fails, on manifolds that carry the ``super-group''.
\item
  \textbf{The spin-zero (``S=0'') boson-singlet superfluid of spinors},
  analysed thoroughly with the new structure.
\item
  \textbf{Renormalizability} of these theories.
\item
  Documentation with provenance and symbol files, a manuscript, and the
  push to the repository.
\end{enumerate}

The word ``Again'' refers to the folder \texttt{Revision\_NoGo} of the
repository, completed earlier the same day, which treated the no-go
theorems in eight dimensions with octonions and triality. That folder
was not modified; this one is new. Its registry of symbols was read in
order to verify that the symbols shared by both folders denote the same
objects (check LL-S2).

\hypertarget{the-request-of-november-16-2026-1932-verbatim}{%
\subsection{2.3 The request of November 16, 2026, 19:32
(verbatim)}\label{the-request-of-november-16-2026-1932-verbatim}}

``explain for an ignorant exactly what your manuscript comment''The
spin-zero superfluid of its spinor charges condenses in an
isospin-triplet, unitary, electrically neutral state, and the pair that
is a singlet of the whole super-loop re- quires two species. '' means.
Explain clearly, correctly, and in entirety. create a pdf file for this
answer; give examples''

During the work on it, at 20:21, the author added: ``when complete, push
everything to repo''.

\hypertarget{what-the-second-request-asks-for-and-how-it-is-answered}{%
\subsection{2.4 What the second request asks for, and how it is
answered}\label{what-the-second-request-asks-for-and-how-it-is-answered}}

\begin{enumerate}
\def\labelenumi{\arabic{enumi}.}
\tightlist
\item
  \textbf{What the sentence means, exactly, for a reader who does not
  know the solution.} The sentence stands in the conclusions of the
  manuscript. The answer is the document
  \texttt{explanations/The\ spin-zero\ superfluid\ of\ the\ spinor\ charges\ -\ an\ explanation}
  (.md, .tex, .pdf). It splits the sentence into seven statements and
  explains each: the particles (the spinor charges as a fermion field),
  ``spin-zero'', ``superfluid'' and ``condenses'', ``isospin-triplet'',
  ``unitary'', ``electrically neutral'', and ``the pair that is a
  singlet of the whole super-loop requires two species''.
\item
  \textbf{Clearly, correctly and in entirety.} Every term is defined
  before it is used, every claim is proved in the document, and the
  document is self-contained (it contains the complete table of the
  symbols). Where the sentence is shorter than the facts, the document
  says so: ``electrically neutral'' holds for every unitary condensate
  in the sense that the pairs have the average charge 0, and in the
  strict sense (every pair neutral, the charge unbroken) for the
  representative along the third isospin axis that the manuscript names;
  the model, which is exactly isospin invariant, does not choose among
  the equivalent directions (Section 8.4 of that document).
\item
  \textbf{A PDF file.} By Hard Rule 1 the document exists as .md, .tex
  and .pdf.
\item
  \textbf{Examples.} Worked examples from atoms (helium), nuclei (the
  spin-zero isospin-triplet pairs of nucleons, the deuteron), helium 3
  and quark matter; seven example condensates computed in full; the
  counting of two-particle states; three figures. Every example is
  computed by the new Rust program \texttt{explain\_superfluid} (folder
  \texttt{explanations/code/explain-rs}, checks X-1 to X-10), and the
  numbers of the superfluid that the document quotes are those of the
  command \texttt{superfluid} of \texttt{superloop}.
\end{enumerate}

\hypertarget{the-requests-of-november-17-2026-their-operative-parts-verbatim}{%
\subsection{2.5 The requests of November 17, 2026 (their operative
parts,
verbatim)}\label{the-requests-of-november-17-2026-their-operative-parts-verbatim}}

At 05:34 the author asked for a test of the sentence of the conclusions
that follows the one explained in stage 4. The operative part of the
request, verbatim:

``Immediately after the sentence cited in your previous task, you state
`In 3 + 1 dimensions its Yukawa form is renormalizable \ldots{}'
{[}\ldots{]} So You must clearly mathematically define , discuss, give
examples (with figures) of your 3+1 supporting manifold for this
statement `In 3 + 1 dimensions its Yukawa form is renormalizable'; you
must do this in 240 seconds or less, and NOT HOURS.''

At 05:36, verbatim: ``none of the figures are new! what are you doing.
sketch the coordinates , similar to sketching x-y-z- Cartesian
coordinates''

At 05:42 the author asked for the work to be done without shortcuts, and
ended, verbatim: ``Put your test answers in a new, Well-Organized pdf
file that is precise, correct and COMPLETE!''

\hypertarget{what-these-requests-ask-for-and-how-they-are-answered}{%
\subsection{2.6 What these requests ask for, and how they are
answered}\label{what-these-requests-ask-for-and-how-they-are-answered}}

The answer is the document
\texttt{explanations/The\ 3+1\ manifold\ of\ the\ super-loop\ and\ the\ renormalizability\ of\ the\ Yukawa\ form}
(.md, .tex, .pdf), new in stage 5.

\begin{enumerate}
\def\labelenumi{\arabic{enumi}.}
\tightlist
\item
  \textbf{Define the \(3+1\) manifold mathematically.} Its Section 2:
  the manifold \(\mathbb R^{1,3}_{\rm LL}\) of the Hermitian
  bi-quaternions, its coordinates, the matrix \(M_{\rm LL}(\mathbf x)\)
  of a point, and the metric \(N_{\rm LL}=\det M_{\rm LL}(\mathbf x)\)
  of signature \((1,3)\).
\item
  \textbf{Discuss it, for the statement.} Its Sections 3 to 5:

  \begin{itemize}
  \tightlist
  \item
    the Lorentz group of the manifold from the unit bi-quaternions of
    the super-loop (double cover, generators, all proper orthochronous
    transformations reached);
  \item
    the spinor charges as two-component spinors on it;
  \item
    the contact and Yukawa forms of the field theory on it;
  \item
    the proof of the statement: mass dimensions, power counting, the
    four divergent functions and the terms that absorb them, the
    completeness of the Lagrangian, the one-loop bubble, why the contact
    form and the eight-dimensional manifolds fail, and what the standard
    theorem adds.
  \end{itemize}
\item
  \textbf{Examples with figures.} Worked examples (a sample point, a
  boost, a rotation, the rotation by \(2\pi\), a spinor and its null
  point, the spinor charges, the bubble at ten cutoffs) and four new
  figures. Figure 1 draws the coordinates like Cartesian axes, as the
  request of 05:36 asks.
\item
  \textbf{A new, well-organized PDF that is precise, correct and
  complete.}

  \begin{itemize}
  \tightlist
  \item
    The document answers each question of the test in a section of its
    own and collects the answers in one table (its Section 6).
  \item
    Every number is computed by the new Rust program
    \texttt{r13\_yukawa} (folder \texttt{explanations/code/r13-rs},
    checks Z-1 to Z-8) and inserted from its data files.
  \item
    The document is self-contained: it contains the complete table of
    the symbols, the reports of its checks, and the reports of the
    checks of \texttt{superloop} that it quotes.
  \end{itemize}
\item
  \textbf{The time limit of 240 seconds} could not be met together with
  the other requirements (D9).
\end{enumerate}

\hypertarget{test-2-the-requests-of-november-17-2026-0641-0645-and-0703-verbatim}{%
\subsection{2.7 Test 2: the requests of November 17, 2026, 06:41, 06:45
and 07:03
(verbatim)}\label{test-2-the-requests-of-november-17-2026-0641-0645-and-0703-verbatim}}

The request of 06:41, verbatim (lines broken for the page; the formula
of the request is reproduced as typed):

\begin{verbatim}
BIG, CRITICAL Test 2, for which you MUST  Clearly and completely answer, with provenance the following
(complete means write out everything; if you use expressions like $X=x^0\sigma_0+\sum_j x^j\sigma_j$ on
one line, then on the next line you MUST explicitly evaluate the matrices ; and so on) ; push to repo
after each item is completed. You are being timed, but there is not time limit (for now): [1] Write out,
explain and discuss the field equations for a self-interacting spinor field that is defined on this
manifold (these field equations are a sort of Dirac equation, specialized to this manifold, and uses
super-loop operators); your solution must be completely formulated, evaluated and explained for future
numerical solution; [2] Analyze these field equations and provide analytical solutions for at least 3
simple cases; [3] write a python script that provides guidance to the user for defining initial and
boundary conditions, and then numerically solves  these field equations after reading initial and
boundary conditions from a user file (query for file name); [4] write a python script that  plots these
numerical solutions; [5] write a full pure rust program that provides guidance to the user for defining
initial and boundary conditions, and then numerically solves  these field equations after reading
initial and boundary conditions from a user file (query for file name); use cargo to build; provide
great instructions for use; [6] write a full pure rust program that plots these rust numerical
solutions; [7] Write out, explain and discuss the construction of the new superfluid boson singlet
(\(S=0\)) state  (SFBS) from the spinor field; [8] Derive , explain and discuss the effective field
equations for superfluid boson singlet  (\(S=0\)) state; [8] Derive , explain and discuss the effective
potential for the superfluid boson singlet  (\(S=0\)) state (similar to the Mexican Hat potential for the
complex Higgs doublet); [9] find the relative minima of this effective potential; spontaneously break
the symmetry, following the example of the Higgs boson; Expand the the superfluid boson singlet (\(S=0\))
state about this minimum, similar to the treatment of the Higgs boson call this approximation SFBSv.
[10] Derive , explain and discuss the kinetic Lagrangian for the superfluid boson singlet  (\(S=0\))
state. [11] Give mass to the Goldstone bosons by substituting SFBSv into the kinetic Lagrangian and
expanding ; diagonalize mass terms and work out the surviving symmetry group; prove this group is related
to U(1)_em and the massless photon.
\end{verbatim}

The request of 06:45, verbatim:

\begin{verbatim}
i should not have to remind you that a gauge group (related to  SL(2,C) ×SU(2)v ×U(1)F LL) acts on
spinors on this manifold, and that YOU MUST EMPLOY A GAUGE COVARIANT DERIVATIVE in all field equations;
do this OR FAIL this TEST NOW.
\end{verbatim}

The request of 07:03, verbatim:

\begin{verbatim}
Remember to Put your test answers in a new, Well-Organized pdf file that is precise, correct and
COMPLETE!
\end{verbatim}

\hypertarget{what-test-2-asks-for-and-how-it-is-answered}{%
\subsection{2.8 What Test 2 asks for, and how it is
answered}\label{what-test-2-asks-for-and-how-it-is-answered}}

The answers are one new document,
\texttt{test2/Test\ 2\ -\ the\ gauge-covariant\ spinor\ field\ on\ the\ 3+1\ manifold\ and\ the\ superfluid\ boson\ singlet\ state}
(.md, .tex, .pdf), which receives one section per item, in the order of
the test, and is committed and carried to the author's computer after
each item (rule 2, Section 3).

\begin{longtable}[]{@{}
  >{\raggedright\arraybackslash}p{(\columnwidth - 4\tabcolsep) * \real{0.0732}}
  >{\raggedright\arraybackslash}p{(\columnwidth - 4\tabcolsep) * \real{0.7317}}
  >{\raggedright\arraybackslash}p{(\columnwidth - 4\tabcolsep) * \real{0.1951}}@{}}
\toprule\noalign{}
\begin{minipage}[b]{\linewidth}\raggedright
item
\end{minipage} & \begin{minipage}[b]{\linewidth}\raggedright
answer
\end{minipage} & \begin{minipage}[b]{\linewidth}\raggedright
state
\end{minipage} \\
\midrule\noalign{}
\endhead
\bottomrule\noalign{}
\endlastfoot
{[}1{]} & Section 2 of the document: the spinor field of the spinor
charges on \(\mathbb R^{1,3}_{\rm LL}\); the gauge group
\(SL(2,\mathbb C)\times SU(2)_{\rm v}\times U(1)^{\rm F}_{\rm LL}\) with
\(SL(2,\mathbb C)\) as Lord's change of representation of the frame; the
covariant derivative with every matrix written out; the field strengths;
why the Lorentz connection is pure gauge on the flat manifold; the
Lagrangian; the exact field equations (operator form, Hamiltonian form,
component form, super-loop form: right multiplication of the spinor
charges by the frame); the gauge-field equations; the mean-field
(Dirac-Bogoliubov-de Gennes) equations with all sixteen components
written out; the gauge-covariant lattice discretization, the time
stepping and their properties; checks GS-1 to GS-16 of the new program
\texttt{t2\_fields}; two figures & answered (stage 6) \\
{[}2{]} & Section 3 of the document: the analysis of the equations (each
mode keeps its norm; the fermion number is conserved exactly when the
pair field is self-consistent; constant potentials are pure gauge except
for their holonomy on a ring; plane waves reduce the equations to
matrices; the square of the covariant Dirac Hamiltonian and the Pauli
equation of slow spinor charges, with the gyromagnetic factor 2 for both
kinds of charge) and four cases solved exactly: free spinor charges in
constant gauge potentials (momentum and energy shifts, the isospin
precession); the uniform polar superfluid (Bogoliubov modes, the
particle-hole oscillation, the gap equation in closed form for massless
charges and its BCS limit, an exact self-consistent lattice solution); a
domain wall of the pair field (bound states at exactly
\(\pm m_{\rm LL}\), the square of the operator and the complete spectrum
of the bound states, the first-order Wilson shift on the lattice); the
Landau levels; checks AS-1 to AS-10 of the new program
\texttt{t2\_analytic}; four figures & answered (stage 7) \\
{[}3{]} & Section 4 of the document: the Python solver
\texttt{test2/python/t2\_solve.py} (guidance, the question for the input
file, the input file of lines \texttt{key\ =\ value} with 50 checked
keys, the tasks evolve, spectrum, ground\_state and quench), its
numerical methods with their conservation properties, and seven examples
with known answers checked by \texttt{t2\_examples\_check.py} &
answered \\
{[}4{]} & Section 5 of the document: the plotter
\texttt{test2/python/t2\_plot.py} (one figure per solution, its panels
chosen by the task; for the seven examples the predictions of the checks
EX-1 to EX-7 drawn with the solutions) and Figures 7 to 13 & answered \\
{[}5{]} & Section 6 of the document: the Rust solver \texttt{t2\_solve}
(member \texttt{t2solve} of the Cargo workspace \texttt{test2/rust}),
with the instructions for installing Rust, building with cargo and using
it; the same input and output files as the Python solver; the comparison
EX-8 of the two sets of solutions & answered \\
{[}6{]} & Section 7 of the document: the Rust plotter \texttt{t2\_plot}
(member \texttt{t2plot} of the workspace \texttt{test2/rust}), which
writes Scalable Vector Graphics with its own small plotting library and
draws the predictions of the examples computed in Rust; Figures 14 to 20
& answered \\
{[}7{]} & Section 8 of the document: the construction of the superfluid
boson singlet state (SFBS) from the spinor field: the 28 pairs of two
spinor charges at one point, the six spin-zero pairs (spin singlet,
isospin triplet by the Pauli principle), the scalar channel selected by
the interaction and by parity, the charges, the super-loop form (the
imaginary part of the product of a spinor charge's bi-quaternion with
its conjugate), the Bose character with Pauli blocking, and the SFBS
state as a coherent state of particle and antiparticle pairs, proved to
be the mean-field ground state; checks SB-1 to SB-5 of the new program
\texttt{t2\_sfbs} (member \texttt{t2sfbs}; exact many-body computations
in Fock spaces of 256 and 65536 states); Figure 21 & answered \\
{[}8{]} & Section 9 of the document: the effective field equations of
the SFBS: the Hubbard-Stratonovich field \(\vec d_{\rm LL}\) (the SFBS
times \(-\tfrac12G_{\rm LL}\), its transformation law), its
gauge-covariant derivative \(\nabla^{\rm B}_{{\rm LL},p}\) inherited
from the spinor field by the Leibniz rule (covariance and commutator),
the effective action and its derivative expansion
\(\mathcal L^{\rm B}_{\rm LL}\), the pair susceptibility of the vacuum
with \(1/G^{\rm crit}_{\rm LL}\), \(Z^{\rm t}_{\rm LL}(0)\),
\(Z^{\rm s}_{\rm LL}(0)\) and \(1/G^{\rm bd}_{\rm LL}\) in closed form,
the effective field equations written out and in first-order form, the
gauge-field equations with the currents of the SFBS, and the SFBS boson
of the normal vacuum (Klein-Gordon form, exact pole, instability beyond
\(G^{\rm crit}_{\rm LL}\)), confirmed by the time-dependent mean-field
equations; checks SB-6 to SB-10, Figure 22 & answered \\
{[}8{]} (second) & Section 10 of the document: the effective potential
of the SFBS: its two invariants (the normal form of a pair field,
Theorem 10.1), the splitting of the Bogoliubov matrix into two sectors
with the gaps \(\Delta^\pm_{\rm LL}\), the potential as the cost of the
field minus the condensation energies \(W^{\rm g}_{\rm LL}\) of the
sectors (increasing and strictly concave), the Ginzburg-Landau form with
\(a^{\rm GL}_{\rm LL}=1/G_{\rm LL}-1/G^{\rm crit}_{\rm LL}\) and
\(b^{\rm GL}_{\rm LL}=Z^{\rm t}_{\rm LL}(0)\), the closed form at
\(\mu_{\rm LL}=0\) with the gap equation as its stationarity, the
massless form with its logarithm, the next term, and the Mexican hat;
checks SB-11 and SB-12, Figure 23 & answered \\
{[}9{]} & Section 11 of the document: the minima of the effective
potential (the two-sector form with the potential
\(\Omega^{\rm s}_{\rm LL}\) of one sector, strictly convex in the
squared gap; Theorem 11.1: every relative minimum is a global one, on
the vacuum manifold of the polar fields of modulus
\(\Delta^{\rm vac}_{\rm LL}\); Theorem 11.2: at \(\mu_{\rm LL}=0\) the
only stationary points are the maximum \(\vec d_{\rm LL}=0\), the minima
and the non-unitary saddle points, with their Hessians; the gradient
flow), the spontaneous breaking of the symmetry (the vacuum
\((0,0,\Delta^{\rm vac}_{\rm LL})\), the generators acting on the field
written out, the broken generators \(F_{\rm LL}\), \(T^1_{\rm LL}\),
\(T^2_{\rm LL}\), the unbroken electric charge
\(Q_{\rm LL}=T^3_{\rm LL}\), the stabilizer \(O(2)\) written out,
Goldstone's theorem), and SFBSv (the parametrization by three Goldstone
fields, the amplitude field and two non-unitary fields, with its
matrices and components written out; the exact potential of SFBSv,
independent of the Goldstone fields; its expansion to the fourth order
with coefficients in closed form; the charges of the fields; the
comparison with a potential of the Higgs form); checks SB-13 and SB-14,
Figure 24 & answered \\
{[}10{]} & Section 12 of the document: the kinetic Lagrangian of the
SFBS (its definition with the gauge-covariant derivative and the reasons
for its form; written out with the couplings to the currents and the
contact matrix \(\mathsf N^{\rm g}_{\rm LL}\) of the four potentials,
which in the vacuum does not contain the potential of the electric
charge; Theorem 12.1: the stiffness coefficients
\(Z^{\rm t}_{\rm LL}(\Delta_{\rm LL})\),
\(Z^{\rm s}_{\rm LL}(\Delta_{\rm LL})\) of the condensed vacuum in
closed form, derived exactly from the energy of the vacuum in constant
gauge potentials, the same for the fermion-number and the isospin
directions, exact for the Goldstone fields by gauge invariance; Theorem
12.2: the exact kernels of the six directions at zero wave number, with
the masses 0 of the Goldstone fields and \(2\Delta^{\rm vac}_{\rm LL}\)
of the amplitude field and the non-unitary fields, below the
pair-breaking threshold; the quadratic Lagrangian of SFBSv; the
time-dependent mean-field oscillation at \(2\Delta^{\rm vac}_{\rm LL}\);
the surface term of the cutoff and dimensional regularization); checks
SB-15 and SB-16, Figure 25 & answered \\
{[}11{]} & Section 13 of the document: the masses of the gauge bosons
(the Lagrangian with the gauge fields; Theorem 13.1: inserting SFBSv,
the gauge transformation of the unitary gauge removes the three
Goldstone fields exactly, and they reappear as the gradient parts of the
potentials of their broken generators; the kinetic Lagrangian in the
unitary gauge written out exactly with its contact matrix (13.4); the
expansion to the second order in the original fields, with the mixing of
the Goldstone fields and the potentials, and in the unitary gauge, where
it is diagonal; the mass matrices (13.7) of the time and the space
components, the charged fields \(W^\pm_{{\rm LL},p}\), the rest masses
\(M^{\rm F}_{\rm LL}\) and \(M^{\rm W}_{\rm LL}\) (13.8), and the
dispersion relations of the transverse and the longitudinal modes
(13.9); Theorem 13.2: the surviving group \(O(2)\) with its connected
part \(U(1)^{\rm Q}_{\rm LL}\), generated by the electric charge,
identified with \(U(1)_{\rm em}\); the mass matrix as the squared length
of the change of the vacuum under a gauge transformation, so that the
photon \(V^3_{{\rm LL},p}\) is the only massless gauge field; the
absence of mixing in the vacuum and the mixing at the non-unitary saddle
point (13.11); the count of the degrees of freedom), the check SB-17
(the unitary gauge, the written-out Lagrangian, and the complete mass
matrices, the photon and the mixing directly from the Bogoliubov
matrices) and Figure 26 & answered \\
\end{longtable}

The demand of 06:45 is met in every field equation of the document: the
derivatives enter only through the covariant derivative
\(\nabla^{\rm LL}_p\), through its lattice form (the link matrices), or,
in the gauge-field equations, through the covariant divergence. The
demand of the request of 06:41 to evaluate every matrix is met by
printing every matrix and every component equation from the program,
which evaluates each printed expression and compares it with the matrix
that it describes (check GS-12).

\hypertarget{the-requests-of-november-18-2026-0533-and-0535-verbatim}{%
\subsection{2.9 The requests of November 18, 2026, 05:33 and 05:35
(verbatim)}\label{the-requests-of-november-18-2026-0533-and-0535-verbatim}}

The request of 05:33, verbatim:

\begin{verbatim}
STOP everything NOW. I do not want you to use the sophomoric RK4 solver. INSTEAD, Use the pure rust
version of cvode on my repo https://github.com/once-ere/rustSolveIt_macos-silicon_SUNDIALS_7_8_0.git.
Redo every calculation that you performed and use cvode instead (or some other applicable SUNDIALS
algorithm). Create a provenance that compares the cvode result with the RK4 result; REMEMBER these
errors, and do NOT do this again!
\end{verbatim}

The request of 05:35, verbatim:

\begin{verbatim}
include all of your new rust code, along with docs, in this 'spacetimes' PROJECT
\end{verbatim}

\hypertarget{what-these-requests-ask-for-and-how-they-are-answered-1}{%
\subsection{2.10 What these requests ask for, and how they are
answered}\label{what-these-requests-ask-for-and-how-they-are-answered-1}}

\begin{enumerate}
\def\labelenumi{\arabic{enumi}.}
\tightlist
\item
  \textbf{CVODE instead of the fixed-step methods.} The first version of
  the programs of Test 2 integrated the time-dependent equations with
  two methods written by hand: the classical Runge-Kutta method of the
  fourth order (RK4) with a fixed step (the task \texttt{evolve} of both
  solvers, the checks AS-2, AS-4 of \texttt{t2\_analytic}, SB-10 and
  SB-16 of \texttt{t2\_sfbs}) and the implicit midpoint rule with a
  fixed step (the task \texttt{quench}). Every one of these integrations
  now uses CVODE of SUNDIALS 7.8.0 in the pure-Rust translation of the
  author's repository: the variable-order, variable-step Adams-Moulton
  formulas with control of the local error
  (\(r^{\rm tol}_{\rm LL}=10^{-12}\), \(a^{\rm tol}_{\rm LL}=10^{-14}\),
  error weights of the complex components; Section 2.11 of the document
  of Test 2). The Rust programs call it through the new module
  \texttt{t2core::cvode}; the Python solver calls the same module
  through the new shared library \texttt{libt2cvode} (crate
  \texttt{test2/rust/t2cvode\_c}), so that there is one integrator
  (decision D22). Every calculation that used a fixed-step method was
  run again: the checks AS-2, AS-4, GS-13 (now a test of CVODE against
  an exact solution), SB-10 and SB-16, the seven examples of both
  solvers, the examples checks EX-1 to EX-8, and all figures that show
  their results.
\item
  \textbf{The comparison.} The new program \texttt{t2\_cvcmp} (crate
  \texttt{test2/rust/t2cvcmp}), with its own copy of the RK4 step of the
  first version, compares the two kinds of integration on three problems
  with exact solutions and on the four time-dependent examples of the
  solvers; the data of the first version are kept in
  \texttt{test2/rk4\_reference/data} (106 files, with SHA-256 digests).
  The comparison is written up, with its figures, in
  \texttt{test2/cvode\_vs\_rk4/P\_CVODE\_vs\_RK4} (.md, .tex, .pdf); the
  document of Test 2 contains its program, its report and its findings
  in its Appendix A.2g. The main finding: the quench of the first
  version was wrong in the third decimal (the mean of the gap 0.500877
  instead of 0.503125, the frequency 1.18449 instead of 1.18898), and
  the stability rule of the first version allowed steps that give large
  errors on random lattice problems.
\item
  \textbf{The code in the project.} The pure-Rust translation of
  SUNDIALS 7.8.0 is part of PROJECT spacetimes, in the folder
  \texttt{vendor/sundials\_rs-7.8.0} (its seven crates, unchanged, with
  the digests of its 283 files and its own provenance
  \texttt{vendor/P\_sundials\_rs-7.8.0}, .md, .tex, .pdf); the new Rust
  code of Test 2 is in the members \texttt{t2core} (\texttt{cvode.rs}),
  \texttt{t2cvode\_c} and \texttt{t2cvcmp} of the workspace
  \texttt{test2/rust}.
\item
  \textbf{Remember the errors.} The error is entry E63 of Section 11,
  and the rule that prevents it is rule 9 of Section 3: a time
  integration is never done with a fixed-step method written by hand.
\end{enumerate}

\hypertarget{the-hard-rules-and-how-each-is-met}{%
\section{3. The hard rules, and how each is
met}\label{the-hard-rules-and-how-each-is-met}}

\begin{longtable}[]{@{}
  >{\raggedright\arraybackslash}p{(\columnwidth - 4\tabcolsep) * \real{0.3333}}
  >{\raggedright\arraybackslash}p{(\columnwidth - 4\tabcolsep) * \real{0.3333}}
  >{\raggedright\arraybackslash}p{(\columnwidth - 4\tabcolsep) * \real{0.3333}}@{}}
\toprule\noalign{}
\begin{minipage}[b]{\linewidth}\raggedright
rule (in the words of this provenance)
\end{minipage} & \begin{minipage}[b]{\linewidth}\raggedright
how it is met
\end{minipage} & \begin{minipage}[b]{\linewidth}\raggedright
where to verify
\end{minipage} \\
\midrule\noalign{}
\endhead
\bottomrule\noalign{}
\endlastfoot
0. Write for clever physics Ph.D.~students who do not know the solution;
they must be able to learn from it and teach it & the manuscript defines
every term before it uses it (generator, Jacobiator, associator,
alternative algebra, loop, Moufang loop, Malcev algebra, derivation,
automorphism, twist charge, Gröbner basis, saturation, unitary state,
Bogoliubov-de Gennes matrix, superficial degree of divergence, \ldots),
proves its seven theorems in the text, and states what is not
established; the explanation of the sentence on the superfluid, and the
document on the \(3+1\) manifold and the Yukawa form, are written for a
reader who knows nothing of the solution, with worked examples &
manuscript, Secs. I--XI; the two documents in \texttt{explanations/} \\
1. Every documentation, symbol and provenance file in .md, .tex and .pdf
& \texttt{README}, \texttt{P0}, \texttt{P1}, \texttt{P2}, the
manuscript, the two documents of \texttt{explanations/} and the document
of the answers to Test 2 in \texttt{test2/} exist in all three formats;
\texttt{check\_all.sh} builds them and fails if one is missing &
\texttt{check\_all.sh}, Section 4 \\
2. Push to the repository as soon as anything is updated & the work was
committed in stages (Section 5), each carried to the author's computer
at once; for Test 2, after each item; the push itself can only be made
by the author, with the command block given at the end of each stage &
the git log \\
3. Working directory \texttt{/Users/nsh/Developer/github/spacetimes} &
the folder \texttt{Revision\_SuperLoop} of that repository & --- \\
4. Repository \texttt{https://github.com/once-ere/spacetimes.git} & as
stated & --- \\
5. One symbol, one object, in the whole project; a master provenance of
the symbols, and its symbols used everywhere & the registry
\texttt{provenance/symbols\_ll.csv} takes 2 symbols of the master
registries with their meanings and adds 201 new ones, all with the
suffix LL or in a type (bold upright, sans-serif, blackboard, fraktur)
that marks them as algebra elements, matrices, number sets or Lie
algebras; the plain-text names of the programs, figures and data files
are registered as well; nine machine checks (LL-S1 to LL-S9) &
\texttt{P1}, data file \texttt{sl\_symbol\_checks.txt} \\
6. Source code for every equation, result, figure and prediction,
commented, with instructions and provenance; Rust first, then C, Python,
Mathematica; execute the code; a carefully designed figure for each
prediction & the Rust program \texttt{superloop} computes every result
except the Gröbner bases (62 checks, all passing); the Gröbner bases are
computed in Python (deviation D3) and every solution found there is
verified exhaustively by the Rust program; a Python script draws the
seven figures from the data files; the Rust program
\texttt{explain\_superfluid} computes every example of the explanation
(10 checks, all passing), and a Python script draws its three figures;
the Rust program \texttt{r13\_yukawa} computes every number of the
document on the \(3+1\) manifold (8 checks, all passing), and a Python
script draws its four figures; the pure Rust program \texttt{t2\_fields}
(folder \texttt{test2/rust}) computes and checks every matrix and
equation of item {[}1{]} of Test 2 (16 checks, all passing), and a
Python script draws its two figures; every number in the documents is
inserted from the data files by a script & \texttt{P2}, \texttt{code/},
\texttt{explanations/code/}, \texttt{test2/} \\
7. Every equation, symbol and definition in the very file that uses it;
never refer the reader elsewhere & every document contains the complete
table of the symbols and the formulas it uses; the check LL-H1 searches
all documents for referrals & data file
\texttt{sl\_self\_contained\_checks.txt} \\
9. (the request of November 18, 2026) Time integrations with SUNDIALS,
never with a fixed-step method written by hand & every time-dependent
calculation of the folder (the tasks \texttt{evolve} and \texttt{quench}
of both solvers, the checks AS-2, AS-4, GS-13, SB-10 and SB-16) calls
CVODE of the vendored pure-Rust SUNDIALS 7.8.0 through
\texttt{t2core::cvode} (the Python solver through \texttt{libt2cvode});
the only remaining copy of the RK4 step is the one inside the comparison
program \texttt{t2\_cvcmp}, which measures the errors of the first
version & \texttt{test2/rust/t2core/src/cvode.rs}; the reports of the
checks; \texttt{test2/cvode\_vs\_rk4/data/cmp\_report.txt} \\
8. Follow the instructions exactly, or explain in a provenance why not
and pause & Section 7 explains every place where the wording of the task
cannot be followed literally (a counterexample to a proved theorem whose
hypotheses hold, a ``group'' that is not associative) and what is
provided instead; the work did not pause, because every instruction
could be carried out in full on that basis & Section 7 \\
\end{longtable}

\hypertarget{what-the-folder-contains}{%
\section{4. What the folder contains}\label{what-the-folder-contains}}

\begin{longtable}[]{@{}
  >{\raggedright\arraybackslash}p{(\columnwidth - 2\tabcolsep) * \real{0.5000}}
  >{\raggedright\arraybackslash}p{(\columnwidth - 2\tabcolsep) * \real{0.5000}}@{}}
\toprule\noalign{}
\begin{minipage}[b]{\linewidth}\raggedright
path
\end{minipage} & \begin{minipage}[b]{\linewidth}\raggedright
content
\end{minipage} \\
\midrule\noalign{}
\endhead
\bottomrule\noalign{}
\endlastfoot
\texttt{manuscript/} & the manuscript ``The Lorentz Super-Loop of the
Bi-Octonions: Non-Associative Counterexamples to the No-Go Theorems and
a Spin-Zero Spinor Superfluid'' (.tex, .pdf, .md) and its reference list
\texttt{sl\_refs.bib} \\
\texttt{provenance/} & this document, the registry of the symbols
(\texttt{P1}), the calculations (\texttt{P2}), each as .md, .tex, .pdf;
\texttt{symbols\_ll.csv}; the generated fragments \\
\texttt{code/superloop-rs/} & the Rust program \texttt{superloop},
commands \texttt{algebra}, \texttt{malcev}, \texttt{manifolds},
\texttt{superfluid}, \texttt{renorm}, \texttt{all} \\
\texttt{code/scripts/} & \texttt{malcev\_groebner.py},
\texttt{make\_figures.py}, \texttt{write\_symbol\_registry.py},
\texttt{check\_symbols.py}, \texttt{check\_self\_contained.py},
\texttt{write\_doc\_tables.py}, \texttt{assemble\_docs.py},
\texttt{build\_manuscript.py}, \texttt{tex2md.py},
\texttt{build\_docs.sh}, \texttt{check\_layout.py},
\texttt{doc\_header.tex}, \texttt{check\_all.sh} \\
\texttt{data/} & 14 data files, the reports of all checks and the
Gröbner bases \\
\texttt{figures/} & 7 figures, .pdf and .png \\
\texttt{explanations/} & the document ``The spin-zero superfluid of the
spinor charges - an explanation'' (.md, .tex, .pdf), with its program
\texttt{code/explain-rs} (binary \texttt{explain\_superfluid}), its
figure script \texttt{code/make\_explain\_figures.py} and its three
figures; the document ``The 3+1 manifold of the super-loop and the
renormalizability of the Yukawa form'' (.md, .tex, .pdf), with its
program \texttt{code/r13-rs} (binary \texttt{r13\_yukawa}), its figure
script \texttt{code/make\_r13\_figures.py} and its four figures; the
data files and reports of both programs in \texttt{data/}, the figures
in \texttt{figures/} \\
\texttt{test2/} & the answers to Test 2 (.md, .tex, .pdf), the Cargo
workspace \texttt{rust/} (library \texttt{t2core}, program
\texttt{t2\_fields}), the figure script
\texttt{python/make\_t2\_figures.py}, the data files and reports in
\texttt{data/}, the figures in \texttt{figures/} \\
\texttt{docs\_src/} & the sources of the documents \\
\end{longtable}

\hypertarget{what-was-done-in-order}{%
\section{5. What was done, in order}\label{what-was-done-in-order}}

\begin{enumerate}
\def\labelenumi{\arabic{enumi}.}
\tightlist
\item
  \textbf{Stage 1 --- the algebra and the classification.} The
  bi-octonions with exact Gaussian-integer coefficients; the places
  where the proofs of the no-go theorems use associativity; the Lorentz
  super-loop, its Malcev algebra, its operators and their eigenvalues;
  Theorems R and L; the classification of the Lorentz- and
  isospin-covariant brackets by Gröbner bases (Theorem M) and its
  exhaustive verification; the manifolds; the superfluid of the spinor
  charges (Theorems SF1 to SF4); power counting and the Yukawa
  completion.
\item
  \textbf{Stage 2 --- the figures}, one for each group of predictions,
  and the script that draws them.
\item
  \textbf{Stage 3a --- one meaning for every name}: the registry of the
  symbols, its checks, a revision of all plain-text names of the
  programs, figures and data files, the new check A-8 (Theorem R), and
  one sign convention for the isospin generators (Section 11, entries E1
  and E2).
\item
  \textbf{Stage 3b --- two corrections in the program and a figure}: the
  credit for the non-Lie Malcev module of \(\mathfrak{sl}(2)\), the
  charge of the condensed pairs in check S-7, and the labels of the
  theorems of the superfluid (Section 11, entries E3 to E5).
\item
  \textbf{Stage 3c --- the documents}: the manuscript, the provenance
  documents, the builds, and the repairs that the builds required
  (Section 11, entries E9 to E13).
\item
  \textbf{Stage 4 --- the explanation} of the sentence on the superfluid
  (request of 19:32, Section 2.3): the program
  \texttt{explain\_superfluid} with the checks X-1 to X-10, three
  figures, the document in \texttt{explanations/}, and the extension of
  the checks of the symbols, of self-containedness and of
  \texttt{check\_all.sh} to the new folder (Section 11, entries E14 to
  E17).
\item
  \textbf{Stage 5 --- the \(3+1\) manifold and the renormalizability of
  the Yukawa form} (requests of November 17, Section 2.5): the program
  \texttt{r13\_yukawa} with the checks Z-1 to Z-8, four new figures, the
  document in \texttt{explanations/}, new registry entries
  (\(\mathsf L_{\rm LL}(\mathbf g)\), \(\chi_{\rm LL}\),
  \(\theta_{\rm LL}\), \(\mathcal L^{\rm Y}_{\rm LL}\) and
  \(\mathcal L^{\rm c}_{\rm LL}\), \(m^\phi_{\rm LL}\),
  \(I^{\rm ren}_{\rm LL}\), \(\mathfrak{so}(4)\)), and the extension of
  the checks and of \texttt{check\_all.sh} to the new document (Section
  11, entries E18 to E24).
\item
  \textbf{Stage 6 --- Test 2, item {[}1{]}} (requests of 06:41, 06:45,
  07:03, Section 2.7): the Cargo workspace \texttt{test2/rust} with the
  library \texttt{t2core} and the program \texttt{t2\_fields} (checks
  GS-1 to GS-16), the explicit matrices and component equations printed
  by the program, two figures, the new document in \texttt{test2/}, the
  registry group 9 (the gauge fields, the covariant derivative, the
  mean-field and lattice equations), the new check LL-S9, and the
  extension of the checks and of \texttt{check\_all.sh} to the folder
  \texttt{test2} (Section 11, entries E25 to E36).
\item
  \textbf{Stage 7 --- Test 2, item {[}2{]}} (Section 2.7): the program
  \texttt{t2\_analytic} (checks AS-1 to AS-10) in the workspace
  \texttt{test2/rust}, which applies the lattice operator of the solvers
  to four exactly solvable cases; additions to the library
  \texttt{t2core} (the RK4 step of the lattice equations, the matrix of
  the free lattice operator on a plane wave, a faster form of the
  Hermitian eigensolver); four figures; Section 3 of the document of
  Test 2; eleven new registry entries of group 9 (the potentials and
  magnetic fields felt by the two isospin components, the precession
  frequency, the probabilities of isospin up and of the hole, the
  constant spinor of a plane wave and the Bogoliubov angle, the domain
  wall and its bound states, the matrix of the square of its operator,
  the Rayleigh quotient, the number of a Landau level), six new
  descriptive column words, and \texttt{check\_all.sh} extended by the
  new program (Section 11, entries E37 to E43).
\item
  \textbf{Stage 8 --- Test 2, item {[}3{]}} (Section 2.8): the Python
  solver \texttt{test2/python/t2\_solve.py} with its self-test, seven
  example input files in \texttt{test2/examples}, the program
  \texttt{test2/python/t2\_examples\_check.py} (checks EX-1 to EX-7 of
  the examples), Section 4 of the document of Test 2 (and the
  corrections of Section 2.11 that it required: the normalization in
  fewer than three dimensions, and the methods of the self-consistent
  tasks), twelve new registry entries of group 9 (the number of active
  directions, the fermion number and the mean-field energy of a lattice
  state, the matrix of the modes, the quantities of the Anderson
  iteration, the wave packet and the branch projectors, the matrix of
  the free lattice operator on a plane wave, the group velocity, the
  magnetic length, the quench and the amplitude mode, and the Bloch
  vectors of the linearized equations), nineteen new descriptive column
  words, and \texttt{check\_all.sh} extended by the solver, its examples
  and their checks (Section 11, entries E44 to E52).
\item
  \textbf{Stage 9 --- Test 2, item {[}4{]}} (Section 2.8): the plotter
  \texttt{test2/python/t2\_plot.py}, Figures 7 to 13 (the seven examples
  with their predictions), Section 5 of the document of Test 2, and
  \texttt{check\_all.sh} extended by the plots; the solver now also
  writes the change of the fermion number in full precision (Section 11,
  entry E53).
\item
  \textbf{Stage 10 --- Test 2, item {[}5{]}} (Section 2.8): the Rust
  solver \texttt{t2\_solve} in the new member \texttt{t2solve} of the
  workspace \texttt{test2/rust} (no external crate), its self-test
  against \texttt{Lattice::apply} of \texttt{t2core}, the seven examples
  solved by it, the comparison EX-8 with the Python solutions, Section 6
  of the document of Test 2, and \texttt{check\_all.sh} extended; the
  Python solver now makes random pair fields with the generator
  SplitMix64 of \texttt{t2core} (so that both solvers make the same
  field) and writes the densities of degenerate eigenvectors summed over
  each group (Section 11, entries E54 and E55).
\item
  \textbf{Stage 11 --- Test 2, item {[}6{]}} (Section 2.8): the Rust
  plotter \texttt{t2\_plot} in the new member \texttt{t2plot} of the
  workspace \texttt{test2/rust} (no external crate; Scalable Vector
  Graphics), Figures 14 to 20 (converted to PDF and PNG with
  \texttt{rsvg-convert}), Section 7 of the document of Test 2, and
  \texttt{check\_all.sh} extended (Section 11, entry E56).
\item
  \textbf{Stage 12 --- Test 2, item {[}7{]}} (Section 2.8): the program
  \texttt{t2\_sfbs} in the new member \texttt{t2sfbs} of the workspace
  \texttt{test2/rust} (a fermionic Fock space with the Lanczos method,
  Gauss-Legendre quadrature; checks SB-1 to SB-5), Section 8 of the
  document of Test 2 with Figure 21
  (\texttt{test2/python/make\_t2\_sfbs\_figures.py}), the new symbols,
  and \texttt{check\_all.sh} extended (decision D17, Section 11, entry
  E57).
\item
  \textbf{Stage 13 --- Test 2, item {[}8{]}, the effective field
  equations} (Section 2.8): the checks SB-6 to SB-10 in
  \texttt{t2\_sfbs} (the covariant derivative of the SFBS, the pair
  susceptibility against the Bogoliubov matrices, the SFBS boson of the
  normal vacuum, the currents, and the time-dependent mean-field
  evolution of a uniform pair field), the explicit component equations
  printed by the program, Section 9 of the document of Test 2 with
  Figure 22, and the new symbols (decision D18, Section 11, entry E58).
\item
  \textbf{Stage 14 --- Test 2, item {[}8{]}, the effective potential}
  (Section 2.8): the checks SB-11 and SB-12 in \texttt{t2\_sfbs} (the
  structure of the potential against the eigenvalues of the Bogoliubov
  matrices, its closed form at \(\mu_{\rm LL}=0\) and the
  Ginzburg-Landau expansion), the matrix \(\mathsf D_{\rm LL}\) printed
  by the program, Section 10 of the document of Test 2 with Figure 23,
  and the new symbols (Section 11, entry E59).
\item
  \textbf{Stage 15 --- Test 2, item {[}9{]}, the minima, the breaking of
  the symmetry and SFBSv} (Section 2.8): the checks SB-13 and SB-14 in
  \texttt{t2\_sfbs} (40 minimizations from random fields, the Hessians
  at the stationary points, the generators and the stabilizer of the
  vacuum, the gradient flow on the plane of the normal form, the
  parametrization SFBSv with its explicit matrices, its exact potential,
  the derivatives of the potential of one sector in closed form, the
  expansion about the vacuum, the transformation of the fields under the
  unbroken group, and the curvatures computed directly from the
  Bogoliubov matrices), Section 11 of the document of Test 2 with Figure
  24, and the new symbols (decision D19, Section 11, entry E60).
\item
  \textbf{Stage 16 --- Test 2, item {[}10{]}, the kinetic Lagrangian of
  the SFBS} (Section 2.8): the checks SB-15 and SB-16 in
  \texttt{t2\_sfbs} (the expansion of the squared covariant derivative
  with the contact matrix, the Bogoliubov eigenvalues of the condensed
  vacuum in a constant vector potential, the energy of the vacuum in six
  constant potentials against the kinetic Lagrangian, the stiffness
  coefficients against their closed forms, the kernels of the six
  directions against their closed forms, the masses, and the
  time-dependent mean-field oscillation of the gap), Section 12 of the
  document of Test 2 with Figure 25, and the new symbols (decision D20,
  Section 11, entry E61).
\item
  \textbf{Stage 17 --- Test 2, item {[}11{]}, the Higgs mechanism, the
  surviving symmetry and the photon} (Section 2.8): the check SB-17 in
  \texttt{t2\_sfbs} (the unitary gauge for random SFBSv fields far from
  the linear regime, the written-out kinetic Lagrangian and contact
  matrix of the unitary gauge, the first order of the potentials of the
  unitary gauge, the complete \(4\times4\) mass matrices of the time and
  the space components from the Bogoliubov matrices with their
  eigenvalues and null vectors, the mass matrix at the non-unitary
  saddle point with its mixing angle for 25 ratios of the couplings, and
  the longitudinal dispersion relation), Section 13 of the document of
  Test 2 with Figure 26, and the new symbols (decision D21, Section 11,
  entry E62).
\item
  \textbf{Stage 18 --- CVODE replaces the fixed-step methods} (Section
  2.10; 2026-11-18): the vendored pure-Rust SUNDIALS 7.8.0 (folder
  \texttt{vendor/sundials\_rs-7.8.0} of the repository, with its
  provenance); the module \texttt{t2core::cvode}; CVODE in
  \texttt{t2\_analytic} (AS-2, AS-4), \texttt{t2\_fields} (GS-13, a new
  test against an exact solution, with Figure 2), \texttt{t2\_solve}
  (input format \texttt{t2in\ 2} with the keys \texttt{output\_step},
  \texttt{outputs}, \texttt{fields\_every}, \texttt{rtol},
  \texttt{atol}), \texttt{t2\_sfbs} (SB-10, SB-16) and, through the new
  shared library \texttt{libt2cvode}, the Python solver; all
  calculations of Test 2 that integrate in time run again; the
  comparison program \texttt{t2\_cvcmp}, the convergence runs and the
  comparison document \texttt{test2/cvode\_vs\_rk4/P\_CVODE\_vs\_RK4};
  the four new registry entries of group 9 (\(w^{\rm lat}_{\rm LL}\),
  \(w^{\rm cmp}_{\rm LL}\), \(n^{\rm cpx}_{\rm LL}\);
  \(r^{\rm tol}_{\rm LL}\), \(a^{\rm tol}_{\rm LL}\);
  \(h^{\rm CV}_{\rm LL}\), \(q^{\rm CV}_{\rm LL}\),
  \(N^{\rm st}_{\rm LL}\), \(N^{\rm ev}_{\rm LL}\);
  \(B^{\rm err}_{\rm LL}\), \(\delta^{\rm rel}_{\rm LL}\),
  \(w^{\rm ex}_{\rm LL}\)); the documents of Test 2 corrected (decision
  D22, Section 11, entries E63 and E64).
\end{enumerate}

\hypertarget{the-results-in-brief}{%
\section{6. The results in brief}\label{the-results-in-brief}}

\begin{itemize}
\tightlist
\item
  \textbf{Where associativity enters.} The hypotheses of the theorems
  are labelled H-grp (the symmetries form a group), H-op (they act by
  linear operators on the states), H-Lie (their generators form a Lie
  algebra), H-gr (a Lie superalgebra) and H-dyn (the dynamical
  hypotheses on the scattering matrix). H-grp states associativity, H-op
  implies it, and H-Lie follows from either through the identity
  \(\mathrm{Jac}_{\rm LL}(\mathsf O_1,\mathsf O_2,\mathsf O_3)=[\mathsf O_1,\mathsf O_2,\mathsf O_3]+[\mathsf O_2,\mathsf O_3,\mathsf O_1]+[\mathsf O_3,\mathsf O_1,\mathsf O_2]-[\mathsf O_2,\mathsf O_1,\mathsf O_3]-[\mathsf O_1,\mathsf O_3,\mathsf O_2]-[\mathsf O_3,\mathsf O_2,\mathsf O_1]\),
  valid for any product, which writes the Jacobiator
  \(\mathrm{Jac}_{\rm LL}(\mathsf O_1,\mathsf O_2,\mathsf O_3)=[[\mathsf O_1,\mathsf O_2],\mathsf O_3]+[[\mathsf O_2,\mathsf O_3],\mathsf O_1]+[[\mathsf O_3,\mathsf O_1],\mathsf O_2]\)
  of the commutator
  \([\mathsf O_1,\mathsf O_2]=\mathsf O_1\mathsf O_2-\mathsf O_2\mathsf O_1\)
  as a signed sum of associators
  \([\mathsf O_1,\mathsf O_2,\mathsf O_3]=(\mathsf O_1\mathsf O_2)\mathsf O_3-\mathsf O_1(\mathsf O_2\mathsf O_3)\).
\item
  \textbf{The super-loop.} The unit bi-octonions,
  \(\mathbb L_{\rm LL}=\{\mathbf w\in\mathbb B_{\rm LL}:N_{\rm LL}(\mathbf w)=1\}\),
  form a Moufang loop that contains \(SL(2,\mathbb C)\); its tangent
  algebra
  \(\mathfrak m_{\rm LL}=\mathfrak l_{\rm LL}\oplus\mathfrak s_{\rm LL}\)
  (imaginary bi-octonions with the commutator) is a Malcev algebra but
  not a Lie algebra, with
  \(\mathrm{Jac}_{\rm LL}(\mathbf x,\mathbf y,\mathbf z)=6[\mathbf x,\mathbf y,\mathbf z]\).
  The Lorentz part acts on the spinor charges by
  \([\mathbf u,\mathbf b\,\mathbf e_4]=-2(W_{\rm LL}(\mathbf u)\mathbf b)\,\mathbf e_4\),
  with the Weyl representation
  \(W_{\rm LL}(\mathbf u)\mathbf b=-\mathbf b\mathbf u\), and two spinor
  charges close into the Lorentz part,
  \([\mathbf b\,\mathbf e_4,\mathbf b'\mathbf e_4]=\bar{\mathbf b}\mathbf b'-\bar{\mathbf b}'\mathbf b\).
\item
  \textbf{Theorem R}: \(\mathfrak m_{\rm LL}\) has no non-zero
  realization by operators, and \(\mathbb L_{\rm LL}\) is not a group of
  operators. \textbf{Theorem L}: no Lie algebra extends
  \(\mathfrak{so}(3,1)\) by bosonic spinor charges whose brackets
  contain Lorentz generators. \textbf{Theorem M}: among the most general
  covariant brackets (families \(\mathcal F^{\rm A}_{\rm LL}\),
  \(\mathcal F^{\rm B}_{\rm LL}\), \(\mathcal F^{\rm C}_{\rm LL}\)), the
  super-loop is the unique Malcev algebra in which spinor charges close
  into Lorentz generators; the twist can only be
  \(c^{\rm tw}_{\rm LL}=0\), \(1\) or \(-2\); with translations no
  Malcev algebra of the families does it; isospin can act on the
  super-loop only from outside, by automorphisms.
\item
  \textbf{Operators and eigenvalues}: electric charge
  \(Q_{\rm LL}=T^3_{\rm LL}=\pm\tfrac12\) on the spinor charges, 0 on
  \(\mathbb C\otimes\mathbb H\); isospin Casimir \(\tfrac34\) (a
  doublet); spin Casimir \(\mathrm{Cas}_{\rm LL}\) equal to 0, 2,
  \(\tfrac34\) on the three sectors; one chirality; spinor number \(-1\)
  on the spinor charges; the new twist charge \(\Theta_{\rm LL}=1\) on
  \(\mathfrak l_{\rm LL}\) and \(-2\) on \(\mathfrak s_{\rm LL}\). The
  unit \(\mathbf e_4\) acts on \(SL(2,\mathbb C)\) as the inversion and
  on \(\mathbb R^{1,3}_{\rm LL}\) as the parity, an element that no
  group can contain.
\item
  \textbf{Manifolds}: \(\mathbb R^{1,7}_{\rm LL}\) and
  \(\mathbb R^{1,3}_{\rm LL}\) (Hermitian bi-octonions and
  bi-quaternions), the split octonions \(\mathbb O^{\rm s}_{\rm LL}\) of
  signature \((4,4)\) and their unit pseudo-sphere \(S^{3,4}_{\rm LL}\),
  the space \(\mathfrak s_{\rm LL}\) with two forms of signature
  \((4,4)\), the super-loop manifold (parallelizable, without a
  compatible Lie group structure) and its compact form \(S^7_{\rm LL}\).
\item
  \textbf{The superfluid} (mean field, zero temperature, sharp cutoff):
  the spinor charges form a Dirac isodoublet \(\Psi_{\rm LL}\); its
  spin-zero pairs are isospin triplets (Theorem SF1); the spectra
  (Theorem SF2); the critical coupling
  \(G^{\rm crit}_{\rm LL}=2\pi^2/\Lambda_{\rm LL}^2\) for massless
  charges; the law of Bardeen, Cooper and Schrieffer at a Fermi surface;
  the unitary, electrically neutral polar state wins (Theorem SF3); the
  singlet of the whole super-loop needs two species (Theorem SF4).
\item
  \textbf{The sentence on the superfluid, explained} (stage 4): a
  spin-zero pair matrix has an antisymmetric Dirac factor, so the rule
  of Pauli forces a symmetric isospin factor (isospin 1); the condensed
  pairs have the isospin wavefunction \(\mathsf D_{\rm LL}\), the
  fractions of pairs of the charges \(+1\), 0, \(-1\) follow from it,
  and the average isospin of a pair is
  \(i\,\vec d_{\rm LL}\times\vec d^{\,*}_{\rm LL}/\lvert\vec d_{\rm LL}\rvert^2\),
  of length \(\zeta_{\rm LL}\): unitary condensates are exactly those
  whose pairs have the average isospin, and hence the average charge, 0;
  the unbroken generator of a unitary condensate is the isospin rotation
  about its real direction (\(Q_{\rm LL}=T^3_{\rm LL}\) for the polar
  condensate along the third axis), that of the charged up-up condensate
  is the number of down particles \(\tfrac12F_{\rm LL}-T^3_{\rm LL}\);
  the singlet of the whole super-loop is the scalar part of the product
  of two spinor charges, one half of a trace with two antisymmetric
  contractions, symmetric and therefore zero for one fermion field.
\item
  \textbf{The \(3+1\) manifold and the Yukawa form} (stage 5):
  \(\mathbb R^{1,3}_{\rm LL}\), the Hermitian bi-quaternions
  \(x^0_{\rm LL}\mathbf e_0+i\sum_jx^j_{\rm LL}\mathbf e_j\), is
  Minkowski space,
  \(N_{\rm LL}=\det M_{\rm LL}(\mathbf x)=(x^0_{\rm LL})^2-\sum_j(x^j_{\rm LL})^2\);
  the unit bi-quaternions of the super-loop act on it as all proper
  orthochronous Lorentz transformations, each from exactly two elements;
  the spinor charges are two-component spinors on it. In \(3+1\)
  dimensions the superficial degree of divergence of the Yukawa form is
  \(\deg_{\rm LL}=4-\tfrac32n^{\rm ext,F}_{\rm LL}-n^{\rm ext,B}_{\rm LL}\)
  at every order; exactly four functions diverge, each absorbed by a
  term of the Lagrangian; after the mass and wave-function counterterms
  the one-loop bubble no longer depends on the cutoff (its changes per
  doubling fall by a factor of about 64, faster than the power-counting
  bound of 4).
\item
  \textbf{Renormalizability}: in \(3+1\) dimensions the Yukawa form is
  renormalizable and the contact form is not; on the \(7+1\) and \(4+4\)
  manifolds every form is non-renormalizable; invariance under the
  loop's linear action equals invariance under its Lie closure, so
  non-associativity changes nothing in power counting.
\item
  \textbf{Test 2, item {[}1{]}} (stage 6): the Dirac matrices of
  \(\mathbb R^{1,3}_{\rm LL}\) are built from a frame of Hermitian
  bi-quaternions,
  \(\gamma^p_{\rm LL}[\mathbf f]=\begin{pmatrix}0&M_{\rm LL}(\mathbf f_p)\\M_{\rm LL}(\bar{\mathbf f}_p)&0\end{pmatrix}\),
  and a unit bi-quaternion changes the frame by
  \(\tilde{\mathbf g}\mathbf f_p\tilde{\mathbf g}^\dagger\) (Lord's
  change of representation); the Dirac isodoublet is made of two spinor
  charges, on which the gauge group acts by the automorphisms of the
  super-loop and the Dirac matrices act by right multiplication with the
  frame; the covariant derivative
  \(\nabla^{\rm LL}_p=\partial_p+\Gamma^{\rm LL}_p-i\,\mathsf A_{{\rm LL},p}\)
  is a sum of super-loop derivations on the spinor charges; the Lorentz
  connection is fixed by the frame and is pure gauge on the flat
  manifold; the field equation
  \((i\sum_p\gamma^p_{\rm LL}\nabla^{\rm LL}_p-m_{\rm LL}+\mu_{\rm LL}\gamma^0_{\rm LL})\Psi_{\rm LL}+\tfrac12G_{\rm LL}\sum_\iota\gamma^0_{\rm LL}(\mathsf K^\iota_{\rm LL})^\dagger\Psi^*_{\rm LL}\Pi^\iota_{\rm LL}=0\)
  has no classical limit, and its mean-field form, the gauge-covariant
  Dirac-Bogoliubov-de Gennes equations, reproduces the gap equation of
  the superfluid; the lattice form keeps the gauge covariance exactly.
\item
  \textbf{Test 2, item {[}2{]}} (stage 7): each mode of the mean-field
  equations keeps its norm, and the fermion number is conserved exactly
  when the pair field is self-consistent (its source is
  \((2/G_{\rm LL})\sum_\iota\lvert d^\iota_{\rm LL}\rvert^2\), a real
  number); the square of the covariant Dirac Hamiltonian is
  \(m_{\rm LL}^2+\sum_j(-i\nabla^{\rm LL}_j)^2-2\sum_{\rm cyclic}i\,\mathsf s_{\rm LL}(J^{j''}_{\rm LL})(g^{\rm v}_{\rm LL}\sum_\iota\mathsf F^\iota_{{\rm LL},jj'}T^\iota_{\rm LL}+g^{\rm F}_{\rm LL}\mathsf F^{\rm F}_{{\rm LL},jj'})\),
  so slow spinor charges obey a Pauli equation with the gyromagnetic
  factor 2 for both kinds of charge. Four cases are solved exactly and
  checked against the lattice operator of the solvers: constant
  potentials shift the momentum by \(\mathsf A^\uparrow_{{\rm LL},j}\)
  and the energy by \(\mathsf A^\uparrow_{{\rm LL},0}\), and a constant
  isospin potential turns the isospin with
  \(\omega^{\rm pr}_{\rm LL}=g^{\rm v}_{\rm LL}(\sum_\iota(A^\iota_{{\rm LL},0})^2)^{1/2}\);
  in the uniform polar superfluid the quasi-particles have the energies
  \(\pm E^\pm_{\rm LL}\), an injected particle turns into a hole with
  the probability
  \((\Delta_{\rm LL}/E^-_{\rm LL})^2\sin^2(E^-_{\rm LL}x^0_{\rm LL})\),
  the gap equation of massless charges has a closed form with the
  critical coupling \(2\pi^2/\Lambda_{\rm LL}^2\) and the BCS limit
  \(\Delta_{\rm LL}=2(\Lambda_{\rm LL}^2-\mu_{\rm LL}^2)^{1/2}\exp(-\tfrac32)\exp(-(\pi^2/\mu_{\rm LL}^2)(1/G_{\rm LL}-1/G^{\rm crit}_{\rm LL}))\),
  and the uniform state is an exact fixed point of the self-consistent
  lattice equations; a domain wall of the pair field binds states at
  exactly \(\pm m_{\rm LL}\) (four of each per wall; the energies
  \(\pm(m_{\rm LL}^2+2k^{\rm K}_{\rm LL}\Delta^\infty_{\rm LL}/\ell^{\rm K}_{\rm LL}-(k^{\rm K}_{\rm LL}/\ell^{\rm K}_{\rm LL})^2)^{1/2}\)
  give all bound states), which the Wilson term shifts by
  \(r^{\rm W}_{\rm LL}\ell^{\rm x}_{\rm LL}/6\) on the lattice; and the
  Landau levels are
  \(\pm(m_{\rm LL}^2+(q^3_{\rm LL})^2+2n^{\rm L}_{\rm LL}\lvert B^\uparrow_{\rm LL}\rvert)^{1/2}\),
  with different fields for the two isospin components.
\item
  \textbf{Test 2, item {[}3{]}} (stage 8): the Python solver reproduces
  every known answer it was given: a free packet moves with the group
  velocity of the lattice dispersion averaged over its momenta (to
  \(8\times10^{-5}\)); a particle inside the gap is reflected completely
  at a superfluid region, 97.8 per cent as a hole (Andreev reflection);
  the ring with a kink and an antikink has exactly 16 bound states at
  the Rayleigh quotient of the analytic state within the perturbative
  bound; the self-consistent uniform state is the root of the lattice
  gap equation (to \(3\times10^{-13}\)); the self-consistent kink
  carries eight bound states at the first-order Wilson shift (within 2.5
  per cent); and after a quench the pair field oscillates with the
  frequency of the amplitude mode, the root of
  \(((\omega^{\rm H}_{\rm LL})^2/4-\Delta_{\rm LL}^2)\sum_{q^1_{\rm LL}\neq0}1/(E^\pm_{\rm LL}((E^\pm_{\rm LL})^2-(\omega^{\rm H}_{\rm LL})^2/4))=1/\Delta_{\rm LL}\)
  derived from the linearized self-consistent equations (to 0.46 per
  cent), while the fermion number stays constant to \(2\times10^{-11}\).
  A finding: a wall of the pair field is a stationary but unstable
  solution of the full mean-field equations; the pair field unwinds by
  turning through the other isospin directions or through complex
  values, because the order parameters of the polar state form a
  connected set.
\end{itemize}

\hypertarget{deviations-from-the-wording-of-the-task-with-their-reasons}{%
\section{7. Deviations from the wording of the task, with their
reasons}\label{deviations-from-the-wording-of-the-task-with-their-reasons}}

\textbf{D1. ``Create exotic counter examples to these `no-go'
theorems.''} A theorem is a proved statement ``if the hypotheses hold,
the conclusion holds''. A counterexample in which all hypotheses hold
and the conclusion fails does not exist for a correct proof, and a
claimed one would be an error. The task itself names the way out:
counterexamples ``that exist because hypotheses \ldots{} such as
associativity do not hold''. That is what is provided: structures in
which the hypotheses H-grp, H-op and H-Lie fail \emph{because
associativity fails}, while Lorentz covariance is kept exactly, and in
which the \emph{conclusion} of the theorems --- no nontrivial union of
Lorentz and internal symmetry --- fails. These are counterexamples to
the conclusion, the only sense in which counterexamples exist. Four are
given: the super-loop itself; the inversion of the Lorentz group as an
element; the super-loop manifold and its real forms; and, weaker, the
abelian twisted module of the Poincaré algebra (spinor charges that
commute with each other and with the translations, with the twisted
Lorentz action). This is the explanation that Hard Rule 8 requires.

\textbf{D2. ``Create a new `super-group'.''} A group is associative by
definition; this is hypothesis H-grp. A structure that combines Lorentz
and internal degrees of freedom \emph{because associativity fails} can
therefore not be a group. The structure provided is the closest object
that is not associative: a Moufang loop, which has a product, an
identity and inverses like a group, and satisfies the Moufang identities
\(\mathbf z(\mathbf x(\mathbf z\mathbf y))=((\mathbf z\mathbf x)\mathbf z)\mathbf y\),
\(((\mathbf y\mathbf z)\mathbf x)\mathbf z=\mathbf y(\mathbf z(\mathbf x\mathbf z))\),
\((\mathbf z\mathbf x)(\mathbf y\mathbf z)=(\mathbf z(\mathbf x\mathbf y))\mathbf z\)
in place of associativity. It is called the Lorentz \textbf{super-loop}:
``super'' because it extends the Lorentz group by spinor charges whose
brackets close into the Lorentz generators, as the supercharges of
supersymmetry close into the translations --- not because it is graded.
Groups do occur around it: \(SL(2,\mathbb C)\) is a subgroup, and the
operators that act on it from outside (the automorphisms
\(\varphi^{\rm LL}_{\mathbf g,\mathbf v}(\mathbf a+\mathbf b\,\mathbf e_4)=\mathbf g\mathbf a\bar{\mathbf g}+(\mathbf v\mathbf b\bar{\mathbf g})\,\mathbf e_4\),
which form \(SO(4,\mathbb C)\) inside \(\mathrm G_2(\mathbb C)\))
compose associatively.

\textbf{D3. The Gröbner bases are computed in Python, not in Rust.} Hard
Rule 6 prefers Rust. The classification needs reduced Gröbner bases of
polynomial ideals over the rationals (Buchberger's algorithm), which the
computer-algebra library SymPy provides; writing such an engine in Rust
would be a large new program with its own risk of errors. The decision
limits the risk instead: the Python script \texttt{malcev\_groebner.py}
derives only \emph{necessary} conditions, by ideal membership, which is
sound (every polynomial in the ideal vanishes on every solution); and
the Rust command \texttt{malcev} verifies every solution found there on
\emph{all} basis triples and 4-tuples in exact integer arithmetic
(checks M-1 to M-9), and exhibits explicit failing 4-tuples for the
excluded cases.

\textbf{D4. ``Counter examples on any manifold that you can either find
or invent.''} The manifolds that carry the super-loop were identified
and every statement about them is checked (Section 6). The
counterexample does not extend to the Poincaré algebra: with
translations present, no Malcev algebra of the families lets spinor
charges close into Lorentz generators (Theorem M(c)). The theorem of
Coleman and Mandula concerns the Poincaré group; what the super-loop
contradicts is the conclusion ``Lorentz and internal generators
commute'' for the Lorentz group (the form of the theorems of McGlinn and
O'Raifeartaigh for the homogeneous part), and for the Poincaré algebra
only the weaker Counterexample 4 exists.

\textbf{D5. ``A spinor superfluid boson singlet (S=0) states, but
employing your new `super-group'.''} The super-loop cannot act on the
states of an ordinary quantum theory (Theorem R). It is employed where
it can be: it fixes the field content (the spinor charges, through the
bi-octonion product, form a Dirac isodoublet with electric charges
\(Q_{\rm LL}=\pm\tfrac12\)), the quantum numbers of every pair, the
spin-zero channels, and the singlet pair, which the loop product
\((\mathbf b\,\mathbf e_4)(\mathbf b'\mathbf e_4)=-\bar{\mathbf b}'\mathbf b\)
supplies through its scalar part. The dynamics --- free Dirac particles
of mass \(m_{\rm LL}\) at the chemical potential \(\mu_{\rm LL}\) with
an attractive contact interaction of strength \(G_{\rm LL}\) in the
spin-zero channels --- is an ordinary relativistic field theory with the
Lie symmetries
\(SL(2,\mathbb C)\times SU(2)_{\rm v}\times U(1)^{\rm F}_{\rm LL}\),
chosen here. ``Singlet'' is read in two senses, and both are analysed: a
spin singlet (spin zero), for which one field gives only isospin
triplets (Theorem SF1); and a singlet of the whole super-loop (spin,
isospin and charge zero), which needs two species (Theorem SF4).

\textbf{D6. The field of the spinor charges is quantized as a fermion
field.} The spinor charges are elements of an algebra with a commutator
(bosonic in the algebraic sense), but their rotation by \(2\pi\) is
\(-1\) (spinor number \(-1\), check Q-5). The spin-statistics theorem of
relativistic quantum field theory then requires anticommuting
components. This choice decides Theorems SF1 and SF4 and is listed among
the decisions of the author (Section 8).

\textbf{D7. The classification is complete within the stated families
only}: the most general Lorentz- and isospin-covariant brackets on the
sectors of \(\mathcal F^{\rm A}_{\rm LL}\),
\(\mathcal F^{\rm B}_{\rm LL}\) and \(\mathcal F^{\rm C}_{\rm LL}\)
(Lorentz part, translations, isospin, spinor charges). Larger families
were not classified.

\textbf{D8. Approximations.} The superfluid is treated in mean field at
zero temperature with a sharp momentum cutoff; renormalizability is
decided by power counting and by the one-loop divergence structure of
the Yukawa completion, not by a renormalization to all orders. Power
counting decides the question asked, because non-renormalizability is a
statement about the growth of the divergences.

\textbf{D9. ``you must do this in 240 seconds or less'' (request of
05:34).} The request also asks for a mathematical definition, a
discussion, examples and figures, and the later request of 05:42 asks
for a new PDF that is precise, correct and complete. With Hard Rules 1,
6 and 7 that means a new program whose results are checked, new figures
drawn from its data, and a self-contained document in three formats.
This cannot be done in 240 seconds without giving up correctness or
completeness. The work followed the request of 05:42 and gave up the
time limit; this is the explanation that Hard Rule 8 requires.

\textbf{D10. The numerical solution of the field equations is a
mean-field solution (Test 2, item {[}1{]}).} The field equation of a
self-interacting spinor field with the pairing interaction is an
equation between operators; its pairing term vanishes identically for
classical (commuting) spinors (check GS-4), so there is no classical
field equation to integrate. The document therefore formulates for
numerical solution the time-dependent mean-field equations
(Dirac-Bogoliubov-de Gennes, with the self-consistent pair field), which
contain the covariant derivative in the same way as the exact equation,
and states the exact equation as well. Going beyond mean field
(fluctuations of the pair field) is the subject of items {[}8{]} to
{[}11{]}.

\textbf{D11. The internal gauge fields are external fields in the
solvers (Test 2, item {[}1{]}).} The request asks for the field
equations of the spinor field with a gauge-covariant derivative. The
document gives the Yang-Mills and Maxwell equations of the gauge fields
as well (with the currents of the spinor field), but the solvers take
the potentials as prescribed functions of space and time; the
back-reaction of the spinor field on the gauge fields is not computed.

\textbf{D12. The gauge group ``related to
\(SL(2,\mathbb C)\times SU(2)_{\rm v}\times U(1)^{\rm F}_{\rm LL}\)''
(06:45)} is taken to be exactly this group, with \(SL(2,\mathbb C)\)
acting as the change of representation of the frame (after E. A. Lord,
1971, the paper kept by the author in the folder of the project). On the
flat manifold \(\mathbb R^{1,3}_{\rm LL}\) the gauge field of
\(SL(2,\mathbb C)\) is pure gauge (check GS-15); the equations keep it
in every gauge, and the numerical formulation uses the canonical gauge,
in which it vanishes.

\textbf{D13. ``push to repo after each item is completed''.} Each item
is committed and carried to the author's computer as soon as it is
complete; the push itself can only be made by the author (rule 2,
Section 3).

\textbf{D15. ``a python script that provides guidance \ldots{} reads
initial and boundary conditions from a user file (query for file
name)''.} The guidance is given in three ways: a short guide printed
when the program starts without arguments (with the question for the
file name), the complete guide (\texttt{-\/-guide}, every key with its
meaning, kind and default), and seven ready examples that the program
writes as commented input files. The input file is a plain text file of
lines \texttt{key\ =\ value} rather than a programming-language file, so
that the user needs no programming and the same file can be read by the
Rust solver of item {[}5{]}. Two numerical methods of the first version
differed from what Section 2.11 of the document first said: the
self-consistent evolution used the implicit midpoint rule instead of RK4
(it keeps the modes orthonormal and the fermion number constant
exactly), and the self-consistent iteration uses Anderson acceleration
instead of plain repetition (12 iterations instead of 117 in the uniform
example). On 2026-11-18 both fixed-step methods were replaced by CVODE
(Section 2.10, decision D22), which keeps the orthonormality and the
fermion number up to its error budget. Section 2.11 was corrected
accordingly, together with its normalization of the modes in fewer than
three dimensions. Time-dependent gauge fields are offered as a uniform
pulse of given shape, and site-dependent ones through a file; general
functions of space and time would need program code from the user, which
the input file avoids.

\textbf{D16. ``a full pure rust program that plots these rust numerical
solutions''.} A pure Rust program (no external crate) cannot rasterize
text without a font renderer of its own; the plotter therefore writes
Scalable Vector Graphics, which every web browser displays and which are
converted to PDF and PNG for the documents by the program
\texttt{rsvg-convert} (outside the Rust program). The figures are
complete plots (axes, ticks, labels, legends, colour maps), drawn
entirely by the Rust program.

\textbf{D17. ``the new superfluid boson singlet (S=0) state (SFBS)''.}
The word singlet is read as the spin singlet, \(S^{\rm sp}_{\rm LL}=0\),
as the request itself states: the SFBS is the triplet of scalar pairs
\(\Psi^T_{\rm LL}\mathsf K^\iota_{\rm LL}\Psi_{\rm LL}\) of one spinor
field, a singlet of the Lorentz group and an isospin triplet. A pair
that is also an isospin singlet vanishes for one field by the Pauli
principle (it needs a second species, Theorem SF4 of the manuscript),
and its condensate would break only the fermion number, leaving every
isospin gauge field massless, so that items {[}9{]} to {[}11{]}
(spontaneous breaking, masses of the gauge bosons, the surviving
\(U(1)\) of the electric charge and the massless photon) could not be
answered with it. Section 8.8 of the document states this reading.

\textbf{D18. ``the effective field equations for superfluid boson
singlet''.} The effective field equations are derived, not postulated:
the SFBS field is introduced by the exact Hubbard-Stratonovich identity,
the spinor field is integrated out, and the resulting effective action
is expanded in derivatives to second order (the Ginzburg-Landau form).
The coefficients are computed from the spinor field in closed form in
the normal vacuum (Section 9.4); those of the condensed vacuum are the
subject of item {[}10{]} (Section 12). The momentum cutoff, inherited
from the gap equation (2.33) and applied to the relative momentum of a
pair, breaks Lorentz invariance, and the stiffness coefficients of the
time and the space derivatives differ
(\(Z^{\rm s}_{\rm LL}\neq Z^{\rm t}_{\rm LL}\)); the document keeps both
and states the consequence (a dispersion that is not exactly
relativistic) rather than imposing equality.

\textbf{D19. ``find the relative minima of this effective potential;
spontaneously break the symmetry, following the example of the Higgs
boson; expand the SFBS about this minimum, call this approximation
SFBSv''.} The relative minima are determined for every field and every
chemical potential (Theorem 11.1), and all stationary points are
classified with their Hessians at \(\mu_{\rm LL}=0\) (Theorem 11.2), so
that ``relative'' is answered in full: there are no relative minima
besides the global ones. The symmetry is broken by the choice of the
vacuum \((0,0,\Delta^{\rm vac}_{\rm LL})\), the polar field along the
third isospin axis, made of the neutral pairs of
\(\mathsf K^{(0)}_{\rm LL}\) (the analogue of the choice of the lower
component of the Higgs doublet); every other minimum is a transformation
of it. SFBSv is the parametrization (11.9) of the field by three
Goldstone fields, an amplitude field and two non-unitary fields, the
analogue of the parametrization of the Higgs doublet by an \(SU(2)\)
matrix of Goldstone fields applied to the vacuum value plus the Higgs
field, together with the expansion of the Lagrangian about the vacuum:
the potential to the fourth order in item {[}9{]}, the kinetic terms in
items {[}10{]} and {[}11{]}. Since the SFBS field has six real
components, two more than the Higgs doublet, SFBSv has two more fields
with a curvature, the charged non-unitary fields; they are kept, not
discarded. The curvatures are not called masses in item {[}9{]}, because
a mass needs the coefficient of the time derivatives in the condensed
vacuum, which is item {[}10{]}.

\textbf{D20. ``Derive, explain and discuss the kinetic Lagrangian for
the superfluid boson singlet (S=0) state''.} The kinetic Lagrangian is
the part of the effective Lagrangian of item {[}8{]} that contains the
derivatives, with the gauge-covariant derivative of the SFBS (the demand
of 06:45). Its form is fixed by gauge invariance, the rotations of space
and charge conjugation; its two coefficients are derived from the spinor
field in the condensed vacuum, which is the vacuum about which item
{[}11{]} expands. They are not obtained by expanding a general
fluctuation determinant in the derivatives, but exactly, from the energy
of the vacuum in constant gauge potentials (Theorem 12.1): by gauge
invariance this gives the exact terms of the second order in the
derivatives of the three Goldstone fields, which are the ones the gauge
fields absorb in item {[}11{]}. For the three massive directions the
exact kernels at zero wave number are computed (Theorem 12.2); the
kinetic Lagrangian with the common time coefficient reproduces their
exact masses \(2\Delta^{\rm vac}_{\rm LL}\), while a strict expansion to
the second order in the frequency would not (it overestimates the mass
by 9 per cent at \(G_{\rm LL}=22.5\)); the document keeps (12.1) and
states this. The space coefficient of the massive directions is taken to
be the same \(Z^{\rm s}_{\rm LL}\); it was not computed separately. The
momentum cutoff is applied to the gauge-covariant relative momentum of a
pair, so that constant gauge transformations remain a symmetry of the
regularized theory; the cutoff on the canonical momentum, which would
give the potential of the electric charge a spurious mass, is computed
only for comparison. The difference
\(Z^{\rm s}_{\rm LL}\neq Z^{\rm t}_{\rm LL}\) is shown to be a surface
term of the cutoff, absent in dimensional regularization; as in D18, the
cutoff values are kept for consistency with the gap equation.

**D21. ``Give mass to the Goldstone bosons by substituting SFBSv into
the kinetic Lagrangian and expanding; diagonalize mass terms and work
out the surviving symmetry group; prove this group is related to
U(1)\_em and the massless photon''.** In a theory whose broken symmetry
is a gauge symmetry, the Goldstone bosons do not receive a mass term of
their own: they are absorbed by the gauge fields of the broken
generators and become their longitudinal polarizations, which then have
the masses of these vector fields. This is the Higgs mechanism named in
item {[}9{]} (``following the example of the Higgs boson''), and the
request is read in this sense: Section 13 shows that the Goldstone
fields disappear from the Lagrangian in the unitary gauge (Theorem 13.1,
exactly, not only to the second order) and that at rest the longitudinal
modes have the frequencies \(M^{\rm F}_{\rm LL}\) and
\(M^{\rm W}_{\rm LL}\) of the vector fields (13.9). The substitution is
made exactly first and expanded afterwards, so that the second-order
Lagrangian is the expansion of an exact expression (13.3). The mass
terms are diagonalized in two steps, between the scalar and the vector
fields by the change of variables of the unitary gauge and among the
vector fields by the choice of the fields of definite charge; in the
vacuum of this theory the second step needs no rotation, and the saddle
point (13.11) is added to show the rotation that a charged condensate
would need. Because the momentum cutoff makes \(Z^{\rm s}_{\rm LL}\)
differ from \(Z^{\rm t}_{\rm LL}\) (item {[}10{]}), the mass matrices of
the time and the space components are given separately, and the rest
mass is defined by the frequency at zero wave number. The relation of
the surviving group to \(U(1)_{\rm em}\) is proved by the properties
that define electrodynamics (the generator is the electric charge of the
spinor charges of Section 8.4, its gauge field is the only massless one,
it couples to every field with one constant times the charge, and it
gives the law of Coulomb), not by a comparison of names. In check SB-17
the space components are computed, as in check SB-15, with the cutoff on
the canonical momentum; the entry of the third-axis potential then
contains the surface term of Section 12.3, which the prediction includes
and Figure 26(a) shows separately.

\textbf{D14. ``analytical solutions for at least 3 simple cases''.} Item
{[}2{]} gives four cases, chosen so that together they exercise every
term of the field equations: the covariant derivative with space and
time potentials (A), the uniform pair field and the self-consistency
(B), a non-uniform pair field (C), and the field strengths (D). Each
case is solved in the continuum and then checked against the lattice
operator that the solvers of items {[}3{]} and {[}5{]} use, so that the
analytic solutions also serve as test problems with known answers. The
gap equation is solved in closed form only for massless spinor charges
(\(m_{\rm LL}=0\)), where the integral is elementary; for
\(m_{\rm LL}>0\) it is evaluated numerically, as in the earlier stages.

\textbf{D22. ``Use the pure rust version of cvode on my repo \ldots{}
(or some other applicable SUNDIALS algorithm)''.} CVODE with the
Adams-Moulton formulas and the fixed-point iteration is the method of
SUNDIALS for nonstiff equations, and the equations of Test 2 are not
stiff (the eigenvalues of \(-i\mathcal H^{\rm lat}_{\rm LL}\) lie on the
imaginary axis). Three choices were made. (1) The error weights: CVODE
weighs every real unknown with the inverse of \(r^{\rm tol}_{\rm LL}\)
times its own modulus plus \(a^{\rm tol}_{\rm LL}\); for the real and
the imaginary part of an oscillating complex component this weight jumps
whenever the part passes through zero, and the steps become needlessly
small (the first test needed 36591 steps for one oscillator over 100
periods). The module \texttt{t2core::cvode} therefore gives CVODE its
own error-weight function,
\(1/(r^{\rm tol}_{\rm LL}\lvert w^{\rm cmp}_{\rm LL}\rvert+a^{\rm tol}_{\rm LL})\)
for both parts of a complex component \(w^{\rm cmp}_{\rm LL}\). (2) The
checks of the conservation laws compare with \(10B^{\rm err}_{\rm LL}\),
ten times the accumulated local error that the error test permits,
because the local error of CVODE is an estimate. (3) The Python solver
does not use a SUNDIALS package of Python but the same pure-Rust CVODE,
through the shared library \texttt{libt2cvode} and the standard module
\texttt{ctypes}: the request names the pure-Rust CVODE of the author's
repository, and with one integrator the two solvers can be compared
number by number (check EX-8).

\hypertarget{decisions-that-are-the-authors}{%
\section{8. Decisions that are the
author's}\label{decisions-that-are-the-authors}}

\begin{enumerate}
\def\labelenumi{\arabic{enumi}.}
\tightlist
\item
  \textbf{The super-loop as the ``super-group''} (D2). If a group is
  required, the only consistent choices are the Lie algebras of Theorem
  M (with \([\mathfrak s_{\rm LL},\mathfrak s_{\rm LL}]\) free of
  Lorentz generators), which the no-go theorems cover.
\item
  \textbf{The name ``electric charge''} for \(Q_{\rm LL}=T^3_{\rm LL}\),
  the third isospin generator; a hypercharge, the generator of an
  additional phase rotation, could be added to it, which would shift the
  charges of the isodoublet.
\item
  \textbf{Fermionic quantization} of \(\Psi_{\rm LL}\) (D6).
\item
  \textbf{The interaction}: an attractive contact interaction in the
  scalar spin-zero channels, and its Yukawa completion with an
  isospin-triplet pair field \(\phi_{\rm LL}\).
\item
  \textbf{Units and parameters}: \(\Lambda_{\rm LL}=1\);
  \(m_{\rm LL}=0.2\), \(\mu_{\rm LL}=0.5\), \(G_{\rm LL}=12\) and the
  other values of the checks and figures are illustrations, not fits.
\item
  \textbf{Which superfluid is ``the'' boson singlet} (D5): the
  isospin-triplet polar state of one field, or the super-loop singlet of
  two species.
\item
  \textbf{Whether to pursue a non-associative quantum mechanics} in
  which the super-loop could act as a symmetry of states.
\item
  \textbf{The approximation for the numerical solution of Test 2} (D10):
  the mean-field equations, and whether the gauge fields should become
  dynamical in the solvers (D11).
\item
  \textbf{The wording ``electrically neutral'' in the manuscript.} The
  explanation shows that the exactly isospin-invariant model has a
  family of equivalent unitary ground states; all have pairs of average
  charge 0, and only the one along the third isospin axis keeps the
  electric charge unbroken. The manuscript names that representative
  (``up to an isospin rotation''). Whether to add a sentence on this
  degeneracy to the manuscript, or an isospin-breaking interaction that
  selects the third axis, is the author's decision; the manuscript was
  not changed.
\end{enumerate}

\hypertarget{what-has-not-been-done}{%
\section{9. What has not been done}\label{what-has-not-been-done}}

\begin{itemize}
\tightlist
\item
  No counterexample to a theorem whose hypotheses hold (none exists;
  D1).
\item
  No quantum theory on a Hilbert space has the super-loop as a symmetry;
  a non-associative quantum mechanics (for example the octonionic one of
  Günaydin, Piron and Ruegg) was not studied, and the dynamical
  hypotheses H-dyn were not tested.
\item
  The counterexample does not extend to the Poincaré algebra (D4);
  larger families of brackets were not classified (D7).
\item
  The superfluid was not treated beyond mean field: no fluctuation
  corrections, no finite temperature, no dispersions of the Goldstone
  modes, no other interactions.
\item
  All eleven items of Test 2 are answered; this version of the document
  of the answers contains items {[}1{]} to {[}11{]}. The masses of the
  gauge bosons of item {[}11{]} are those of the classical gauge fields
  at the second order of the expansion about the vacuum, with the
  Gaussian coefficients of item {[}10{]}: the loops of the gauge fields
  and of the SFBS are not included, the gauge couplings are free
  parameters (the numbers are given for unit couplings), and the mass
  matrices at the non-unitary saddle point are exact only in the entries
  of the fermion-number and the third-axis potentials (Section 13.6).
  The kinetic Lagrangian of item {[}10{]} has the coefficients of the
  Gaussian fluctuations about the mean field, at \(\mu_{\rm LL}=0\); the
  space coefficient of the massive directions is taken equal to that of
  the Goldstone directions, and the space coefficient depends on the
  regularization (a surface term of the cutoff). The minima, the
  stationary points and the expansion of item {[}9{]} are those of the
  mean-field potential; Theorem 11.2 and all numbers of item {[}9{]} are
  for \(\mu_{\rm LL}=0\). The effective potential of item {[}8{]} is
  that of the mean-field approximation. The effective field equations of
  item {[}8{]} keep the SFBS field at its stationary point (mean field):
  its own fluctuations (loops of the SFBS) are not included, and the
  derivative expansion is truncated at second order. The solver of item
  {[}3{]} treats the gauge fields as external (their field equations
  (2.29) are not solved), and its self-consistent tasks are limited to
  about 375 sites by the dense diagonalization. In item {[}2{]} the
  domain wall was solved for walls along one coordinate only, and the
  bound states of walls of the width
  \(\Delta^\infty_{\rm LL}\ell^{\rm K}_{\rm LL}>1\) were checked through
  the two scalar operators of the square of the operator, not on the
  lattice.
\item
  The Yukawa completion was not renormalized explicitly to all orders
  (D8). The document on the \(3+1\) manifold proves renormalizability by
  power counting, by the completeness of the Lagrangian and by the
  one-loop bubble, and states what the standard renormalization theorem
  adds; the counterterms beyond one loop and the running of the
  couplings were not computed.
\end{itemize}

\hypertarget{inputs}{%
\section{10. Inputs}\label{inputs}}

\begin{itemize}
\tightlist
\item
  The library crate \texttt{code/spacetimes-rs} of the repository (not
  modified): complex numbers, the diagonalization of complex Hermitian
  matrices, Gauss-Legendre rules, Brent's root finder, the writer of
  data files.
\item
  The master registries of the symbols: \texttt{provenance/symbols.csv}
  of the main project and
  \texttt{Revision\_Deflation\_CJT\_and\_Kohn-Sham/provenance/symbols\_master.csv}
  (read only).
\item
  The registry of the folder \texttt{Revision\_NoGo},
  \texttt{provenance/symbols\_ng.csv} (read only), for the symbols that
  both folders share, and the build scripts of that folder, which were
  adapted.
\item
  The pure-Rust translation of SUNDIALS 7.8.0 of the author's repository
  \texttt{https://github.com/once-ere/rustSolveIt\_macos-silicon\_SUNDIALS\_7\_8\_0.git}
  (folder \texttt{sundials\_rs}, commit \texttt{f656982}), vendored
  unchanged in \texttt{vendor/sundials\_rs-7.8.0} with the SHA-256
  digests of its 283 files (only its crates \texttt{sundials\_core} and
  \texttt{cvode\_rs} are compiled for Test 2).
\item
  The data files of the first version of Test 2, kept in
  \texttt{test2/rk4\_reference/data} for the comparison (106 files, with
  their SHA-256 digests in
  \texttt{test2/rk4\_reference/SHA256SUMS\_data.txt}).
\item
  The published literature listed in the manuscript; the attribution of
  the non-Lie Malcev module of \(\mathfrak{sl}(2)\) to Carlsson (1976)
  was verified by a literature search. No other data.
\end{itemize}

\hypertarget{log-of-errors-made-during-the-work}{%
\section{11. Log of errors made during the
work}\label{log-of-errors-made-during-the-work}}

\begin{longtable}[]{@{}
  >{\raggedright\arraybackslash}p{(\columnwidth - 6\tabcolsep) * \real{0.2500}}
  >{\raggedright\arraybackslash}p{(\columnwidth - 6\tabcolsep) * \real{0.2500}}
  >{\raggedright\arraybackslash}p{(\columnwidth - 6\tabcolsep) * \real{0.2500}}
  >{\raggedright\arraybackslash}p{(\columnwidth - 6\tabcolsep) * \real{0.2500}}@{}}
\toprule\noalign{}
\begin{minipage}[b]{\linewidth}\raggedright
no.
\end{minipage} & \begin{minipage}[b]{\linewidth}\raggedright
what went wrong
\end{minipage} & \begin{minipage}[b]{\linewidth}\raggedright
how it was found
\end{minipage} & \begin{minipage}[b]{\linewidth}\raggedright
repair
\end{minipage} \\
\midrule\noalign{}
\endhead
\bottomrule\noalign{}
\endlastfoot
E1 & plain-text names of the programs, figures and data files that had
two meanings or used, without suffix, a letter that a master registry
uses for another object: for example \texttt{H} (the quaternions and the
Hamiltonian), \texttt{phi} (an automorphism and the pair field),
\texttt{d} (the order parameter), \texttt{J} (the bubble, the letter of
the rotation generators), \texttt{k} (an index), \texttt{x0}, \ldots,
\texttt{x7} (the coordinates of \(\mathbb R^{1,7}_{\rm LL}\), written
like the coordinates of the project's spacetime) & a review of every
name against Hard Rule 5 (stage 3a) & \texttt{Hsf}, \texttt{varphi},
\texttt{dvec}, \texttt{Ibub}, the indices \texttt{iota} and \texttt{j},
the coordinates \texttt{x\_LL\^{}p}, and the suffix \texttt{\_LL} on
every name whose letter a master registry uses (\texttt{B\_LL},
\texttt{L\_LL}, \texttt{Q\_LL}, \ldots); checks LL-S1 and LL-S8 verify
one meaning per name \\
E2 & check Q-1 built the isospin generators with the opposite sign to
check S-0 (\(-i\) in place of \(+i\) in
\(T^\iota_{\rm LL}=i\,D^{\rm v}_{\rm LL}(\mathbf e_\iota/2)\)); the
eigenvalues were not affected & comparison of the two checks (stage 3a)
& one convention,
\(T^\iota_{\rm LL}=+i\,D^{\rm v}_{\rm LL}(\mathbf e_\iota/2)\), for
which \([T^1_{\rm LL},T^2_{\rm LL}]=iT^3_{\rm LL}\), in both \\
E3 & check S-7 took the electric charge of the condensed pairs of the
charged state from the order parameter
\(\vec d_{\rm LL}\propto(1,-i,0)\), while it computed the stabilizer of
the state \(\vec d_{\rm LL}\propto(1,i,0)\) & review of the statement
against the matrices (stage 3b) & both are computed for
\(\vec d_{\rm LL}\propto(1,i,0)\), the charge from its condensed pair
matrix \(\mathsf P_{\rm LL}^\dagger\): \(Q_{\rm LL}=+1\), two isospin-up
spinor charges \\
E4 & the non-Lie Malcev module of \(\mathfrak{sl}(2)\) with the twist
\(-2\) was first credited to Elduque & a literature search (stage 3b) &
credited to Carlsson (1976); Elduque (1990) is cited for the
classification of Malcev modules \\
E5 & the theorems of the superfluid were first labelled \texttt{S1} to
\texttt{S4}, the labels of provenance documents of the main project &
the check of self-containedness, which searches for these labels as
places to read (stage 3b) & the theorems are labelled SF1 to SF4 \\
E6 & registry meanings used bare letters (a product written
\texttt{XY-YX}, a Casimir \texttt{s(s+1)}, a function \texttt{f}, a gap
\texttt{Delta}), and four items that the registry of
\texttt{Revision\_NoGo} also contains (\texttt{diag}, the transpose, the
tensor product, the expectation value) were not declared as shared &
checks LL-S6 and LL-S2 & meanings reworded; the shared items declared \\
E7 & manuscript drafts used a bare letter g in the introduction, letters
a, b, c as indices in the identity between Jacobiator and associators,
and a bare letter D for a derivation; the word ``left-handed'' for the
spinor charges depends on a convention & checks LL-S6 and LL-S7, review
& the sentence reworded; the six associators written out;
\(\mathcal D_{\rm LL}\); the chirality stated through the action of the
boosts and the block \(\gamma^{\rm ch}_{\rm LL}=+1\) \\
E8 & labels in four figures (tables, charges, manifolds,
renormalizability) overlapped & inspection of every figure & positions
and texts changed; every figure inspected again \\
E9 & plain-text names inside the registry meanings (\texttt{R\^{}8},
\texttt{M\_2(C)}, \ldots) were not escaped for LaTeX and stopped the
build of the manuscript & the build (stage 3c) & they are written as
code and converted to typewriter type for LaTeX \\
E10 & the check of the symbols rejected formulas of the manuscript
sources that were broken across two lines inside the suffix & checks
LL-S5 and LL-S6 & every run of white space in a formula counts as one
space, as for LaTeX \\
E11 & a page of the long table of the symbols ran 3.4 pt beyond its
height: TeX lets the last line of a page hang below it, and the closing
rule of the table was added after that line & the log check of the build
& the largest allowed depth is zero inside the long tables and restored
after them \\
E12 & two comments in the program were outdated: a theorem label
\texttt{S1} in the module of the superfluid, and the statement in
\texttt{Cargo.toml} that the library crate supplies the exact ranks (the
module \texttt{exact.rs} of the program computes them) & review of the
source before the description of the programs was written & both
comments corrected \\
E13 & the generic pair matrix \(\mathsf K\) of Theorem SF1, and the
coordinates \texttt{x5}, \texttt{x6}, \texttt{x7} in the meaning of the
taken symbol \(D\), were not explained in the registry & review of the
table of symbols & both are now explained in their entries \\
E14 & the first version of the program \texttt{explain\_superfluid}
typed in the action of the isospin generators on \(\vec d_{\rm LL}\) as
a rotation with the opposite sense to the one that follows from the pair
matrices; with the factor \(2i\) of \(F_{\rm LL}\) this would have named
the number of up particles, not of down particles, as the unbroken
generator of the up-up condensate (the report then said only ``one
combination of \(T^3_{\rm LL}\) and \(F_{\rm LL}\)'') & deriving the
statement of the explanation by hand (stage 4) & the action of every
generator on \(\vec d_{\rm LL}\) is now derived from its action on the
pair matrices (check X-7), and the unbroken generator of each charged
condensate is named and checked \\
E15 & in the first drafts of the program of the explanation the report
used \texttt{q} (the momentum) for the charge of a pair, an unregistered
name for a unit vector, and a condition of check X-9 that did not test
what its statement says & checks LL-S8 and review (stage 4) &
\texttt{Q\_LL(Kch)}, words, and the condition rewritten (the species
Casimir is 0 for \(\mathsf E_{\rm LL}\) and 2 for
\(\sigma^{\rm LL}_0\)) \\
E16 & the explanation placed computed numbers directly after a formula
instead of inside it; the Markdown converter then did not recognize the
formula, and three Greek letters were set in a text font without these
glyphs; the explanation also used bare Latin letters for the isospin
quantum number and for two amplitudes & the layout check of the build
and check LL-S6 (stage 4) & the numbers stand inside their formulas; the
letters were replaced by words \\
E17 & three cells of the table of the example condensates were wider
than their columns & the layout check of the build (stage 4) & the
column widths were changed \\
E18 & the first draft of check Z-6 (the mass dimensions) contained
unfinished code & review of the program before its first run (stage 5) &
rewritten: the dimensions follow, in doubled integers, from the
condition that every term of the Lagrangian density has the dimension
\(D+1\), and are compared with the closed forms \\
E19 & the first version of check Z-8 required that each doubling of the
cutoff divide the change of the renormalized bubble by about 4, the
power-counting bound; the computed factor was about 64 & the check
failed (stage 5) & the term of order \((p^{\rm E}_{\rm LL})^4\) of the
directional average was derived,
\(m_{\rm LL}^4/[(k^{\rm E}_{\rm LL})^2+m_{\rm LL}^2]^5\), and is now
verified numerically; the check requires the factor 64 and states the
bound 4 separately \\
E20 & figure labels with the abbreviated fractions
\texttt{\textbackslash{}frac12} and \texttt{\textbackslash{}frac32}
stopped the figure script, and several labels overlapped curves, axes or
the frame & the script and inspection of every figure (stage 5) &
fractions written in full; labels moved, given white backgrounds, or
placed in free corners; the four curves of Figure 3(b) shifted sideways
so that equal values stay visible \\
E21 & the report, the comments and the figures used letters without
their suffix or with a second meaning: \texttt{X} and \texttt{M} for
matrices, \texttt{x\^{}0} for coordinates, \texttt{e\^{}} for the
exponential, the parameter of the transformations without its suffix,
\texttt{P} for the sample point (the letter P is a symbol of the main
project, a generator), \texttt{R\_j}, \texttt{B\_j} and \texttt{G} for
the generator matrices (\texttt{G} is the contact coupling), \texttt{k}
for the loop momentum, and the angle variable of the quadrature named
\texttt{chi} like the spinor & review against the registry (stage 5) &
the registered names \texttt{M\_LL(x)}, \texttt{x\_LL\^{}p},
\texttt{exp}, \texttt{theta}, \texttt{V\_LL} and \texttt{kE}; the words
``sample point''; the angle variable renamed \\
E22 & the report of check Z-7 called the enumerated sets of numbers
``one-particle-irreducible graph classes''; such a class need not
contain an irreducible graph & review of the statement against the code
(stage 5) & the classes are defined in the report and in the document,
which explains why the list of divergent classes bounds the irreducible
divergent functions; the check now covers every largest number of
vertices from 2 to 6, as the data show \\
E23 & the comment at the head of the program said that every tolerance
other than \(10^{-12}\) is stated in the check, but checks Z-3, Z-4, Z-5
and Z-8 did not state theirs, and a comment in Z-8 described the factor
4 of the first version & review of the program (stage 5) & the
statements of the checks name their tolerances; the comment corrected \\
E24 & the first draft of the document quoted the conclusions with an
omission, called the six spin-zero pair matrices ``of the scalar
channel'' (three are scalar, three pseudoscalar), and had one equation
and one table column wider than the page & review and the layout check
of the build (stage 5) & the three sentences quoted in full; the pair
matrices described correctly; the equation split into two lines and the
column widened \\
E25 & the first version of check GS-3 changed the frame with
\(\mathbf g\) in place of
\(\tilde{\mathbf g}=(\mathbf g^\dagger)^{-1}\); the check failed & the
check (stage 6) & the change of representation (2.9) of the document:
the frame changes with \(\tilde{\mathbf g}\), the upper block of the
spinor with \(M_{\rm LL}(\tilde{\mathbf g})\), the lower block with
\(M_{\rm LL}(\mathbf g)\) \\
E26 & the first statement of check GS-13 gave the loss of modulus per
RK4 step as the loss of the squared modulus (twice too large) & the
check failed (stage 6) & the loss of the modulus, about a sixth power of
the energy times the time step divided by 144, with the squared modulus
stated separately \\
E27 & the statement of check GS-8 wrote the hole block with the
conjugate potential inside and a complex conjugation outside, which is
not the matrix that the program computes & review of the statements
(stage 6) & the hole block written as computed,
\texttt{-(h(-q\ -\ Acal)\ -\ Acal\_0\ -\ mu)\^{}*} \\
E28 & the printed hole equations wrote the conjugate order parameter
with a double superscript, which stopped the build of the document & the
build (stage 6) & written \((d^\iota_{\rm LL})^*\) \\
E29 & the report used \texttt{A\_LL} for two different potentials and
for the isospin matrix, \texttt{H\_BdG(q)} for the matrix with gauge
fields, \texttt{S(...)} for two different matrices, and \texttt{D\_LL}
for the bi-quaternion of covariant derivatives (\texttt{D\_LL} is the
Lorentz derivation) & review against the registry and check LL-S8 (stage
6) & \texttt{Av\_LL}, \texttt{AF\_LL}, \texttt{atau\_LL},
\texttt{H\_DB(q)}, \texttt{Sspin\_LL}, \texttt{sgen\_LL},
\texttt{Dbq\_LL}, all registered \\
E30 & the first component equations were written in the canonical gauge
only, without the terms of the spin connection & review against the
request of 06:45 (stage 6) & the expansions of the covariant derivatives
and the sixteen mean-field equations contain the spin connection, and
check GS-12 evaluates them with it \\
E31 & the connection was first written with a bold Greek letter, which
the mathematical font of the documents does not have in bold (it was
printed like the components) & inspection of a test build (stage 6) &
the connection is the bi-quaternion \(\mathbf u^\omega_{{\rm LL},p}\) \\
E32 & the first draft of the document used letters without suffix: the
angle and the length of the isovector in the link formula, a letter for
a \(2\times2\) matrix and for matrix entries, a letter for an energy, a
letter for a mode, the quasi-particle operators written like Dirac
matrices, the map from spinor charges to the field without suffix, and
the label of the spin zero of the test & checks LL-S6 and LL-S7 (stage
6) & registered symbols (\(\theta_{\rm LL}\),
\(M_{\rm LL}(\tilde{\mathbf g})\), \(\xi^{\rm LL}_p\),
\(E^{\rm qp}_{{\rm LL},n}\), \(\hat a^{\rm qp}_{{\rm LL},n}\),
\(\Psi_{\rm LL}[\mathbf b^-_{\rm LL},\mathbf b^+_{\rm LL}]\)), or words
(``spin-zero'') \\
E33 & ten formulas and one cell of the table of symbols were wider than
the page, and the sixteen mean-field equations in one display left a
page almost empty & the layout check of the build and inspection (stage
6) & formulas split into lines; the pairing block written as
\(\sigma^{\rm LL}_0\) times a \(4\times4\) matrix; the blocks of the
spin connection with a factor \(\tfrac12\); one display per equation \\
E34 & the new check LL-S9 counted the letters of the line spacing of the
printed equations (a layout command) as symbols & the check (stage 6) &
the check removes this layout command before it reads the letters \\
E35 & in the first drawing of Figure 1 of the document of Test 2 the
legend and two labels covered curves and points & inspection of the
figure (stage 6) & legend below the panel, labels moved into empty
regions \\
E36 & an editing script wrote parts of this provenance document with
escape sequences of the scripting language in place of backslashes of
the formulas & a scan of the sources for control characters (stage 6) &
the document was restored from the repository and edited again; every
source of the folder is scanned for control characters \\
E37 & the first time step of check AS-2 (0.02, 800 steps) would have
given an RK4 phase error of about \(8\times10^{-7}\), above the
tolerance \(10^{-8}\); the error was estimated before the first run &
the estimate of the summed phase error (stage 7) & the time step 0.004
with 4000 steps (expected error \(1.2\times10^{-9}\)), and the expected
size printed with every RK4 check \\
E38 & the first statements of the checks of item {[}2{]} used single
letters without the suffix LL, several of which the registry gives other
meanings: for the width of the wall and for the precession frequency,
for \(\varepsilon_{\rm LL}-\mu_{\rm LL}\), for the profile and the
spinor of a bound state, for \((\mathsf K^3_{\rm LL})^\dagger\), for the
third isospin component, and two Greek letters for the parameter of a
wall & review against rule 5 (stage 7) & new registry entries
(\(\ell^{\rm K}_{\rm LL}\), \(\omega^{\rm pr}_{\rm LL}\),
\(\varphi^{\rm K}_{\rm LL}\), \(\chi^{\rm K}_{\rm LL}\),
\(k^{\rm K}_{\rm LL}\), \ldots), \(Q_{\rm LL}\) for the charge of a
component, and \((\mathsf K^3_{\rm LL})^\dagger\) and
\(\Delta^\infty_{\rm LL}\ell^{\rm K}_{\rm LL}\) written out \\
E39 & the first tolerance of the density comparison of check AS-9 (0.999
for the Bhattacharyya coefficient) was a guess, and it failed (0.99888)
& the check (stage 7) & the tolerances of AS-9 are now the second-order
bounds of Kato-Temple and Davis-Kahan, computed from the residual and
the spectral distance \\
E40 & in the first figure of case B the arrow for the gap pointed at the
momentum 0 (a local maximum of \(E^-_{\rm LL}\)) instead of the Fermi
momentum, and in the first figure of case C a label covered the curve &
inspection of the figures (stage 7) & the arrow is placed at the
computed minimum of \(E^-_{\rm LL}\); the label moved \\
E41 & the first output of check AS-8 labelled the two sectors of
\(\mathsf Q^{\rm K}_{\rm LL}\) with the opposite sign & review of the
report (stage 7) & the label printed with the sign of the operator
\(\mp d\Delta_{\rm LL}/dx^1_{\rm LL}\) \\
E42 & the first draft of Section 3 of the document of Test 2 used
unregistered letters in formulas (a name for the \(2\times2\) matrix of
the particle-hole oscillation, bare letters for a profile, a residual
and a distance, bold units for the components of a spinor) and the first
draft of Section 11 of this provenance document cited equation numbers
of the document of Test 2 & the checks LL-S6, LL-S7 and review against
rule 7 (stage 7) & written out in words or with registered symbols;
Section 11 states every formula itself \\
E43 & the labels of the figures of item {[}2{]} first used notation that
the formula renderer of the plotting library does not accept (a
sans-serif letter without braces, a fraction without braces) & the
figure script stopped (stage 7) & braces added \\
E44 & the first couplings of the self-consistent examples (2.0 and 1.6)
gave gaps larger than the band width of the lattice (7.4 for the uniform
example); the couplings were then chosen by solving the lattice gap
equation (3.21) for gaps of 0.4 to 0.6 & & \\
E45 & the first self-consistent kink did not converge in 400 iterations:
the pair field turned out of its isospin direction; the analysis showed
that a wall of the pair field is unstable, and the option
\texttt{pair\_constraint\ =\ line} with the check of the full condition
(4.1) was added & & \\
E46 & the first check of the group velocity of the free packet failed
(deviation 1.6e-3, allowed first 2e-2, then 2e-3) because the packet
contained the other branch at its other momenta (Zitterbewegung); the
packet is now projected on its branch at every momentum, and the
deviation is 4e-7 & & \\
E47 & the first expected velocity of a hole packet had the wrong sign (a
hole of the branch \(+\varepsilon^{\rm lat}_{\rm LL}\) moves with
\(+d\varepsilon^{\rm lat}_{\rm LL}/dq^1_{\rm LL}\)) & & \\
E48 & the first cyclotron example put the centre of the orbit below the
packet, so that the orbit reached the wall, and compared the motion with
a classical orbit, although the packet consists of the lowest Landau
levels; the packet was moved and the comparison is made with the Landau
levels & & \\
E49 & the first quench was integrated with RK4, which would not conserve
the orthonormality and the fermion number (estimated change
\(10^{-4}\)); the first quench was also too large for the linear
prediction, which was first taken as twice the gap instead of the root
of the linearized equations of the finite ring (4.11) & & \\
E50 & the first check EX-3 required less than \(10^{-6}\) inside the
superfluid region, but the part of the packet above the gap passes the
interface (\(7\times10^{-6}\)); the first check EX-6 expected 16 bound
states, but one wall has eight (3.28) & & \\
E51 & the first form of the self-test report replaced the name of the
print function inside the function, which Python rejects & & \\
E52 & the first draft of Section 4 used letters without the suffix LL in
formulas (the coefficients of a two-level state, the number of steps, a
lattice index, the potential of the pulse, the Landau energies), and its
key table and one formula were wider than the page & & \\
E53 & the first quench output wrote the fermion number with twelve
significant digits only, so that its change (\(2\times10^{-11}\) of 96)
was invisible in the data file and in the first figure; the solver now
writes the change itself (\texttt{N\_F\_change}), and the figure legends
were moved where they cover no data & & \\
E54 & the first comparison of the two solvers failed for the task
\texttt{spectrum}: the densities of single eigenvectors inside a
degenerate eigenspace depend on the basis that each eigensolver chooses;
the solvers now write the densities summed over each group of degenerate
eigenvalues, which are independent of the basis & & \\
E55 & the first comparison also failed for the history of the Anderson
iteration of the uniform example (difference \(2\times10^{-7}\) at an
intermediate iterate, with the same fixed point to \(10^{-15}\)); the
histories are now compared with the tolerance \(10^{-5}\) and the
solutions with \(10^{-7}\), for the reason stated in check EX-8 & & \\
E56 & in the first drawings of the Rust plotter three legends covered
data (the crossing of the probabilities of the Andreev reflection, the
curves of the self-consistent kink and some eigenvalues), and in the
two-dimensional figure the axis labels of the second column of panels
overlapped the first; the legends and panels were moved & & \\
E57 & the first draft of Section 8 used letters without the suffix LL
for auxiliary quantities of its proofs (a small parameter, the entries
of a \(2\times2\) matrix, a rapidity, the operators of two states, an
index condition); the proofs were rewritten with registered symbols,
with particular Lorentz transformations, and with index conditions
stated in words & & \\
E58 & in the first draft of Section 9, auxiliary letters without the
suffix LL appeared in derivations (a generic stiffness coefficient, the
modulus of an integration variable, a unit vector and two generic
isovectors, a Kronecker symbol), the plain-text names of the two pair
energies clashed with an existing name, and three displays were wider
than the page; the derivations were rewritten with the registered
symbols (\(Z^{\rm t}_{\rm LL}\), \(-Z^{\rm s}_{\rm LL}\),
\(\vec A^{\rm v}_{{\rm LL},p}\), Fresnel integrals of the real and
imaginary parts), the names were changed to \texttt{eps(+)},
\texttt{eps(-)}, and the displays were split; the first version of check
SB-10 asked the energy of the time stepping to be conserved to
\(10^{-10}\), below the truncation error of the fourth-order method
(\(1.5\times10^{-8}\)), and was corrected to \(10^{-6}\) & & \\
E59 & the first draft of Section 10 used the big-O notation (a bare
letter) in formulas, primes on \(\Xi_{\rm LL}\) for its derivatives, a
symbol for the \(4\times4\) Dirac part of the plane-wave Hamiltonian,
and two abbreviations of numbers; they were replaced by words, by the
superscripts \((1)\), \((2)\), \((3)\), by the Dirac matrices written
out, and by the derivatives stated in words; in Figure 23 the first
drawing of the Mexican hat on a square grid had ragged edges and hid the
profile, and it was redrawn on a polar grid with a quarter cut away &
& \\
E60 & the first draft of Section 11 used letters without the suffix LL
in proofs and formulas (a parameter of a line of fields, the six real
coordinates of the field and a matrix of a generator in the proof of
Goldstone's theorem, the symbol of a variation, a charge number, the
variable of a figure), and the derivatives of the potential of one
sector were written with nested braces in the superscript, which the
check LL-S5 does not read as a suffixed symbol; they were replaced by
the registered \(\kappa^{\rm st}_{\rm LL}\), by the real and imaginary
parts of \(d^\iota_{\rm LL}\) and the first-order change
\(\vec d^{\rm tr}_{\rm LL}\), by words, and by the superscripts
\({\rm s(2)}\), \({\rm s(3)}\), \({\rm s(4)}\). In Figure 24(a) the
first drawing of the 40 minimization paths, projected from six
dimensions on the plane of the two gaps, was a tangle of lines, partly
outside the window; it was replaced by the lines of the exact gradient
flow on the invariant plane of the normal form (11.5), and the paths are
shown by their convergence in panel (c). The third and fourth numerical
derivatives of check SB-14 first had errors larger than their tolerance
and are now Richardson-extrapolated; the minimizer first ran to its
iteration limit because its stopping test asked for a gradient below the
noise of the numerical gradient, and it now stops at \(10^{-10}\) or
when a step no longer lowers the potential & & \\
E61 & the first draft of Section 12 used letters without the suffix LL
in derivations (three generic complex vectors in the expansion of a
squared length, two generic real vectors in Lagrange's identity, two
abbreviations in the expansion of the square roots of (12.7), the shift
of the momentum by the third-axis potential, the amplitude of a
deviation of the field, and a positive number in the closed form of the
kernel, written with the letter of the registered twist charge
\(\Theta_{\rm LL}\)), the frequency without its scripts, and a reference
to the decision D18 of this document, which the document of Test 2 may
not make (its reasons must stand in it); they were replaced by words, by
the registered symbols \(\delta^{\rm f}_{\rm LL}\),
\(\vec u^{\rm f}_{\rm LL}\), \(c^{\rm w}_{\rm LL}\),
\(\vec a^{\rm F}_{\rm LL}\), by the real and the imaginary part of
\(\vec d_{\rm LL}\), and by \(\omega^{\rm B}_{\rm LL}\). In the report
of check SB-16 the plain-text name \texttt{eps} was used for the
relative displacement of the starting gap, which clashed with
\texttt{eps} for \(\varepsilon_{\rm LL}\); it was renamed \texttt{disp},
the contact matrix \texttt{N} was renamed \texttt{N\_g}, and the
momentum and the potential of the eigenvalue formula, first written
\texttt{p} and \texttt{a}, are now \texttt{q} and \texttt{a\_F}. The
labels of Figure 25(c) first used the frequency without its scripts. One
display of the proof of Theorem 12.2 was wider than the page and was
split. The check of the energy in constant potentials first failed its
tolerance (relative deviations \(2.6\times10^{-6}\) and
\(1.8\times10^{-5}\) from the terms of the fourth order at the potential
0.02); the second differences are now taken at 0.005 and 0.01 and
extrapolated. The \(16\times16\) unit matrix, also used by the folder
Revision\_NoGo, was declared shared & & \\
E62 & in the first draft of Section 13 the mass matrices were written
with a styled letter that was not yet registered (check LL-S7); it was
registered as \(\mathsf M^{\rm t}_{\rm LL}\),
\(\mathsf M^{\rm s}_{\rm LL}\) together with the other new symbols of
item {[}11{]}. After the equations of Section 13 were renumbered, the
report of check SB-17 and the legend of Figure 26(c) still named the
dispersion relation of the longitudinal mode by its old number (13.7);
they were corrected to (13.9). The first version of the first-order test
of check SB-17 used potentials of the size 0.3, for which the rotation
of the potentials by the Goldstone angles, a product of an angle and a
potential, is not small, and the test failed its tolerance; it now uses
fields and potentials of the size \(10^{-4}\), with the tolerance
\(10^{-3}\) relative to their size. In Figure 26(c) the first legend
covered the curves; it was moved and framed, and the frequency axis was
limited to 1.08. In the report of check SB-17 the eigenvalues were first
written with \texttt{Delta\^{}2}, which names \(\Delta_{\rm LL}\), for
the square of the gap of the vacuum; it is now \texttt{Delta\_vac\^{}2}.
The first labels of Figure 26 wrote the potentials without the index
\(p\) and named the photon by the potential \(A^3_{{\rm LL},p}\) instead
of the field \(V^3_{{\rm LL},p}\) of the unitary gauge; they were
corrected & & \\
E63 & from the first version on, every time integration of Test 2 (the
task \texttt{evolve} of both solvers, the checks AS-2, AS-4, SB-10 and
SB-16 by RK4 with a fixed step; the task \texttt{quench} by the implicit
midpoint rule with a fixed step) was done with a fixed-step method
written by hand, although the author's repository provides CVODE of
SUNDIALS in pure Rust. The author stopped the work on 2026-11-18
(Section 2.9). The consequences, measured by the program
\texttt{t2\_cvcmp}: the quench was wrong in the third decimal (mean of
the gap 0.500877 instead of 0.503125, frequency 1.18449 instead of
1.18898, amplitude 0.00587 instead of 0.00708; the conserved functional
\(E^{\rm MF}_{\rm LL}\) changed by \(2\times10^{-3}\) of its value), the
norm of the evolve examples was conserved only to \(3\times10^{-7}\),
and the stability rule \(\ell^{\rm t}_{\rm LL}B^{\rm H}_{\rm LL}\le2.5\)
of the first version allows steps that give the relative error 0.37 on a
random lattice problem. All these calculations were redone with CVODE
(stage 18), and rule 9 of Section 3 now forbids fixed-step methods
written by hand & & \\
E64 & during the change to CVODE: the first test of the module
\texttt{t2core::cvode} (one oscillator over 100 periods) used the
standard error weights of CVODE for every real unknown and needed 36591
steps for the deviation \(1.8\times10^{-8}\); with the error weights of
complex components (decision D22) it needs 13643 steps for
\(3.3\times10^{-9}\); the first version of check GS-13 asked that the
error fall by a factor of 10 from one tolerance to the next, which the
stability-limited steps of the Adams-Moulton formulas do not give
(1.2e-6 to 1.7e-7), and it now compares each error with the error budget
\(B^{\rm err}_{\rm LL}\); the first conservation check of the solver
compared the change of the norm with the number of steps times
\(r^{\rm tol}_{\rm LL}\), which is not a bound (2.56e-9 against 2.3e-9
in the example \texttt{cyclotron}), and it now uses
\(10B^{\rm err}_{\rm LL}\) with the factor \(\sqrt2\) and the term of
\(a^{\rm tol}_{\rm LL}\); the reports of the solvers first wrote the
mode as \texttt{w} and the mode matrix as \texttt{W}, which are the
names of the bi-octonion \(\mathbf w\) and of nothing registered, and
now write \texttt{w\_lat} and \texttt{W\_modes}; the report of check
GS-13 first named the complex components \texttt{w\_k} and the
eigenpairs \texttt{E\_n}, \texttt{v\_n}, and now uses \texttt{w\_cmp},
\texttt{E\_qp} and words & & \\
\end{longtable}

\hypertarget{the-symbols}{%
\section{12. The symbols}\label{the-symbols}}

\textbf{1 Indices, conventions and number sets}

\begin{longtable}[]{@{}
  >{\raggedright\arraybackslash}p{(\columnwidth - 6\tabcolsep) * \real{0.2200}}
  >{\raggedright\arraybackslash}p{(\columnwidth - 6\tabcolsep) * \real{0.5000}}
  >{\raggedright\arraybackslash}p{(\columnwidth - 6\tabcolsep) * \real{0.0800}}
  >{\raggedright\arraybackslash}p{(\columnwidth - 6\tabcolsep) * \real{0.2000}}@{}}
\toprule\noalign{}
\begin{minipage}[b]{\linewidth}\raggedright
symbol
\end{minipage} & \begin{minipage}[b]{\linewidth}\raggedright
meaning
\end{minipage} & \begin{minipage}[b]{\linewidth}\raggedright
origin
\end{minipage} & \begin{minipage}[b]{\linewidth}\raggedright
plain-text names
\end{minipage} \\
\midrule\noalign{}
\endhead
\bottomrule\noalign{}
\endlastfoot
\(i,j\) & index of space-like coordinates; values 1, 2, 3 (ordinary
space) unless a sum is said to run over all space-like coordinates of
the spacetime. In Revision\_SuperLoop only \(j\) is used as an index,
with the values 1, 2, 3, and with primes (\(j'\), \(j''\)) where two or
three such indices are needed: it labels the three imaginary quaternion
units \(\mathbf e_j\), which correspond to the three space directions of
\(\mathbb R^{1,3}_{\rm LL}\), the coordinates \(x^j_{\rm LL}\), the
momentum components \(q^j_{\rm LL}\), the rotation and boost generators
\(J^j_{\rm LL}\), \(K^j_{\rm LL}\), and the Pauli matrices
\(\sigma^{\rm LL}_j\) and Dirac matrices \(\alpha^j_{\rm LL}\),
\(\gamma^j_{\rm LL}\). The letter \(i\) is the imaginary unit (entry
\(i\)); the isospin index is \(\iota\) (a different letter). & taken &
\texttt{j} \texttt{j\textquotesingle{}}
\texttt{j\textquotesingle{}\textquotesingle{}} \\
\(D\) & the number of space dimensions in which the particles of matter
have momentum: 7 (all space coordinates of 7+1), 4 (fields that do not
depend on \(x5,x6,x7\)) or 3 (fields that depend on ordinary space
only). \(V_D=\pi^{D/2}/\Gamma_{\rm E}(D/2+1)\) is the volume and
\(A_D=D\,V_D\) the area of the unit ball in \(D\) dimensions
(\(\Gamma_{\rm E}\): Euler's gamma function). In Revision\_SuperLoop
\(D=3\) for the four-dimensional world of the superfluid and \(D=7\) for
the eight-dimensional manifolds (on \(4+4\) power counting sees the
total dimension 8 as well); the value 4, the coordinates \(x5,x6,x7\) of
the eight-dimensional spacetime of the project (on which such fields do
not depend), and the volume and the area of the unit ball do not occur.
Section 12.5 of the document of Test 2 continues \(D\) away from 3, in
dimensional regularization (the momentum integrals taken in \(D\) space
dimensions). & taken & \texttt{D} \\
\(i\) & the imaginary unit, \(i^2=-1\); in \(\mathbb B_{\rm LL}\) it is
the imaginary unit of the complex coefficients: it commutes and
associates with every bi-octonion and is not one of the units
\(\mathbf e_1,\dots,\mathbf e_7\) & new & \\
\(\mathbb R\), \(\mathbb C\) & the real and the complex numbers (plain
text: \texttt{R\^{}8}, \texttt{C\^{}3}, \texttt{C\^{}8} are the real or
complex vector spaces of that dimension) & new & \texttt{R} \texttt{C}
\texttt{R\^{}8} \texttt{C\^{}3} \texttt{C\^{}8} \\
\(p\), \(p'\) & a summation index over the basis of an algebra (values
\(0,\dots,7\), or \(0,\dots,3\) for the quaternions): in
Revision\_SuperLoop it labels the units \(\mathbf e_p\) of
\(\mathbb B_{\rm LL}\), the coefficients \(\xi^{\rm LL}_p(\mathbf x)\),
the coordinates \(x^p_{\rm LL}\) of \(\mathbb R^{1,7}_{\rm LL}\) and
\(\mathbb R^{1,3}_{\rm LL}\) and the Dirac matrices
\(\gamma^p_{\rm LL}\) (\(p=0,1,2,3\)); \(p'\) is a second index of the
same kind (plain text \texttt{p1}, \texttt{p2}, \texttt{p3}: three such
indices) & new & \texttt{p} \texttt{p\textquotesingle{}} \texttt{p1}
\texttt{p2} \texttt{p3} \\
\(\iota\), \(\iota'\), \(\iota''\) & the isospin index, values 1, 2, 3:
it labels the isospin generators \(T^\iota_{\rm LL}\), the Pauli
matrices \(\tau^{\rm LL}_\iota\) on the isospin index of a field, the
pair channels \(\mathsf K^\iota_{\rm LL}\),
\(\mathsf K'^\iota_{\rm LL}\), the pair operators \(\Pi^\iota_{\rm LL}\)
and the components \(d^\iota_{\rm LL}\) of the order parameter; it is an
internal direction, not a direction in space (that is \(j\)) & new &
\texttt{iota} \\
\(\epsilon^{\rm LL}_{jj'j''}\) & the Levi-Civita symbol: \(+1\) if
\((j,j',j'')\) is an even permutation of \((1,2,3)\), \(-1\) if it is an
odd one, \(0\) otherwise; also with the isospin indices,
\(\epsilon^{\rm LL}_{\iota\iota'\iota''}\) & new & \texttt{eps3} \\
\(\sigma^{\rm LL}_0\), \(\sigma^{\rm LL}_j\) & the \(2\times2\) unit
matrix \(\sigma^{\rm LL}_0\) and the Pauli matrices
\(\sigma^{\rm LL}_1=\begin{pmatrix}0&1\\1&0\end{pmatrix}\),
\(\sigma^{\rm LL}_2=\begin{pmatrix}0&-i\\i&0\end{pmatrix}\),
\(\sigma^{\rm LL}_3=\begin{pmatrix}1&0\\0&-1\end{pmatrix}\), acting on
the two components of a Weyl spinor or on the species index & new &
\texttt{sigma\_0} \texttt{sigma\_j} \texttt{sigma\_1} \texttt{sigma\_2}
\texttt{sigma\_3} \\
\(\tau^{\rm LL}_\iota\) & the same three Pauli matrices acting on the
isospin index of a field (the index of the two components up and down of
an isodoublet) & new & \texttt{tau\_iota} \texttt{tau\_1}
\texttt{tau\_2} \texttt{tau\_3} \\
\(\mathsf E_{\rm LL}\) & the real antisymmetric \(2\times2\) matrix
\(i\sigma^{\rm LL}_2=\begin{pmatrix}0&1\\-1&0\end{pmatrix}\), with
\(\mathsf E_{\rm LL}^{-1}=\mathsf E_{\rm LL}^T=-\mathsf E_{\rm LL}\); it
acts on a Weyl, isospin or species index, as stated at each place & new
& \texttt{eps2} \\
\(\mathsf 1_2\), \(\mathsf 1_4\), \(\mathsf 1_8\), \(\mathsf 1_{16}\) &
the \(2\times2\), the \(4\times4\), the \(8\times8\) and the
\(16\times16\) unit matrices & new & \\
\end{longtable}

\textbf{2 The bi-octonions and the super-loop}

\begin{longtable}[]{@{}
  >{\raggedright\arraybackslash}p{(\columnwidth - 6\tabcolsep) * \real{0.2200}}
  >{\raggedright\arraybackslash}p{(\columnwidth - 6\tabcolsep) * \real{0.5000}}
  >{\raggedright\arraybackslash}p{(\columnwidth - 6\tabcolsep) * \real{0.0800}}
  >{\raggedright\arraybackslash}p{(\columnwidth - 6\tabcolsep) * \real{0.2000}}@{}}
\toprule\noalign{}
\begin{minipage}[b]{\linewidth}\raggedright
symbol
\end{minipage} & \begin{minipage}[b]{\linewidth}\raggedright
meaning
\end{minipage} & \begin{minipage}[b]{\linewidth}\raggedright
origin
\end{minipage} & \begin{minipage}[b]{\linewidth}\raggedright
plain-text names
\end{minipage} \\
\midrule\noalign{}
\endhead
\bottomrule\noalign{}
\endlastfoot
\(\mathbb H\), \(\mathbb O\) & the quaternions and the octonions: the
real combinations of \(\mathbf e_0,\dots,\mathbf e_3\) and of
\(\mathbf e_0,\dots,\mathbf e_7\); \({\rm Im}\,\mathbb H\) and
\({\rm Im}\,\mathbb O\) are their imaginary parts, the real combinations
of \(\mathbf e_1,\mathbf e_2,\mathbf e_3\) and of
\(\mathbf e_1,\dots,\mathbf e_7\) (plain text \texttt{Im\ H},
\texttt{Im\ O}; \texttt{C\ (x)\ H} is \(\mathbb C\otimes\mathbb H\), the
bi-quaternions) & new & \texttt{H} \texttt{O} \\
\(\mathbb B_{\rm LL}\) & the bi-octonions \(\mathbb C\otimes\mathbb O\):
the complex combinations
\(\mathbf x=\sum_{p=0}^7\xi^{\rm LL}_p(\mathbf x)\,\mathbf e_p\) of the
eight units, with the product of the units extended bilinearly; real
dimension 16 & new & \texttt{B\_LL} \\
\(\mathbf e_0,\dots,\mathbf e_7\); \(\mathbf e_p\), \(\mathbf e_j\),
\(\mathbf e_\iota\) & the units of \(\mathbb B_{\rm LL}\), of
\(\mathbb O\) and of \(\mathbb H\): \(\mathbf e_0=1\);
\(\mathbf e_1,\mathbf e_2,\mathbf e_3\) the quaternion units,
\(\mathbf e_j^2=-1\),
\(\mathbf e_1\mathbf e_2=-\mathbf e_2\mathbf e_1=\mathbf e_3\) and
cyclically; \(\mathbf e_4\) the unit of the Cayley-Dickson doubling,
\(\mathbf e_4^2=-1\); \(\mathbf e_5=\mathbf e_1\mathbf e_4\),
\(\mathbf e_6=\mathbf e_2\mathbf e_4\),
\(\mathbf e_7=\mathbf e_3\mathbf e_4\) (\(\mathbf e_\iota\): a
quaternion unit labelled by the isospin index; plain text \texttt{ie1},
\ldots, \texttt{ie7}: the bi-octonions
\(i\mathbf e_1,\dots,i\mathbf e_7\)) & new & \texttt{e0} \texttt{e1}
\texttt{e2} \texttt{e3} \texttt{e4} \texttt{e5} \texttt{e6} \texttt{e7}
\texttt{e\_p} \texttt{e\_p\textquotesingle{}} \texttt{e\_j}
\texttt{e\_iota} \texttt{e\_p1} \texttt{e\_p2} \texttt{e\_p3}
\texttt{ie1} \texttt{ie2} \texttt{ie3} \texttt{ie4} \texttt{ie5}
\texttt{ie6} \texttt{ie7} \\
\(\xi^{\rm LL}_p(\mathbf x)\) & the complex coefficient of the unit
\(\mathbf e_p\) in the bi-octonion \(\mathbf x\) & new & \texttt{xi\_p}
\texttt{xi\_0} \texttt{xi\_1} \texttt{xi\_2} \texttt{xi\_3} \\
\(\mathbf x\), \(\mathbf y\), \(\mathbf z\), \(\mathbf w\) & elements of
\(\mathbb B_{\rm LL}\) (bi-octonions); \(\mathbf w\) is also a point of
the super-loop manifold \(\mathbb L_{\rm LL}\) & new & \texttt{x}
\texttt{y} \texttt{z} \texttt{w} \\
\(\mathbf a\), \(\mathbf b\), \(\mathbf c\), \(\mathbf d\);
\(\mathbf b'\) & bi-quaternions: elements of
\(\mathbb C\otimes\mathbb H\), the complex combinations of
\(\mathbf e_0,\mathbf e_1,\mathbf e_2,\mathbf e_3\); every bi-octonion
is written \(\mathbf x=\mathbf a+\mathbf b\,\mathbf e_4\)
(Cayley-Dickson form), and \(\mathbf b\,\mathbf e_4\),
\(\mathbf b'\,\mathbf e_4\) are spinor charges & new & \texttt{aH}
\texttt{bS} \texttt{cH} \texttt{dS} \texttt{bS\textquotesingle{}} \\
\(\bar{\mathbf x}\), \(\mathbf x^*\), \(\mathbf x^\dagger\) & the three
conjugations (also of bi-quaternions): \(\bar{\mathbf x}\) changes the
signs of the coefficients of \(\mathbf e_1,\dots,\mathbf e_7\) (it is
complex-linear,
\(\overline{\mathbf x\mathbf y}=\bar{\mathbf y}\,\bar{\mathbf x}\), and
\(\overline{\mathbf a+\mathbf b\mathbf e_4}=\bar{\mathbf a}-\mathbf b\mathbf e_4\));
\(\mathbf x^*\) replaces every complex coefficient by its complex
conjugate; \(\mathbf x^\dagger=(\bar{\mathbf x})^*\) (plain text:
\texttt{xbar}, \texttt{x\^{}*}, \texttt{x\^{}dag}, and likewise
\texttt{aHbar}, \texttt{gL\^{}dag}, \ldots) & new & \texttt{xbar}
\texttt{x\^{}*} \texttt{x\^{}dag} \texttt{aHbar} \texttt{bSbar}
\texttt{bS\textquotesingle{}bar} \texttt{cHbar} \texttt{dSbar}
\texttt{gLbar} \texttt{hLbar} \texttt{rPbar}
\texttt{rP\textquotesingle{}bar} \texttt{bS\^{}*}
\texttt{bS\textquotesingle{}\^{}*} \texttt{gL\^{}*} \texttt{gL\^{}dag}
\texttt{uL\^{}dag} \\
\(N_{\rm LL}(\mathbf x)\), \(N_{\rm LL}(\mathbf x,\mathbf y)\) & the
complex norm
\(N_{\rm LL}(\mathbf x)=\sum_p\xi^{\rm LL}_p(\mathbf x)^2\), with
\(\mathbf x\bar{\mathbf x}=\bar{\mathbf x}\mathbf x=N_{\rm LL}(\mathbf x)\,\mathbf e_0\),
and its polarization
\(N_{\rm LL}(\mathbf x,\mathbf y)=\sum_p\xi^{\rm LL}_p(\mathbf x)\,\xi^{\rm LL}_p(\mathbf y)\),
\(N_{\rm LL}(\mathbf x,\mathbf x)=N_{\rm LL}(\mathbf x)\)
(complex-valued; for bi-quaternions the determinant of the \(2\times2\)
matrix \(M_{\rm LL}\)) & new & \texttt{N\_LL} \\
\([\mathbf x,\mathbf y]\); \([\mathbf x,\mathbf y,\mathbf z]\) & the
commutator \(\mathbf x\mathbf y-\mathbf y\mathbf x\) and the associator
\((\mathbf x\mathbf y)\mathbf z-\mathbf x(\mathbf y\mathbf z)\) of
bi-octonions (the same brackets of other elements are defined at their
place: the bracket of an algebra \(\mathcal F^{\rm A}_{\rm LL}\),
\(\mathcal F^{\rm B}_{\rm LL}\), \(\mathcal F^{\rm C}_{\rm LL}\), and
the commutator of two matrices or operators, the product in one order
minus the product in the other) & new & \texttt{{[}x,y{]}}
\texttt{{[}x,y,z{]}} \\
\(\mathrm{Jac}_{\rm LL}\) & the Jacobiator
\(\mathrm{Jac}_{\rm LL}(\mathbf x,\mathbf y,\mathbf z)=[[\mathbf x,\mathbf y],\mathbf z]+[[\mathbf y,\mathbf z],\mathbf x]+[[\mathbf z,\mathbf x],\mathbf y]\)
of a bracket; it vanishes identically exactly in a Lie algebra & new &
\texttt{Jac} \\
\(L^{\rm LL}_{\mathbf x}\), \(R^{\rm LL}_{\mathbf x}\) & the left and
right multiplications \(\mathbf y\mapsto\mathbf x\mathbf y\) and
\(\mathbf y\mapsto\mathbf y\mathbf x\) by a bi-octonion, real
\(16\times16\) matrices (on \(\mathbb O\) or
\(\mathbb O^{\rm s}_{\rm LL}\): \(8\times8\)) & new & \texttt{L\_x}
\texttt{L\_y} \texttt{L\_e1} \texttt{L\_e4} \texttt{L\_{[}e1,e4{]}}
\texttt{L\_(e\_p)} \texttt{L\_(e\_p\textquotesingle{})} \texttt{R\_x}
\texttt{R\_y} \\
\(\mathbb L_{\rm LL}\) & the Lorentz super-loop: the unit bi-octonions,
\(\mathbb L_{\rm LL}=\{\mathbf w\in\mathbb B_{\rm LL}:N_{\rm LL}(\mathbf w)=1\}\),
with the product of \(\mathbb B_{\rm LL}\); a Moufang loop (not a
group), and a manifold of real dimension 14 & new & \texttt{L\_LL} \\
\(\mathbf g\), \(\mathbf h\) & elements of \(SL(2,\mathbb C)\),
i.e.~unit bi-quaternions (\(N_{\rm LL}(\mathbf g)=1\)), a subgroup of
\(\mathbb L_{\rm LL}\) & new & \texttt{gL} \texttt{hL} \\
\(\mathbf v\) & a unit (bi-)quaternion that rotates the isospin of the
spinor charges (from the left) & new & \texttt{vI} \\
\(\mathfrak m_{\rm LL}\), \(\mathfrak l_{\rm LL}\),
\(\mathfrak s_{\rm LL}\) & the Malcev algebra of the super-loop,
\(\mathfrak m_{\rm LL}=\mathbb C\otimes{\rm Im}\,\mathbb O\) with the
commutator (real dimension 14), its Lorentz part
\(\mathfrak l_{\rm LL}=\mathbb C\otimes{\rm Im}\,\mathbb H\) (6, a Lie
algebra isomorphic to \(\mathfrak{so}(3,1)\)), and the space of spinor
charges
\(\mathfrak s_{\rm LL}=(\mathbb C\otimes\mathbb H)\,\mathbf e_4\) (8);
\(\mathfrak m_{\rm LL}=\mathfrak l_{\rm LL}\oplus\mathfrak s_{\rm LL}\)
& new & \texttt{m\_LL} \texttt{l\_LL} \texttt{s\_LL} \\
\(J^j_{\rm LL}\), \(K^j_{\rm LL}\) & the rotation and boost generators
of \(\mathfrak l_{\rm LL}\), \(J^j_{\rm LL}=\mathbf e_j/2\) and
\(K^j_{\rm LL}=i\,\mathbf e_j/2\) & new & \texttt{J\_LL\^{}j}
\texttt{K\_LL\^{}j} \texttt{J\_LL\^{}j\textquotesingle{}}
\texttt{K\_LL\^{}j\textquotesingle{}}
\texttt{J\_LL\^{}j\textquotesingle{}\textquotesingle{}}
\texttt{K\_LL\^{}j\textquotesingle{}\textquotesingle{}}
\texttt{J\_LL\^{}3} \\
\(\mathbf u\), \(\mathbf u'\), \(\mathbf u_1\), \(\mathbf u_2\) &
elements of the Lorentz part \(\mathfrak l_{\rm LL}\) (Lorentz
generators): imaginary bi-quaternions & new & \texttt{uL}
\texttt{uL\textquotesingle{}} \texttt{uL1} \texttt{uL2} \\
\(\mathbf t\), \(\mathbf t'\) & elements of \({\rm Im}\,\mathbb H\)
(real imaginary quaternions): isospin generators, which act on the
spinor charges from the left & new & \texttt{tI}
\texttt{tI\textquotesingle{}} \\
\(W_{\rm LL}(\mathbf u)\) & the Weyl representation of
\(\mathfrak l_{\rm LL}\) on bi-quaternions,
\(W_{\rm LL}(\mathbf u)\mathbf b=-\mathbf b\,\mathbf u\): a Lie
representation,
\([W_{\rm LL}(\mathbf u_1),W_{\rm LL}(\mathbf u_2)]=W_{\rm LL}([\mathbf u_1,\mathbf u_2])\)
& new & \texttt{W\_LL} \\
\(\rho_{\rm LL}(\mathbf u)\) & the inner action of a Lorentz element on
the spinor charges,
\([\mathbf u,\mathbf b\,\mathbf e_4]=(\rho_{\rm LL}(\mathbf u)\mathbf b)\,\mathbf e_4\)
with \(\rho_{\rm LL}=-2W_{\rm LL}\); not a representation & new &
\texttt{rho\_LL} \\
\(\Theta_{\rm LL}\) & the twist charge: the linear map of
\(\mathfrak m_{\rm LL}\) that is \(1\) on \(\mathfrak l_{\rm LL}\) and
\(-2\) on \(\mathfrak s_{\rm LL}\); the inner action is
\([\mathbf u,\cdot\,]=D_{\rm LL}(\mathbf u)\circ\Theta_{\rm LL}\) (as a
number: the eigenvalue \(1\) or \(-2\), or the value of
\(c^{\rm tw}_{\rm LL}\) in the families) & new & \texttt{Theta\_LL}
\texttt{Theta\_LL\^{}2} \\
\(D_{\rm LL}(\mathbf u)\), \(D^{\rm v}_{\rm LL}(\mathbf t)\) & the
Lorentz and isospin derivations of \(\mathbb B_{\rm LL}\):
\(D_{\rm LL}(\mathbf u)(\mathbf a+\mathbf b\mathbf e_4)=[\mathbf u,\mathbf a]+(-\mathbf b\mathbf u)\mathbf e_4\)
and
\(D^{\rm v}_{\rm LL}(\mathbf t)(\mathbf a+\mathbf b\mathbf e_4)=(\mathbf t\mathbf b)\mathbf e_4\)
(both linear in their argument) & new & \texttt{D\_LL}
\texttt{Dv\_LL} \\
\(\mathcal D_{\rm LL}\) & a derivation of \(\mathbb B_{\rm LL}\): a
linear map with
\(\mathcal D_{\rm LL}(\mathbf x\mathbf y)=(\mathcal D_{\rm LL}\mathbf x)\mathbf y+\mathbf x(\mathcal D_{\rm LL}\mathbf y)\),
the infinitesimal form of an automorphism & new & \\
\(\varphi^{\rm LL}_{\mathbf g,\mathbf v}\) & the automorphism
\(\mathbf a+\mathbf b\mathbf e_4\mapsto\mathbf g\mathbf a\bar{\mathbf g}+(\mathbf v\mathbf b\bar{\mathbf g})\mathbf e_4\)
of \(\mathbb B_{\rm LL}\) (\(\mathbf g\) a Lorentz transformation,
\(\mathbf v\) an isospin rotation) & new & \texttt{varphi} \\
\(\varrho_{\rm LL}\); \(\mathsf O_1\), \(\mathsf O_2\), \(\mathsf O_3\)
& a linear map from \(\mathfrak m_{\rm LL}\) to the linear operators
(matrices) on a vector space that turns brackets into commutators,
\(\varrho_{\rm LL}([\mathbf x,\mathbf y])=\varrho_{\rm LL}(\mathbf x)\varrho_{\rm LL}(\mathbf y)-\varrho_{\rm LL}(\mathbf y)\varrho_{\rm LL}(\mathbf x)\)
(Theorem R); \(\mathsf O_1,\mathsf O_2,\mathsf O_3\) three such
operators & new & \\
\(\mathrm{Kil}_{\rm LL}\); \(\mathrm{ad}\) & the Killing form
\(\mathrm{Kil}_{\rm LL}(\mathbf x,\mathbf y)=\mathrm{tr}(\mathrm{ad}\,\mathbf x\,\mathrm{ad}\,\mathbf y)\)
of a Lie algebra, where \(\mathrm{ad}\,\mathbf x\) is the matrix of the
map \(\mathbf y\mapsto[\mathbf x,\mathbf y]\) & new & \texttt{Killing}
\texttt{ad} \\
\end{longtable}

\textbf{3 Operators on the super-loop and their eigenvalues}

\begin{longtable}[]{@{}
  >{\raggedright\arraybackslash}p{(\columnwidth - 6\tabcolsep) * \real{0.2200}}
  >{\raggedright\arraybackslash}p{(\columnwidth - 6\tabcolsep) * \real{0.5000}}
  >{\raggedright\arraybackslash}p{(\columnwidth - 6\tabcolsep) * \real{0.0800}}
  >{\raggedright\arraybackslash}p{(\columnwidth - 6\tabcolsep) * \real{0.2000}}@{}}
\toprule\noalign{}
\begin{minipage}[b]{\linewidth}\raggedright
symbol
\end{minipage} & \begin{minipage}[b]{\linewidth}\raggedright
meaning
\end{minipage} & \begin{minipage}[b]{\linewidth}\raggedright
origin
\end{minipage} & \begin{minipage}[b]{\linewidth}\raggedright
plain-text names
\end{minipage} \\
\midrule\noalign{}
\endhead
\bottomrule\noalign{}
\endlastfoot
\(Q_{\rm LL}\) & the electric charge, \(Q_{\rm LL}=T^3_{\rm LL}\):
eigenvalues \(0\) on \(\mathbb C\otimes\mathbb H\) and \(\pm\tfrac12\)
on the spinor charges & new & \texttt{Q\_LL} \\
\(T^\iota_{\rm LL}\), \((T_{\rm LL})^2\) & the isospin generators
\(T^\iota_{\rm LL}=i\,D^{\rm v}_{\rm LL}(\mathbf e_\iota/2)\),
Hermitian, \([T^1_{\rm LL},T^2_{\rm LL}]=iT^3_{\rm LL}\) and cyclically,
and the isospin Casimir
\((T_{\rm LL})^2=\sum_\iota(T^\iota_{\rm LL})^2\) & new &
\texttt{T\_LL\^{}iota} \texttt{T\_LL\^{}iota\textquotesingle{}}
\texttt{T\_LL\^{}iota\textquotesingle{}\textquotesingle{}}
\texttt{T\_LL\^{}3} \texttt{T\_LL\^{}2} \\
\(\mathrm{Cas}_{\rm LL}\) & the Casimir operator of the rotations,
\(\mathrm{Cas}_{\rm LL}=-\sum_jD_{\rm LL}(J^j_{\rm LL})^2\); its
eigenvalue is 0 for spin 0, 3/4 for spin 1/2 and 2 for spin 1 & new &
\texttt{Cas\_LL} \\
\end{longtable}

\textbf{4 The covariant families and their classification}

\begin{longtable}[]{@{}
  >{\raggedright\arraybackslash}p{(\columnwidth - 6\tabcolsep) * \real{0.2200}}
  >{\raggedright\arraybackslash}p{(\columnwidth - 6\tabcolsep) * \real{0.5000}}
  >{\raggedright\arraybackslash}p{(\columnwidth - 6\tabcolsep) * \real{0.0800}}
  >{\raggedright\arraybackslash}p{(\columnwidth - 6\tabcolsep) * \real{0.2000}}@{}}
\toprule\noalign{}
\begin{minipage}[b]{\linewidth}\raggedright
symbol
\end{minipage} & \begin{minipage}[b]{\linewidth}\raggedright
meaning
\end{minipage} & \begin{minipage}[b]{\linewidth}\raggedright
origin
\end{minipage} & \begin{minipage}[b]{\linewidth}\raggedright
plain-text names
\end{minipage} \\
\midrule\noalign{}
\endhead
\bottomrule\noalign{}
\endlastfoot
\(\mathcal F^{\rm A}_{\rm LL}\), \(\mathcal F^{\rm B}_{\rm LL}\),
\(\mathcal F^{\rm C}_{\rm LL}\) & the three families of covariant
brackets: \(\mathcal F^{\rm A}_{\rm LL}\) on
\(\mathfrak l_{\rm LL}\oplus\mathfrak s_{\rm LL}\) (dimension 14),
\(\mathcal F^{\rm B}_{\rm LL}\) on
\(\mathfrak l_{\rm LL}\oplus\mathfrak p_{\rm LL}\oplus\mathfrak s_{\rm LL}\)
(18), \(\mathcal F^{\rm C}_{\rm LL}\) on
\(\mathfrak l_{\rm LL}\oplus\mathfrak{su}(2)\oplus\mathfrak s_{\rm LL}\)
(17) & new & \texttt{F\_A} \texttt{F\_B} \texttt{F\_C} \\
\(\mathfrak p_{\rm LL}\); \(\mathbf r\), \(\mathbf r'\) & the
translations of the family \(\mathcal F^{\rm B}_{\rm LL}\): the real
space of the Hermitian bi-quaternions \(\mathbf r=\mathbf r^\dagger\)
(dimension 4), on which \(\mathfrak l_{\rm LL}\) acts by
\([\mathbf u,\mathbf r]=\mathbf u\mathbf r+\mathbf r\mathbf u^\dagger\)
& new & \texttt{p\_LL} \texttt{rP} \texttt{rP\textquotesingle{}} \\
\(c^{\rm tw}_{\rm LL}\), \(c^{\rm sl}_{\rm LL}\),
\(c^{\rm PP}_{\rm LL}\), \(c^{\rm Ps}_{\rm LL}\),
\(c^{\rm sP}_{\rm LL}\), \(c^{\rm TT}_{\rm LL}\),
\(c^{\rm Ts}_{\rm LL}\), \(c^{\rm sT}_{\rm LL}\) & the coefficients of
the families:
\([\mathbf u,\mathbf b\mathbf e_4]=c^{\rm tw}_{\rm LL}(W_{\rm LL}(\mathbf u)\mathbf b)\mathbf e_4\)
(the twist),
\([\mathbf b\mathbf e_4,\mathbf b'\mathbf e_4]\supset c^{\rm sl}_{\rm LL}(\bar{\mathbf b}\mathbf b'-\bar{\mathbf b}'\mathbf b)\),
\([\mathbf r,\mathbf r']=c^{\rm PP}_{\rm LL}(\mathbf r\bar{\mathbf r}'-\mathbf r'\bar{\mathbf r})\),
\([\mathbf r,\mathbf b\mathbf e_4]=c^{\rm Ps}_{\rm LL}(\mathbf b^*\bar{\mathbf r})\mathbf e_4\),
\([\mathbf b\mathbf e_4,\mathbf b'\mathbf e_4]\supset c^{\rm sP}_{\rm LL}\,i(\bar{\mathbf b}\mathbf b'^*-\bar{\mathbf b}'\mathbf b^*)\),
\([\mathbf t,\mathbf t']=c^{\rm TT}_{\rm LL}(\mathbf t\mathbf t'-\mathbf t'\mathbf t)\),
\([\mathbf t,\mathbf b\mathbf e_4]=c^{\rm Ts}_{\rm LL}(\mathbf t\mathbf b)\mathbf e_4\),
\([\mathbf b\mathbf e_4,\mathbf b'\mathbf e_4]\supset\mathrm{Re}[c^{\rm sT}_{\rm LL}(\mathbf b\bar{\mathbf b}'-\mathbf b'\bar{\mathbf b})]\);
complex except \(c^{\rm sP}_{\rm LL}\), \(c^{\rm TT}_{\rm LL}\) (real);
plain text \texttt{ctw\_re}, \texttt{ctw\_im}: real and imaginary part
of \(c^{\rm tw}_{\rm LL}\), and likewise & new & \texttt{c\_tw}
\texttt{c\_sl} \texttt{c\_PP} \texttt{c\_Ps} \texttt{c\_sP}
\texttt{c\_TT} \texttt{c\_Ts} \texttt{c\_sT} \texttt{csP} \texttt{cTT}
\texttt{ctw\_re} \texttt{ctw\_im} \texttt{csl\_re} \texttt{csl\_im}
\texttt{cPP\_re} \texttt{cPP\_im} \texttt{cPs\_re} \texttt{cPs\_im}
\texttt{cTs\_re} \texttt{cTs\_im} \texttt{csT\_re} \texttt{csT\_im} \\
\(\lambda^{\rm sc}_{\rm LL}\) & a complex number \(\neq0\) by which the
spinor charges are rescaled,
\(\mathbf b\mapsto\lambda^{\rm sc}_{\rm LL}\mathbf b\) & new &
\texttt{lambda\_sc} \\
\(\mathcal I^{\rm Jac}_{\rm LL}\), \(\mathcal I^{\rm Mal}_{\rm LL}\);
\(w^{\rm sat}_{\rm LL}\) & the ideals of polynomials in the real
coefficients of a family generated by the Jacobi and by the Malcev
identity; \(w^{\rm sat}_{\rm LL}\) the auxiliary variable of the
saturation (the points of the variety at which a given polynomial does
not vanish are the points of the variety of the ideal enlarged by that
polynomial times \(w^{\rm sat}_{\rm LL}\), minus 1) & new &
\texttt{w\_sat} \\
\end{longtable}

\textbf{5 The manifolds}

\begin{longtable}[]{@{}
  >{\raggedright\arraybackslash}p{(\columnwidth - 6\tabcolsep) * \real{0.2200}}
  >{\raggedright\arraybackslash}p{(\columnwidth - 6\tabcolsep) * \real{0.5000}}
  >{\raggedright\arraybackslash}p{(\columnwidth - 6\tabcolsep) * \real{0.0800}}
  >{\raggedright\arraybackslash}p{(\columnwidth - 6\tabcolsep) * \real{0.2000}}@{}}
\toprule\noalign{}
\begin{minipage}[b]{\linewidth}\raggedright
symbol
\end{minipage} & \begin{minipage}[b]{\linewidth}\raggedright
meaning
\end{minipage} & \begin{minipage}[b]{\linewidth}\raggedright
origin
\end{minipage} & \begin{minipage}[b]{\linewidth}\raggedright
plain-text names
\end{minipage} \\
\midrule\noalign{}
\endhead
\bottomrule\noalign{}
\endlastfoot
\(\mathbb R^{1,7}_{\rm LL}\), \(\mathbb R^{1,3}_{\rm LL}\);
\(x^p_{\rm LL}\) & the Hermitian bi-octonions
\(\mathbf x=\mathbf x^\dagger=x^0_{\rm LL}\mathbf e_0+i\sum_{p=1}^7x^p_{\rm LL}\mathbf e_p\)
(real \(x^p_{\rm LL}\)): the \(7+1\) Minkowski space with the real norm
\(N_{\rm LL}(\mathbf x)=(x^0_{\rm LL})^2-\sum_p(x^p_{\rm LL})^2\); and
its subspace of Hermitian bi-quaternions (\(p\le3\)), Minkowski space of
\(3+1\) dimensions & new & \texttt{R\^{}(1,7)\_LL}
\texttt{R\^{}(1,3)\_LL} \texttt{x\_LL\^{}0} \texttt{x\_LL\^{}p}
\texttt{x\_LL\^{}1} \texttt{x\_LL\^{}2} \texttt{x\_LL\^{}3}
\texttt{x\_LL\^{}j} \\
\(\mathbb O^{\rm s}_{\rm LL}\) & the split octonions inside
\(\mathbb B_{\rm LL}\): the real space
\(\mathbb H+(i\mathbb H)\mathbf e_4\) (real combinations of
\(\mathbf e_0,\dots,\mathbf e_3,i\mathbf e_4,\dots,i\mathbf e_7\)),
closed under the product, on which \(N_{\rm LL}\) has signature
\((4,4)\): the \(4+4\) manifold (plain text \texttt{R\^{}(4,4)}) & new &
\texttt{Os} \texttt{R\^{}(4,4)} \\
\(S^7_{\rm LL}\), \(S^{3,4}_{\rm LL}\) & the unit elements
(\(N_{\rm LL}=1\)) of \(\mathbb O\) (the seven-sphere) and of
\(\mathbb O^{\rm s}_{\rm LL}\) (the pseudo-sphere): the two real forms
of \(\mathbb L_{\rm LL}\), both Moufang loops & new & \texttt{S\^{}7}
\texttt{S\^{}(3,4)} \\
\(X^{\rm LL}_{\mathbf x}\) & the frame field
\(X^{\rm LL}_{\mathbf x}(\mathbf w)=\mathbf w\mathbf x\) on
\(\mathbb L_{\rm LL}\) for imaginary \(\mathbf x\) (a vector field: the
velocity of the curve \(\mathbf w(\theta_{\rm LL})\) with
\(d\mathbf w/d\theta_{\rm LL}=\mathbf w\mathbf x\)) & new &
\texttt{X\_x} \texttt{X\_y} \\
\(\mathsf L_{\rm LL}(\mathbf g)\) & the real \(4\times4\) matrix of the
Lorentz transformation
\(\mathbf x\mapsto\mathbf g\mathbf x\mathbf g^\dagger\) of
\(\mathbb R^{1,3}_{\rm LL}\) in the coordinates
\(x^0_{\rm LL},\dots,x^3_{\rm LL}\), for a unit bi-quaternion
\(\mathbf g\): its entry in row \(p\) and column \(p'\)
(\(p,p'=0,\dots,3\)) is
\(\tfrac12\,\mathrm{Re}\,\mathrm{tr}(\sigma^{\rm LL}_pM_{\rm LL}(\mathbf g)\,\sigma^{\rm LL}_{p'}M_{\rm LL}(\mathbf g)^\dagger)\),
with \(\sigma^{\rm LL}_0\) the unit matrix & new & \texttt{Lmat} \\
\(\chi_{\rm LL}\) & a two-component spinor (a column of two complex
numbers) on which a unit bi-quaternion \(\mathbf g\) acts by
\(\chi_{\rm LL}\mapsto M_{\rm LL}(\mathbf g)\chi_{\rm LL}\); the
Hermitian matrix \(\chi_{\rm LL}\chi_{\rm LL}^\dagger\) is the matrix
\(M_{\rm LL}(\mathbf x)\) of a point \(\mathbf x\) of the future light
cone of \(\mathbb R^{1,3}_{\rm LL}\) & new & \texttt{chi} \\
\(\theta_{\rm LL}\) & the real parameter of a one-parameter family of
transformations or of a curve (an angle, a rapidity) & new &
\texttt{theta} \\
\(\mathcal V_{\rm LL}(\mathbf u)\) & the vector action of a Lorentz
element on \(\mathbb R^{1,3}_{\rm LL}\),
\(\mathcal V_{\rm LL}(\mathbf u)\mathbf x=\mathbf u\mathbf x+\mathbf x\mathbf u^\dagger\),
a real \(4\times4\) matrix with entries
\(\mathcal V_{\rm LL}(\mathbf u)^p{}_{p'}\) in the coordinates
\(x^p_{\rm LL}\) (plain text \texttt{V\_LL\^{}p\_p\textquotesingle{}}) &
new & \texttt{V\_LL} \\
\end{longtable}

\textbf{6 The spin-zero superfluid}

\begin{longtable}[]{@{}
  >{\raggedright\arraybackslash}p{(\columnwidth - 6\tabcolsep) * \real{0.2200}}
  >{\raggedright\arraybackslash}p{(\columnwidth - 6\tabcolsep) * \real{0.5000}}
  >{\raggedright\arraybackslash}p{(\columnwidth - 6\tabcolsep) * \real{0.0800}}
  >{\raggedright\arraybackslash}p{(\columnwidth - 6\tabcolsep) * \real{0.2000}}@{}}
\toprule\noalign{}
\begin{minipage}[b]{\linewidth}\raggedright
symbol
\end{minipage} & \begin{minipage}[b]{\linewidth}\raggedright
meaning
\end{minipage} & \begin{minipage}[b]{\linewidth}\raggedright
origin
\end{minipage} & \begin{minipage}[b]{\linewidth}\raggedright
plain-text names
\end{minipage} \\
\midrule\noalign{}
\endhead
\bottomrule\noalign{}
\endlastfoot
\(M_{\rm LL}(\mathbf b)\) & the complex \(2\times2\) matrix of a
bi-quaternion, \(M_{\rm LL}(\mathbf e_0)=\sigma^{\rm LL}_0\),
\(M_{\rm LL}(\mathbf e_j)=-i\sigma^{\rm LL}_j\), extended
complex-linearly;
\(M_{\rm LL}(\mathbf b\mathbf b')=M_{\rm LL}(\mathbf b)M_{\rm LL}(\mathbf b')\),
\(\det M_{\rm LL}(\mathbf b)=N_{\rm LL}(\mathbf b)\) (plain text
\texttt{M\_2(C)}: all complex \(2\times2\) matrices) & new &
\texttt{M\_LL} \texttt{M\_2(C)} \\
\(\psi'_{\rm LL}\) & the Weyl isodoublet of a spinor charge
\(\mathbf b\mathbf e_4\): the \(2\times2\) matrix
\(\psi'_{\rm LL}=\mathsf E_{\rm LL}^{-1}M_{\rm LL}(\mathbf b)^T\), whose
first column (isospin up) and second column (isospin down) are Weyl
spinors & new & \texttt{psi\textquotesingle{}\_LL} \\
\(\Psi_{\rm LL}\) & the Dirac isodoublet of the spinor charges: a field
of 8 complex Grassmann components (Dirac \(\otimes\) isospin), whose
\(\gamma^{\rm ch}_{\rm LL}=+1\) half is built from
\(\mathfrak s_{\rm LL}\) and whose other half from the complex-conjugate
super-loop & new & \texttt{Psi\_LL} \texttt{Psi\_LL\^{}T}
\texttt{Psi\_LL\^{}dag} \\
\(\gamma^p_{\rm LL}\), \(\gamma^{\rm ch}_{\rm LL}\), \(C_{\rm LL}\),
\(\alpha^j_{\rm LL}\), \(\beta_{\rm LL}\) & the \(4\times4\) Dirac
matrices of \(\mathbb R^{1,3}_{\rm LL}\) in the chiral basis,
\(\gamma^0_{\rm LL}=\begin{pmatrix}0&\sigma^{\rm LL}_0\\\sigma^{\rm LL}_0&0\end{pmatrix}\),
\(\gamma^j_{\rm LL}=\begin{pmatrix}0&\sigma^{\rm LL}_j\\-\sigma^{\rm LL}_j&0\end{pmatrix}\);
the chirality matrix
\(\gamma^{\rm ch}_{\rm LL}=i\gamma^0_{\rm LL}\gamma^1_{\rm LL}\gamma^2_{\rm LL}\gamma^3_{\rm LL}=\mathrm{diag}(-1,-1,1,1)\);
the charge-conjugation matrix
\(C_{\rm LL}=i\gamma^2_{\rm LL}\gamma^0_{\rm LL}\);
\(\alpha^j_{\rm LL}=\gamma^0_{\rm LL}\gamma^j_{\rm LL}\),
\(\beta_{\rm LL}=\gamma^0_{\rm LL}\) (on Dirac \(\otimes\) isospin each
acts as \(\otimes\sigma^{\rm LL}_0\)) & new & \texttt{gamma\_LL\^{}p}
\texttt{gamma\_LL\^{}p\textquotesingle{}} \texttt{gamma\_LL\^{}0}
\texttt{gamma\_LL\^{}1} \texttt{gamma\_LL\^{}2} \texttt{gamma\_LL\^{}3}
\texttt{gamma\_LL\^{}j} \texttt{gamma\_LL\^{}j\textquotesingle{}}
\texttt{gamma\_LL\^{}j\textquotesingle{}\textquotesingle{}}
\texttt{gamma\_ch} \texttt{C\_LL} \texttt{alpha\_LL\^{}j}
\texttt{alpha\_LL\^{}1} \texttt{alpha\_LL\^{}2} \texttt{alpha\_LL\^{}3}
\texttt{beta\_LL} \\
\(\mathsf K^\iota_{\rm LL}\), \(\mathsf K'^\iota_{\rm LL}\),
\(\mathsf K^{\rm s}_{\rm LL}\) & the \(8\times8\) pair matrices of the
spin-zero channels: scalar
\(\mathsf K^\iota_{\rm LL}=C_{\rm LL}\gamma^{\rm ch}_{\rm LL}\otimes\mathsf E_{\rm LL}\tau^{\rm LL}_\iota\)
(\(0^+\)), pseudoscalar
\(\mathsf K'^\iota_{\rm LL}=C_{\rm LL}\otimes\mathsf E_{\rm LL}\tau^{\rm LL}_\iota\)
(\(0^-\)), and the \(16\times16\) singlet of two species
\(\mathsf K^{\rm s}_{\rm LL}=C_{\rm LL}\gamma^{\rm ch}_{\rm LL}\otimes\mathsf E_{\rm LL}\otimes\mathsf E_{\rm LL}\)
(the last factor on the species index); the plain-text name Kch alone,
and the letter \(\mathsf K\) without scripts (in Theorem SF1 and in the
explanation of the folder explanations), stand for any pair matrix
(\(8\times8\) for one field, \(16\times16\) for two species), the pair
of two spinor charges being \(\Psi^T_{\rm LL}\mathsf K\Psi_{\rm LL}\);
\(\mathsf K_{12}\) is the entry of \(\mathsf K\) in row 1 and column 2 &
new & \texttt{Kch} \texttt{Kch\_iota}
\texttt{Kch\_iota\textquotesingle{}} \texttt{Kch\_1} \texttt{Kch\_2}
\texttt{Kch\_3} \texttt{Kch\textquotesingle{}\_iota}
\texttt{Kch\textquotesingle{}\_1} \texttt{Kch\textquotesingle{}\_2}
\texttt{Kch\textquotesingle{}\_3} \texttt{Kch\_singlet} \\
\(\Pi^\iota_{\rm LL}\) & the pair operators
\(\Pi^\iota_{\rm LL}=\sum_{\vec q_{\rm LL}}\Psi_{\rm LL}(\vec q_{\rm LL})^T\mathsf K^\iota_{\rm LL}\Psi_{\rm LL}(-\vec q_{\rm LL})\),
which annihilate two spinor charges of opposite momenta & new &
\texttt{Pi\_iota} \\
\(\vec q_{\rm LL}\), \(q_{\rm LL}\), \(q^j_{\rm LL}\) & the momentum of
a particle in the space of \(\mathbb R^{1,3}_{\rm LL}\), its length and
its components & new & \texttt{q} \\
\(m_{\rm LL}\), \(\mu_{\rm LL}\), \(\Lambda_{\rm LL}\) & the mass of the
spinor charges, the chemical potential, and the momentum cutoff
(\(q_{\rm LL}<\Lambda_{\rm LL}\); \(\Lambda_{\rm LL}=1\) sets the units,
except where it varies) & new & \texttt{m} \texttt{mu}
\texttt{Lambda} \\
\(\varepsilon_{\rm LL}\), \(E^\pm_{\rm LL}\), \(\tilde E^\pm_{\rm LL}\)
& the energy \(\varepsilon_{\rm LL}=(q_{\rm LL}^2+m_{\rm LL}^2)^{1/2}\),
the quasi-particle energies of the scalar channel
\(E^\pm_{\rm LL}=[(\varepsilon_{\rm LL}\pm\mu_{\rm LL})^2+\Delta_{\rm LL}^2]^{1/2}\),
a function of the gap (written \(E^\pm_{\rm LL}(\Delta^+_{\rm LL})\) and
\(E^\pm_{\rm LL}(\Delta^-_{\rm LL})\) when \(\Delta_{\rm LL}\) is
replaced by one of the two gaps of a general state), and those of the
pseudoscalar channel,
\((\tilde E^\pm_{\rm LL})^2=\varepsilon_{\rm LL}^2+\mu_{\rm LL}^2+\Delta_{\rm LL}^2\pm2[\mu_{\rm LL}^2\varepsilon_{\rm LL}^2+m_{\rm LL}^2\Delta_{\rm LL}^2]^{1/2}\)
& new & \texttt{eps} \texttt{E\^{}+} \texttt{E\^{}-} \texttt{E\^{}(+-)}
\texttt{E\_minus\_polar} \texttt{E\_plus\_polar} \texttt{Etilde\^{}+}
\texttt{Etilde\^{}-} \texttt{Etilde\^{}(+-)} \texttt{Etilde\_minus}
\texttt{Etilde\_plus} \\
\(G_{\rm LL}\), \(G^{\rm crit}_{\rm LL}\) & the coupling of the
attractive spin-zero contact interaction (dimension mass\(^{-2}\)) and
its critical value at \(\mu_{\rm LL}=0\), above which the vacuum
condenses & new & \texttt{G} \texttt{G\_crit}
\texttt{G\_over\_Gcrit\_m0} \texttt{G\_crit\_mu0\_m0} \\
\(\mathcal H^{\rm sf}_{\rm LL}\) & the Hamiltonian of the superfluid,
\(\sum_{\vec q_{\rm LL}}\Psi^\dagger_{\rm LL}(\mathsf h_{\rm LL}-\mu_{\rm LL})\Psi_{\rm LL}-\tfrac14G_{\rm LL}\sum_\iota(\Pi^\iota_{\rm LL})^\dagger\Pi^\iota_{\rm LL}\)
& new & \texttt{Hsf} \\
\(\mathsf h_{\rm LL}(\vec q_{\rm LL})\) & the \(8\times8\) Hamiltonian
matrix of a plane wave,
\(\sum_j\alpha^j_{\rm LL}q^j_{\rm LL}+\beta_{\rm LL}m_{\rm LL}\) (times
\(\sigma^{\rm LL}_0\) on the isospin index); eigenvalues
\(\pm\varepsilon_{\rm LL}\) & new & \texttt{h(q)} \\
\(\vec d_{\rm LL}\), \(d^\iota_{\rm LL}\); \(\zeta_{\rm LL}\) & the
order parameter, the isospin vector
\(\vec d_{\rm LL}=(d^1_{\rm LL},d^2_{\rm LL},d^3_{\rm LL})\in\mathbb C^3\)
with
\(d^\iota_{\rm LL}=-\tfrac12G_{\rm LL}\langle\Pi^\iota_{\rm LL}\rangle\),
and its non-unitarity
\(\zeta_{\rm LL}=\lvert\vec d_{\rm LL}\times\vec d^{\,*}_{\rm LL}\rvert/\lvert\vec d_{\rm LL}\rvert^2\in[0,1]\)
(plain text \texttt{cross(dvec,\ dvec\^{}*)}: the cross product) & new &
\texttt{dvec} \texttt{d\_iota} \texttt{abs\_dvec\_at\_min}
\texttt{zeta} \\
\(\Delta_{\rm LL}\), \(\Delta^\pm_{\rm LL}\) & the gap:
\(\Delta_{\rm LL}=\lvert\vec d_{\rm LL}\rvert\) for a unitary state, and
the two gaps
\((\Delta^\pm_{\rm LL})^2=\lvert\vec d_{\rm LL}\rvert^2(1\pm\zeta_{\rm LL})\)
of a general state & new & \texttt{Delta} \texttt{Delta\_+}
\texttt{Delta\_-} \texttt{Delta\_(+-)} \\
\(\mathsf P_{\rm LL}\); \(\mathsf H_{\rm LL}(\vec q_{\rm LL})\);
\(\mathcal N_{\rm LL}(\vec q_{\rm LL})\) & the pairing block
\(\mathsf P_{\rm LL}=\sum_\iota d^\iota_{\rm LL}(\mathsf K^\iota_{\rm LL})^\dagger\),
the \(16\times16\) Bogoliubov-de Gennes matrix
\(\mathsf H_{\rm LL}=\begin{pmatrix}\mathsf h_{\rm LL}(\vec q_{\rm LL})-\mu_{\rm LL}&\mathsf P_{\rm LL}\\\mathsf P_{\rm LL}^\dagger&-(\mathsf h_{\rm LL}(-\vec q_{\rm LL})-\mu_{\rm LL})^T\end{pmatrix}\)
and the Nambu spinor
\(\mathcal N_{\rm LL}=(\Psi_{\rm LL}(\vec q_{\rm LL}),\Psi_{\rm LL}^{*}(-\vec q_{\rm LL}))\)
on which it acts & new & \texttt{Ppair} \texttt{H\_BdG} \\
\(\mathsf D_{\rm LL}\) & the isospin part of the pairing block of the
scalar channel, the \(2\times2\) matrix
\(\mathsf D_{\rm LL}=\sum_\iota d^\iota_{\rm LL}(\mathsf E_{\rm LL}\tau^{\rm LL}_\iota)^\dagger\),
so that
\(\mathsf P_{\rm LL}=(C_{\rm LL}\gamma^{\rm ch}_{\rm LL})^\dagger\otimes\mathsf D_{\rm LL}\);
it is also the isospin wavefunction of the condensed pairs (rows and
columns: isospin up, down); plain text \texttt{Dpair} (not
\texttt{D\_LL}, which is the Lorentz derivation) & new &
\texttt{Dpair} \\
\(\Omega_{\rm LL}\) & the grand potential per volume at zero
temperature, measured from the normal state:
\(\Omega_{\rm LL}=\lvert\vec d_{\rm LL}\rvert^2/G_{\rm LL}-\tfrac12\int\!d^3q_{\rm LL}(2\pi)^{-3}\sum(\text{positive eigenvalues of }\mathsf H_{\rm LL})\)
minus its value at \(\vec d_{\rm LL}=0\) & new & \texttt{Omega}
\texttt{Omega\_min} \texttt{Omega\_minus\_Omega0}
\texttt{Omega\_at\_polar\_gap} \\
\(k^{\rm F}_{\rm LL}\), \(\nu^{\rm F}_{\rm LL}\) & the Fermi momentum
\((\mu_{\rm LL}^2-m_{\rm LL}^2)^{1/2}\) and the density of states of one
species and one spin state at the Fermi surface,
\(\nu^{\rm F}_{\rm LL}=k^{\rm F}_{\rm LL}\mu_{\rm LL}/(2\pi^2)\) & new &
\texttt{kF} \texttt{nuF} \\
\(F_{\rm LL}\); \(U(1)^{\rm F}_{\rm LL}\); \(SU(2)_{\rm v}\) & the
generator of the fermion number
\(\Psi_{\rm LL}\to e^{i\theta_{\rm LL}}\Psi_{\rm LL}\) (it multiplies
\(\vec d_{\rm LL}\) by \(2i\)) and its group; the isospin group, which
rotates the isospin index of \(\Psi_{\rm LL}\) and \(\vec d_{\rm LL}\)
as a vector & new & \texttt{F\_LL} \texttt{U(1)\_F\_LL}
\texttt{SU(2)\_v} \\
\(\Psi^{(1)}_{\rm LL}\), \(\Psi^{(2)}_{\rm LL}\) & two species of the
field \(\Psi_{\rm LL}\), combined into one field of 16 components with a
species index & new & \\
\end{longtable}

\textbf{7 Power counting and renormalizability}

\begin{longtable}[]{@{}
  >{\raggedright\arraybackslash}p{(\columnwidth - 6\tabcolsep) * \real{0.2200}}
  >{\raggedright\arraybackslash}p{(\columnwidth - 6\tabcolsep) * \real{0.5000}}
  >{\raggedright\arraybackslash}p{(\columnwidth - 6\tabcolsep) * \real{0.0800}}
  >{\raggedright\arraybackslash}p{(\columnwidth - 6\tabcolsep) * \real{0.2000}}@{}}
\toprule\noalign{}
\begin{minipage}[b]{\linewidth}\raggedright
symbol
\end{minipage} & \begin{minipage}[b]{\linewidth}\raggedright
meaning
\end{minipage} & \begin{minipage}[b]{\linewidth}\raggedright
origin
\end{minipage} & \begin{minipage}[b]{\linewidth}\raggedright
plain-text names
\end{minipage} \\
\midrule\noalign{}
\endhead
\bottomrule\noalign{}
\endlastfoot
\([\ \cdot\ ]\) (one argument) & the mass dimension of a field or a
coupling, in units in which \(\hbar\) and the speed of light are 1 (a
bracket with two or three arguments is a commutator or an associator) &
new & \texttt{{[}Psi\_LL{]}} \texttt{{[}phi{]}} \texttt{{[}G{]}}
\texttt{{[}g\_Y{]}} \texttt{{[}lam4{]}} \\
\(\phi_{\rm LL}\), \(g^{\rm Y}_{\rm LL}\), \(\lambda^{(1)}_{\rm LL}\),
\(\lambda^{(2)}_{\rm LL}\) & the Yukawa completion: a complex
isospin-triplet pair field
\(\phi_{\rm LL}=(\phi^1_{\rm LL},\phi^2_{\rm LL},\phi^3_{\rm LL})\)
coupled to \(\Psi^T_{\rm LL}\mathsf K^\iota_{\rm LL}\Psi_{\rm LL}\) with
the Yukawa coupling \(g^{\rm Y}_{\rm LL}\), and the couplings
\(\lambda^{(1)}_{\rm LL}\), \(\lambda^{(2)}_{\rm LL}\) of the quartic
terms \((\phi^\dagger_{\rm LL}\phi_{\rm LL})^2\) and
\(\lvert\phi^T_{\rm LL}\phi_{\rm LL}\rvert^2\) (plain text
\texttt{lam4}: either of them) & new & \texttt{phi} \texttt{g\_Y}
\texttt{lam4} \texttt{lam4\_1} \texttt{lam4\_2} \\
\(\mathcal L^{\rm Y}_{\rm LL}\), \(\mathcal L^{\rm c}_{\rm LL}\) & the
Lagrangian densities of the Yukawa form and of the contact form of the
field theory of the spinor charges on \(\mathbb R^{1,3}_{\rm LL}\) & new
& \texttt{Lagr\_Y} \texttt{Lagr\_c} \\
\(m^\phi_{\rm LL}\) & the mass of the pair field \(\phi_{\rm LL}\) in
the Yukawa form & new & \texttt{m\_phi} \\
\(I^{\rm ren}_{\rm LL}\) & the renormalized one-loop bubble,
\(I^{\rm ren}_{\rm LL}(p^{\rm E}_{\rm LL})=I^{\rm bub}_{\rm LL}(p^{\rm E}_{\rm LL})-I^{\rm bub}_{\rm LL}(0)-(p^{\rm E}_{\rm LL})^2\,\partial I^{\rm bub}_{\rm LL}/\partial(p^{\rm E}_{\rm LL})^2\)
(the derivative taken at \(p^{\rm E}_{\rm LL}=0\)): the bubble minus its
mass and wave-function counterterms & new & \texttt{Iren} \\
\(\deg_{\rm LL}\) & the superficial degree of divergence of a
one-particle-irreducible graph: the power of the cutoff with which it
grows when all its loop momenta grow together & new & \texttt{deg} \\
\(n^{\rm loop}_{\rm LL}\), \(n^{\rm int,F}_{\rm LL}\),
\(n^{\rm int,B}_{\rm LL}\), \(n^{\rm ext,F}_{\rm LL}\),
\(n^{\rm ext,B}_{\rm LL}\), \(n^{\rm v}_{\rm LL}\),
\(n^{\rm v,Y}_{\rm LL}\), \(n^{\rm v,4}_{\rm LL}\) & the numbers of
loops, of internal fermion and boson lines, of external fermion and
boson lines, of vertices, of Yukawa vertices and of quartic vertices of
a graph & new & \texttt{nloop} \texttt{nint\_F} \texttt{nint\_B}
\texttt{next\_F} \texttt{next\_B} \texttt{nvert} \texttt{nvert\_Y}
\texttt{nvert\_4} \\
\(I^{\rm bub}_{\rm LL}\); \(k^{\rm E}_{\rm LL}\), \(p^{\rm E}_{\rm LL}\)
& the one-loop bubble of the Yukawa completion,
\(I^{\rm bub}_{\rm LL}(p^{\rm E}_{\rm LL})=\int_{k^{\rm E}_{\rm LL}<\Lambda_{\rm LL}}\frac{d^4k^{\rm E}_{\rm LL}}{(2\pi)^4}\frac{k^{\rm E}_{\rm LL}\cdot(k^{\rm E}_{\rm LL}+p^{\rm E}_{\rm LL})+m_{\rm LL}^2}{[(k^{\rm E}_{\rm LL})^2+m_{\rm LL}^2][(k^{\rm E}_{\rm LL}+p^{\rm E}_{\rm LL})^2+m_{\rm LL}^2]}\),
with the Euclidean loop and external four-momenta
\(k^{\rm E}_{\rm LL}\), \(p^{\rm E}_{\rm LL}\) (and their lengths) & new
& \texttt{Ibub} \texttt{Ibub\_0} \texttt{Ibub\_pE}
\texttt{Ibub\_pE\_minus\_Ibub\_0} \texttt{kE} \texttt{pE} \\
\end{longtable}

\textbf{8 Groups, Lie algebras, operations and labels}

\begin{longtable}[]{@{}
  >{\raggedright\arraybackslash}p{(\columnwidth - 6\tabcolsep) * \real{0.2200}}
  >{\raggedright\arraybackslash}p{(\columnwidth - 6\tabcolsep) * \real{0.5000}}
  >{\raggedright\arraybackslash}p{(\columnwidth - 6\tabcolsep) * \real{0.0800}}
  >{\raggedright\arraybackslash}p{(\columnwidth - 6\tabcolsep) * \real{0.2000}}@{}}
\toprule\noalign{}
\begin{minipage}[b]{\linewidth}\raggedright
symbol
\end{minipage} & \begin{minipage}[b]{\linewidth}\raggedright
meaning
\end{minipage} & \begin{minipage}[b]{\linewidth}\raggedright
origin
\end{minipage} & \begin{minipage}[b]{\linewidth}\raggedright
plain-text names
\end{minipage} \\
\midrule\noalign{}
\endhead
\bottomrule\noalign{}
\endlastfoot
\(SL(2,\mathbb C)\), \(Spin(3,1)^+\), \(SO(4,\mathbb C)\),
\(\mathrm G_2\), \(\mathrm G_2(\mathbb C)\), \(SO(1,7)\), \(Spin(7)\),
\(Spin(7,\mathbb C)\), \(SL(2,\mathbb R)\), \(Spin(2,1)\), \(SU(2)\),
\(U(1)\), \(\mathbb Z_2\) & groups: the complex \(2\times2\) matrices of
determinant 1 (the unit bi-quaternions), equal to the connected double
cover of the Lorentz group of \(3+1\) dimensions; the complex orthogonal
group of four dimensions; the compact exceptional group \(\mathrm G_2\)
(automorphisms of \(\mathbb O\)) and its complex form (automorphisms of
\(\mathbb B_{\rm LL}\)); the Lorentz group of \(7+1\) dimensions; the
spin groups of seven dimensions (compact and complex); the real
\(2\times2\) matrices of determinant 1, the double cover of the Lorentz
group of \(2+1\) dimensions; the special unitary group of two
dimensions; the phases; the group of two elements & new &
\texttt{SL(2,C)} \texttt{Spin(3,1)\^{}+} \texttt{SO(4,C)} \texttt{G2}
\texttt{G2(C)} \texttt{SO(1,7)} \texttt{SO(3,1)} \texttt{Spin(7)}
\texttt{Spin(7,C)} \texttt{SL(2,R)} \texttt{Spin(2,1)} \texttt{SU(2)}
\texttt{U(1)} \texttt{Z2} \\
\(\mathfrak{so}(3,1)\), \(\mathfrak{so}(1,7)\), \(\mathfrak{so}(4,1)\),
\(\mathfrak{so}(8)\), \(\mathfrak{so}(4,4)\),
\(\mathfrak{so}(8,\mathbb C)\), \(\mathfrak{so}(3)\),
\(\mathfrak{so}(4)\), \(\mathfrak{so}(2,1)\),
\(\mathfrak{sl}(2,\mathbb R)\), \(\mathfrak{sl}(2)\),
\(\mathfrak{su}(2)\), \(\mathfrak{spin}(7)\), \(\mathfrak g_2\),
\(\mathfrak g_2(\mathbb C)\) & Lie algebras: of the Lorentz groups of
\(3+1\) and \(7+1\) dimensions; of the de Sitter group (signature of the
Killing form: 4 positive, 6 negative); of the linear maps that preserve
a definite real, a \((4,4)\) and a complex quadratic form in eight
dimensions; of the rotations of three-dimensional space and of
four-dimensional Euclidean space; of the Lorentz group of \(2+1\)
dimensions, equal to \(\mathfrak{sl}(2,\mathbb R)\);
\(\mathfrak{sl}(2)\) the simple Lie algebra of dimension 3 over any
field; \(\mathfrak{su}(2)\) isospin; \(\mathfrak{spin}(7)\) of dimension
21; the exceptional Lie algebra \(\mathfrak g_2\) (compact, dimension
14) and its complex form (real dimension 28) & new & \texttt{so(3,1)}
\texttt{so(1,7)} \texttt{so(4,1)} \texttt{so(8)} \texttt{so(4,4)}
\texttt{so(8,C)} \texttt{so(3)} \texttt{so(4)} \texttt{so(2,1)}
\texttt{sl(2,R)} \texttt{sl(2)} \texttt{su(2)} \texttt{spin(7)}
\texttt{g2} \texttt{g2(C)} \\
\(\mathrm{tr}\), \(\det\), \(\mathrm{diag}\), \(\mathrm{Re}\),
\(\mathrm{Im}\) & trace, determinant, the diagonal matrix with the
listed entries, real part, imaginary part (also: the imaginary part of
an algebra, the span of its imaginary units); plain text \texttt{exp},
\texttt{ln}, \texttt{sqrt}, \texttt{sum}, \texttt{int}, \texttt{span},
\texttt{dim}: the exponential, natural logarithm, square root, a sum, an
integral, the set of real linear combinations, a dimension & new &
\texttt{tr} \texttt{det} \texttt{diag} \texttt{Re} \texttt{Im}
\texttt{exp} \texttt{ln} \texttt{sqrt} \texttt{sum} \texttt{int}
\texttt{span} \texttt{dim} \\
\({}^T\); \(\otimes\), \(\oplus\); \(\circ\); \(\langle\ \rangle\); h.c.
& transpose; tensor product and direct sum (plain text \texttt{(x)} for
\(\otimes\)); composition of maps; expectation value in the ground
state; the Hermitian conjugate of the preceding terms & new &
\texttt{(x)} \texttt{\^{}T} \texttt{o} \\
A-1, \ldots, R-6; G-1, \ldots, G-12; LL-S1, \ldots, LL-H1; X-1, \ldots,
X-10; Z-1, \ldots, Z-8; Theorems R, L, M, SF1, \ldots, SF4 & the
identifiers of the machine checks (letters: A associativity, L Lorentz
super-loop, Q operators and eigenvalues, C the inversion element, M
Malcev classification, G Groebner bases, N manifolds, S superfluid, R
renormalizability, LL-S symbols, LL-H self-containedness, X the examples
of the explanation of the spin-zero superfluid in the folder
explanations, Z the 3+1 manifold and the renormalizability of the Yukawa
form in the folder explanations) and the labels of the theorems proved
in the manuscript & new & \\
H-grp, H-op, H-Lie, H-gr, H-dyn & the labels of the hypotheses of the
no-go theorems: the symmetries form a group (H-grp), act as linear
operators on the space of states (H-op), have generators that form a Lie
algebra (H-Lie) or a Lie superalgebra (H-gr), and the dynamical
hypotheses on the scattering matrix (H-dyn) & new & \\
\end{longtable}

\textbf{9 Test 2: the gauge-covariant field equations of the spinor
field and their solutions}

\begin{longtable}[]{@{}
  >{\raggedright\arraybackslash}p{(\columnwidth - 6\tabcolsep) * \real{0.2200}}
  >{\raggedright\arraybackslash}p{(\columnwidth - 6\tabcolsep) * \real{0.5000}}
  >{\raggedright\arraybackslash}p{(\columnwidth - 6\tabcolsep) * \real{0.0800}}
  >{\raggedright\arraybackslash}p{(\columnwidth - 6\tabcolsep) * \real{0.2000}}@{}}
\toprule\noalign{}
\begin{minipage}[b]{\linewidth}\raggedright
symbol
\end{minipage} & \begin{minipage}[b]{\linewidth}\raggedright
meaning
\end{minipage} & \begin{minipage}[b]{\linewidth}\raggedright
origin
\end{minipage} & \begin{minipage}[b]{\linewidth}\raggedright
plain-text names
\end{minipage} \\
\midrule\noalign{}
\endhead
\bottomrule\noalign{}
\endlastfoot
\(\eta^{\rm LL}_{pp'}\) & the Minkowski metric of
\(\mathbb R^{1,3}_{\rm LL}\) in the coordinates \(x^p_{\rm LL}\):
\(\eta^{\rm LL}_{00}=1\),
\(\eta^{\rm LL}_{11}=\eta^{\rm LL}_{22}=\eta^{\rm LL}_{33}=-1\) and
\(\eta^{\rm LL}_{pp'}=0\) for \(p\neq p'\); it is its own inverse and
raises and lowers the indices \(p\), \(p'\) & new & \texttt{eta} \\
\(x_{\rm LL}\), \(\vec x_{\rm LL}\); \(\partial_p\) & a point of
\(\mathbb R^{1,3}_{\rm LL}\) given by its coordinates,
\(x_{\rm LL}=(x^0_{\rm LL},x^1_{\rm LL},x^2_{\rm LL},x^3_{\rm LL})\)
(its Hermitian bi-quaternion is
\(x^0_{\rm LL}\mathbf e_0+i\sum_jx^j_{\rm LL}\mathbf e_j\)), and its
space part \(\vec x_{\rm LL}=(x^1_{\rm LL},x^2_{\rm LL},x^3_{\rm LL})\);
\(x^0_{\rm LL}\) is the time; \(\partial_p\) is the partial derivative
with respect to \(x^p_{\rm LL}\) (plain text \texttt{d\_p}) & new &
\texttt{x\_LL} \texttt{xvec} \texttt{d\_p} \\
\(\mathbf f_p\) & the frame of \(\mathbb R^{1,3}_{\rm LL}\): four
Hermitian bi-quaternions \(\mathbf f_0,\dots,\mathbf f_3\) with
\(N_{\rm LL}(\mathbf f_p,\mathbf f_{p'})=\eta^{\rm LL}_{pp'}\), possibly
depending on the point \(x_{\rm LL}\) (a vierbein); the canonical frame
is \(\mathbf f_0=\mathbf e_0\), \(\mathbf f_j=i\,\mathbf e_j\), for
which \(M_{\rm LL}(\mathbf f_p)=\sigma^{\rm LL}_p\),
\(M_{\rm LL}(\bar{\mathbf f}_0)=\sigma^{\rm LL}_0\) and
\(M_{\rm LL}(\bar{\mathbf f}_j)=-\sigma^{\rm LL}_j\) & new &
\texttt{f\_p} \texttt{f\_0} \texttt{f\_j} \\
\(\gamma^p_{\rm LL}[\mathbf f]\) & the Dirac matrices of a frame, the
\(4\times4\) matrices with the \(2\times2\) blocks
\(\gamma^p_{\rm LL}[\mathbf f]=\begin{pmatrix}0&M_{\rm LL}(\mathbf f_p)\\M_{\rm LL}(\bar{\mathbf f}_p)&0\end{pmatrix}\);
they satisfy
\(\gamma^p_{\rm LL}[\mathbf f]\gamma^{p'}_{\rm LL}[\mathbf f]+\gamma^{p'}_{\rm LL}[\mathbf f]\gamma^p_{\rm LL}[\mathbf f]=2N_{\rm LL}(\mathbf f_p,\mathbf f_{p'})\,\mathsf 1_4\);
for the canonical frame they are the matrices \(\gamma^p_{\rm LL}\) &
new & \texttt{gamma\_LL\^{}p{[}f{]}} \\
\(\tilde{\mathbf g}\) & for a unit bi-quaternion \(\mathbf g\):
\(\tilde{\mathbf g}=(\mathbf g^\dagger)^{-1}\), which equals the complex
conjugate \(\mathbf g^*\);
\(M_{\rm LL}(\tilde{\mathbf g})=(M_{\rm LL}(\mathbf g)^\dagger)^{-1}\);
the change of representation of the Dirac matrices by \(\mathbf g\)
changes the frame by
\(\mathbf f_p\mapsto\tilde{\mathbf g}\,\mathbf f_p\,\tilde{\mathbf g}^\dagger\)
& new & \texttt{gLt} \texttt{gLt\^{}dag} \\
\(\mathsf S_{\rm LL}(\mathbf g)\) & the spinor matrix of the Lorentz
transformation \(\mathbf g\): the \(4\times4\) matrix
\(\mathsf S_{\rm LL}(\mathbf g)=\mathrm{diag}(M_{\rm LL}(\tilde{\mathbf g}),M_{\rm LL}(\mathbf g))\)
of \(2\times2\) blocks (on the eight components of \(\Psi_{\rm LL}\):
\(\mathsf S_{\rm LL}(\mathbf g)\otimes\sigma^{\rm LL}_0\), plain text
\texttt{Sspin});
\(\mathsf S_{\rm LL}(\mathbf g)^{-1}\gamma^p_{\rm LL}\mathsf S_{\rm LL}(\mathbf g)=\sum_{p'}[\mathsf L_{\rm LL}(\mathbf g)]_{pp'}\gamma^{p'}_{\rm LL}\),
where \([\mathsf L_{\rm LL}(\mathbf g)]_{pp'}\) is the entry in row
\(p\) and column \(p'\) & new & \texttt{Sspin\_LL} \texttt{Sspin} \\
\(\mathsf s_{\rm LL}(\mathbf u)\) & a Lorentz generator \(\mathbf u\) of
\(\mathfrak l_{\rm LL}\) acting on the Dirac spinor: the \(4\times4\)
matrix
\(\mathsf s_{\rm LL}(\mathbf u)=\mathrm{diag}(-M_{\rm LL}(\mathbf u)^\dagger,M_{\rm LL}(\mathbf u))\),
the derivative of \(\mathsf S_{\rm LL}\) at the unit;
\(\mathsf s_{\rm LL}(J^j_{\rm LL})=\tfrac12\gamma^{j'}_{\rm LL}\gamma^{j''}_{\rm LL}\)
for cyclic \((j,j',j'')\) and
\(\mathsf s_{\rm LL}(K^j_{\rm LL})=\tfrac12\gamma^0_{\rm LL}\gamma^j_{\rm LL}\)
& new & \texttt{sgen\_LL} \\
\(\mathbf b^+_{\rm LL}\), \(\mathbf b^-_{\rm LL}\) & the two spinor
charges \(\mathbf b^\pm_{\rm LL}\mathbf e_4\) of which the Dirac
isodoublet is made: the upper block (chirality \(-1\)) of
\(\Psi_{\rm LL}\) is \(\psi'_{\rm LL}(\mathbf b^-_{\rm LL})\) and the
lower block (chirality \(+1\)) is
\(\psi'_{\rm LL}(\mathbf b^+_{\rm LL})\), where
\(\psi'_{\rm LL}(\mathbf y)=\mathsf E_{\rm LL}^{-1}M_{\rm LL}(\mathbf y)^T\)
for a bi-quaternion \(\mathbf y\);
\(\Psi_{\rm LL}[\mathbf b^-_{\rm LL},\mathbf b^+_{\rm LL}]\) is
\(\Psi_{\rm LL}\) written with them (plain text
\texttt{Psi(bSm,\ bSp)}); under the gauge group
\(\mathbf b^+_{\rm LL}\mapsto\exp(i\theta^{\rm F}_{\rm LL})\,\mathbf v\,\mathbf b^+_{\rm LL}\,\bar{\mathbf g}\),
the automorphism \(\varphi^{\rm LL}_{\mathbf g,\mathbf v}\) on the
spinor charge times a phase, and
\(\mathbf b^-_{\rm LL}\mapsto\exp(i\theta^{\rm F}_{\rm LL})\,\mathbf v\,\mathbf b^-_{\rm LL}\,\mathbf g^\dagger\),
its complex conjugate & new & \texttt{bSp} \texttt{bSm} \\
\(\theta^{\rm F}_{\rm LL}\) & the phase of a local transformation of the
fermion number,
\(\Psi_{\rm LL}\mapsto\exp(i\theta^{\rm F}_{\rm LL})\Psi_{\rm LL}\), a
real function of \(x_{\rm LL}\) & new & \texttt{thF} \\
\(\mathsf U_{\rm LL}\),
\(\mathsf R^{\rm v}_{\rm LL}(\mathsf U_{\rm LL})\) & a local isospin
rotation: a matrix \(\mathsf U_{\rm LL}\) of \(SU(2)\), a function of
\(x_{\rm LL}\), that acts on the isospin index of \(\Psi_{\rm LL}\)
(\(\mathsf U_{\rm LL}=M_{\rm LL}(\mathbf v)\) for a unit real quaternion
\(\mathbf v\)); \(\mathsf R^{\rm v}_{\rm LL}(\mathsf U_{\rm LL})\) is
the real orthogonal \(3\times3\) matrix with
\(\mathsf U_{\rm LL}\tau^{\rm LL}_\iota\mathsf U_{\rm LL}^\dagger=\sum_{\iota'}[\mathsf R^{\rm v}_{\rm LL}(\mathsf U_{\rm LL})]_{\iota'\iota}\tau^{\rm LL}_{\iota'}\),
the rotation of an isospin vector & new & \texttt{U\_LL}
\texttt{U\_LL\^{}dag} \texttt{Rv} \\
\(\mathcal U_{\rm LL}\) & a local gauge transformation of
\(\Psi_{\rm LL}\) by the gauge group
\(SL(2,\mathbb C)\times SU(2)_{\rm v}\times U(1)^{\rm F}_{\rm LL}\):
\(\mathcal U_{\rm LL}=\exp(i\theta^{\rm F}_{\rm LL})\,\mathsf S_{\rm LL}(\mathbf g)\otimes\mathsf U_{\rm LL}\),
with \(\mathbf g\), \(\mathsf U_{\rm LL}\) and
\(\theta^{\rm F}_{\rm LL}\) functions of \(x_{\rm LL}\) & new &
\texttt{Ucal} \\
\(A^\iota_{{\rm LL},p}\), \(A^{\rm F}_{{\rm LL},p}\);
\(g^{\rm v}_{\rm LL}\), \(g^{\rm F}_{\rm LL}\) & the gauge potentials,
real fields on \(\mathbb R^{1,3}_{\rm LL}\): \(A^\iota_{{\rm LL},p}\) of
the isospin group \(SU(2)_{\rm v}\) (three fields for each direction
\(p\)) and \(A^{\rm F}_{{\rm LL},p}\) of the fermion number
\(U(1)^{\rm F}_{\rm LL}\); \(g^{\rm v}_{\rm LL}\) and
\(g^{\rm F}_{\rm LL}\) are their couplings (real numbers); a constant
\(g^{\rm F}_{\rm LL}A^{\rm F}_{{\rm LL},0}\) acts like a chemical
potential & new & \texttt{Av\_LL} \texttt{Av\_LL\^{}iota}
\texttt{Av\_LL\^{}iota\textquotesingle{}}
\texttt{Av\_LL\^{}iota\textquotesingle{}\textquotesingle{}}
\texttt{AF\_LL} \texttt{g\_v} \texttt{g\_F} \\
\(\mathsf a_{{\rm LL},p}\) & the isospin potential as a Hermitian
\(2\times2\) matrix on the isospin index,
\(\mathsf a_{{\rm LL},p}=\sum_\iota A^\iota_{{\rm LL},p}\tau^{\rm LL}_\iota/2\)
& new & \texttt{atau\_LL} \\
\(\mathbf t^{\rm A}_{{\rm LL},p}\) & the isospin potential as a real
imaginary quaternion,
\(\mathbf t^{\rm A}_{{\rm LL},p}=\sum_\iota A^\iota_{{\rm LL},p}\mathbf e_\iota/2\)
(an element of \({\rm Im}\,\mathbb H\), like \(\mathbf t\)), with
\(M_{\rm LL}(\mathbf t^{\rm A}_{{\rm LL},p})=-i\,\mathsf a_{{\rm LL},p}\);
it acts on the spinor charges through the isospin derivation
\(D^{\rm v}_{\rm LL}\) & new & \texttt{tA} \texttt{tA\_p} \\
\(\mathbf u^\omega_{{\rm LL},p}\); \(\omega^j_{{\rm LL},p}\),
\(\varpi^j_{{\rm LL},p}\) & the Lorentz connection: for each direction
\(p\) an element of \(\mathfrak l_{\rm LL}\),
\(\mathbf u^\omega_{{\rm LL},p}=\sum_j(\omega^j_{{\rm LL},p}J^j_{\rm LL}+\varpi^j_{{\rm LL},p}K^j_{\rm LL})\),
with the real rotation components \(\omega^j_{{\rm LL},p}\) and the real
boost components \(\varpi^j_{{\rm LL},p}\);
\(\mathbf u^{\omega*}_{{\rm LL},p}\) is its complex conjugate (plain
text \texttt{uLw\_p\^{}*}) & new & \texttt{uLw} \texttt{uLw\_p}
\texttt{uLw\^{}*} \texttt{omega\_LL} \texttt{varpi\_LL} \\
\(\Gamma^{\rm LL}_p\) & the spin connection, the \(8\times8\) matrix
\(\Gamma^{\rm LL}_p=\mathsf s_{\rm LL}(\mathbf u^\omega_{{\rm LL},p})\otimes\sigma^{\rm LL}_0\)
(it acts on the Dirac index and leaves the isospin index alone) & new &
\texttt{Gamma\_LL} \\
\(\mathsf A_{{\rm LL},p}\) & the gauge matrix of the internal symmetries
on the eight components of \(\Psi_{\rm LL}\),
\(\mathsf A_{{\rm LL},p}=g^{\rm v}_{\rm LL}\sum_\iota A^\iota_{{\rm LL},p}T^\iota_{\rm LL}+g^{\rm F}_{\rm LL}A^{\rm F}_{{\rm LL},p}\mathsf 1_8\),
Hermitian (\(\mathsf 1_8\) because the fermion number of
\(\Psi_{\rm LL}\) is 1) & new & \texttt{Acal} \texttt{Acal\_0}
\texttt{Acal\_j} \texttt{Acal\_p} \\
\(\nabla^{\rm LL}_p\), \(\nabla^{{\rm LL}*}_p\) & the gauge-covariant
derivative of \(\Psi_{\rm LL}\), the \(8\times8\) matrix of differential
operators
\(\nabla^{\rm LL}_p=\partial_p+\Gamma^{\rm LL}_p-i\,\mathsf A_{{\rm LL},p}\),
and the covariant derivative of a field that transforms like
\(\Psi^*_{\rm LL}\) (the hole component),
\(\nabla^{{\rm LL}*}_p=\partial_p+(\Gamma^{\rm LL}_p)^*+i\,(\mathsf A_{{\rm LL},p})^*\)
& new & \texttt{nabla\_LL} \texttt{nabla*\_LL} \\
\(\mathbf D_{\rm LL}\) & the bi-quaternion of covariant derivatives,
\(\mathbf D_{\rm LL}=\sum_p\mathbf f_p\nabla^{\rm LL}_p\) (bold upright
like every bi-quaternion of this folder; it is not the Lorentz
derivation \(D_{\rm LL}(\mathbf u)\), which is italic and has an
argument), with
\(\sum_p\gamma^p_{\rm LL}\nabla^{\rm LL}_p=\begin{pmatrix}0&M_{\rm LL}(\mathbf D_{\rm LL})\\M_{\rm LL}(\bar{\mathbf D}_{\rm LL})&0\end{pmatrix}\)
& new & \texttt{Dbq\_LL} \\
\(\mathsf F^\iota_{{\rm LL},pp'}\),
\(\mathsf F^{\rm F}_{{\rm LL},pp'}\), \(\mathsf R^{\rm LL}_{pp'}\) & the
field strengths and the curvature:
\(\mathsf F^\iota_{{\rm LL},pp'}=\partial_pA^\iota_{{\rm LL},p'}-\partial_{p'}A^\iota_{{\rm LL},p}+g^{\rm v}_{\rm LL}\sum_{\iota',\iota''}\epsilon^{\rm LL}_{\iota\iota'\iota''}A^{\iota'}_{{\rm LL},p}A^{\iota''}_{{\rm LL},p'}\),
\(\mathsf F^{\rm F}_{{\rm LL},pp'}=\partial_pA^{\rm F}_{{\rm LL},p'}-\partial_{p'}A^{\rm F}_{{\rm LL},p}\),
and the \(8\times8\) curvature of the spin connection,
\(\mathsf R^{\rm LL}_{pp'}=\partial_p\Gamma^{\rm LL}_{p'}-\partial_{p'}\Gamma^{\rm LL}_p+\Gamma^{\rm LL}_p\Gamma^{\rm LL}_{p'}-\Gamma^{\rm LL}_{p'}\Gamma^{\rm LL}_p\);
the commutator of two covariant derivatives is
\(\mathsf R^{\rm LL}_{pp'}-i\,g^{\rm v}_{\rm LL}\sum_\iota\mathsf F^\iota_{{\rm LL},pp'}T^\iota_{\rm LL}-i\,g^{\rm F}_{\rm LL}\mathsf F^{\rm F}_{{\rm LL},pp'}\)
& new & \texttt{Fv\_LL} \texttt{FF\_LL} \texttt{Rc\_LL} \\
\(\Psi_{{\rm LL},1},\dots,\Psi_{{\rm LL},8}\); \(\Psi^\pm_{\rm LL}\);
\(\bar\Psi_{\rm LL}\) & the eight components of \(\Psi_{\rm LL}\):
components 1 and 2 are the first Dirac component with isospin up and
down, 3 and 4 the second, 5 and 6 the third, 7 and 8 the fourth;
\(\Psi^-_{\rm LL}\) (components 1 to 4, chirality \(-1\)) and
\(\Psi^+_{\rm LL}\) (components 5 to 8, chirality \(+1\)) are the two
Weyl isodoublets;
\(\bar\Psi_{\rm LL}=\Psi^\dagger_{\rm LL}(\gamma^0_{\rm LL}\otimes\sigma^{\rm LL}_0)\)
is the Dirac adjoint & new & \texttt{Psi\_LL\^{}+}
\texttt{Psi\_LL\^{}-} \\
\(\psi^{\rm p}_{\rm LL}\), \(\psi^{\rm h}_{\rm LL}\) & the particle and
the hole component (eight complex functions each, in the order of the
components of \(\Psi_{\rm LL}\)) of a mode of the mean-field equations:
\(\psi^{\rm p}_{\rm LL}\) transforms like \(\Psi_{\rm LL}\) and
\(\psi^{\rm h}_{\rm LL}\) like \(\Psi^*_{\rm LL}\); a component is
written \(\psi^{\rm p}_{{\rm LL},k}\), and a numbered mode
\((\psi^{\rm p}_{{\rm LL},n},\psi^{\rm h}_{{\rm LL},n})\) & new &
\texttt{psi\_p} \texttt{psi\_h} \\
\(\mathcal H^{\rm DB}_{\rm LL}\), \(h^\nabla_{\rm LL}\) & the
Dirac-Bogoliubov-de Gennes operator of the mean-field equations: the
\(2\times2\) block matrix (blocks of \(8\times8\)) with the upper left
block
\(h^\nabla_{\rm LL}-i\Gamma^{\rm LL}_0-\mathsf A_{{\rm LL},0}-\mu_{\rm LL}\),
the upper right block \(\mathsf P_{\rm LL}\), the lower left block
\(\mathsf P_{\rm LL}^\dagger\) and the lower right block
\(-(h^\nabla_{\rm LL}-i\Gamma^{\rm LL}_0-\mathsf A_{{\rm LL},0}-\mu_{\rm LL})^*\),
with the covariant Dirac Hamiltonian
\(h^\nabla_{\rm LL}=-i\sum_j\alpha^j_{\rm LL}\nabla^{\rm LL}_j+\beta_{\rm LL}m_{\rm LL}\)
(\(\alpha^j_{\rm LL}\) and \(\beta_{\rm LL}\) acting on the Dirac
index); on a plane wave \(\exp(i\sum_jq^j_{\rm LL}x^j_{\rm LL})\) with
constant fields it is a \(16\times16\) matrix
\(\mathcal H^{\rm DB}_{\rm LL}(\vec q_{\rm LL})\) (plain text
\texttt{H\_DB(q)}), which without gauge fields equals
\(\mathsf H_{\rm LL}(\vec q_{\rm LL})\) & new & \texttt{H\_DB}
\texttt{h\_nabla} \\
\(\hat a^{\rm qp}_{{\rm LL},n}\) & the annihilation operator of the
quasi-particle of a numbered mode of the mean-field equations (constant
in time):
\(\Psi_{\rm LL}=\sum_n(\psi^{\rm p}_{{\rm LL},n}\hat a^{\rm qp}_{{\rm LL},n}+(\psi^{\rm h}_{{\rm LL},n})^*(\hat a^{\rm qp}_{{\rm LL},n})^\dagger)\);
the quasi-particle vacuum is annihilated by every
\(\hat a^{\rm qp}_{{\rm LL},n}\) & new & \texttt{a\_qp} \\
\(E^{\rm qp}_{{\rm LL},n}\) & the energy of a numbered mode of the
mean-field equations in a stationary state, an eigenvalue of
\(\mathcal H^{\rm DB}_{\rm LL}\); the mean-field ground state is the
vacuum of the quasi-particles of the modes with
\(E^{\rm qp}_{{\rm LL},n}>0\) & new & \texttt{E\_qp} \\
\(\Pi^\iota_{\rm LL}(x_{\rm LL})\) & the local pair operator
\(\Pi^\iota_{\rm LL}(x_{\rm LL})=\Psi_{\rm LL}(x_{\rm LL})^T\mathsf K^\iota_{\rm LL}\Psi_{\rm LL}(x_{\rm LL})\),
which annihilates two spinor charges at the point \(x_{\rm LL}\); its
integral over space, with the plane waves normalized in a box, is the
operator \(\Pi^\iota_{\rm LL}\) of the momenta & new &
\texttt{Pi\_iota(x)} \\
\(\mathcal J^{{\rm F},p}_{\rm LL}\), \(\mathcal J^{\iota,p}_{\rm LL}\) &
the currents that are the sources of the gauge fields,
\(\mathcal J^{{\rm F},p}_{\rm LL}=g^{\rm F}_{\rm LL}\bar\Psi_{\rm LL}\gamma^p_{\rm LL}\Psi_{\rm LL}\)
and
\(\mathcal J^{\iota,p}_{\rm LL}=g^{\rm v}_{\rm LL}\bar\Psi_{\rm LL}\gamma^p_{\rm LL}T^\iota_{\rm LL}\Psi_{\rm LL}\)
(\(\gamma^p_{\rm LL}\) acting on the Dirac index) & new &
\texttt{JF\_LL} \texttt{Jv\_LL} \\
\(\mathcal L^\nabla_{\rm LL}\) & the gauge-invariant Lagrangian density
of Test 2: the contact form \(\mathcal L^{\rm c}_{\rm LL}\) with the
covariant derivative \(\nabla^{\rm LL}_p\) in place of \(\partial_p\),
plus the Lagrangian density of the gauge fields & new &
\texttt{Lagr\_nabla} \\
\(\ell^{\rm x}_{\rm LL}\), \(\ell^{\rm t}_{\rm LL}\),
\(r^{\rm W}_{\rm LL}\) & the lattice spacing (the same in the three
directions), a time step of the lattice equations (a fixed step of the
first version of Test 2, or the interval between the output times of an
integration by CVODE, which chooses its own internal steps), and the
coefficient of the Wilson term of the lattice equations & new &
\texttt{dx} \texttt{dt} \texttt{r\_W} \\
\(n^{\rm x}_{\rm LL}\), \(n^{\rm y}_{\rm LL}\), \(n^{\rm z}_{\rm LL}\);
\(\vec s_{\rm LL}\); \(\vec\delta^j_{\rm LL}\) & the numbers of lattice
sites in the directions \(x^1_{\rm LL}\), \(x^2_{\rm LL}\),
\(x^3_{\rm LL}\); a lattice site
\(\vec s_{\rm LL}=(s^1_{\rm LL},s^2_{\rm LL},s^3_{\rm LL})\), whose
integer coordinates \(s^j_{\rm LL}\) run from 0 to the number of sites
of their direction minus 1, at the point
\(x^j_{\rm LL}=s^j_{\rm LL}\ell^{\rm x}_{\rm LL}\); and the step to the
next site in the direction \(j\), the unit vector
\(\vec\delta^j_{\rm LL}\) & new & \texttt{nx} \texttt{ny} \texttt{nz}
\texttt{site} \texttt{delta\_j} \\
\(\mathsf U^{\rm link}_{{\rm LL},j}(\vec s_{\rm LL})\) & the link matrix
of the lattice from the site \(\vec s_{\rm LL}\) to the site
\(\vec s_{\rm LL}+\vec\delta^j_{\rm LL}\), the \(2\times2\) isospin
matrix
\(\exp(-i\ell^{\rm x}_{\rm LL}(g^{\rm v}_{\rm LL}\mathsf a_{{\rm LL},j}+g^{\rm F}_{\rm LL}A^{\rm F}_{{\rm LL},j}\sigma^{\rm LL}_0))\)
of the potentials at the site: the parallel transport of the isospin and
of the phase over one lattice step & new & \texttt{Ulink}
\texttt{Ulink\_j} \\
\(\mathcal H^{\rm lat}_{\rm LL}\), \(\varepsilon^{\rm lat}_{\rm LL}\),
\(B^{\rm H}_{\rm LL}\) & the lattice form of
\(\mathcal H^{\rm DB}_{\rm LL}\) (a matrix of 16 rows per site), the
energy of a free particle on the lattice, and the bound on the largest
eigenvalue magnitude of \(\mathcal H^{\rm lat}_{\rm LL}\) from which the
time step is chosen & new & \texttt{H\_lat} \texttt{eps\_lat}
\texttt{Hbound} \\
\(\mathsf G^{\rm site}_{\rm LL}\), \(\mathsf X^{\rm ph}_{\rm LL}\) & the
gauge transformation of the lattice equations, which acts at each site
as \(\exp(i\theta^{\rm F}_{\rm LL})\mathsf U_{\rm LL}\) on the isospin
index of the particle component and as its complex conjugate on the hole
component; and the exchange of the particle and the hole component,
\(\mathsf X^{\rm ph}_{\rm LL}=\begin{pmatrix}0&\mathsf 1_8\\\mathsf 1_8&0\end{pmatrix}\),
with which the particle-hole symmetry reads
\(\mathsf X^{\rm ph}_{\rm LL}(\mathcal H^{\rm DB}_{\rm LL})^*\mathsf X^{\rm ph}_{\rm LL}=-\mathcal H^{\rm DB}_{\rm LL}\)
& new & \texttt{Gsite} \texttt{Xph} \\
\(\mathsf A^\uparrow_{{\rm LL},p}\),
\(\mathsf A^\downarrow_{{\rm LL},p}\) & the two eigenvalues of
\(\mathsf A_{{\rm LL},p}\) when the isospin potentials point along the
third isospin direction (\(A^1_{{\rm LL},p}=A^2_{{\rm LL},p}=0\)):
\(\mathsf A^\uparrow_{{\rm LL},p}=g^{\rm F}_{\rm LL}A^{\rm F}_{{\rm LL},p}+\tfrac12g^{\rm v}_{\rm LL}A^3_{{\rm LL},p}\)
on the isospin-up components (\(Q_{\rm LL}=+\tfrac12\)) and
\(\mathsf A^\downarrow_{{\rm LL},p}=g^{\rm F}_{\rm LL}A^{\rm F}_{{\rm LL},p}-\tfrac12g^{\rm v}_{\rm LL}A^3_{{\rm LL},p}\)
on the isospin-down components (\(Q_{\rm LL}=-\tfrac12\)); constant
potentials of a space direction shift the momentum of a component, those
of the time direction its energy & new & \texttt{Acal\_up}
\texttt{Acal\_down} \texttt{Acal\_up\_0} \texttt{Acal\_up\_1}
\texttt{Acal\_down\_0} \texttt{Acal\_down\_1} \\
\(\omega^{\rm pr}_{\rm LL}\) & the angular frequency with which the
isospin of a mode turns in a constant isospin potential of the time
direction,
\(\omega^{\rm pr}_{\rm LL}=g^{\rm v}_{\rm LL}\big(\sum_\iota(A^\iota_{{\rm LL},0})^2\big)^{1/2}\);
the axis of the rotation is the direction of the potential in isospin
space & new & \texttt{omega\_pr} \\
\(P^\uparrow_{\rm LL}\), \(P^{\rm h}_{\rm LL}\) & the probability of
isospin up of a mode, the fraction of its norm
\(\int d^3\vec x_{\rm LL}(\lvert\psi^{\rm p}_{\rm LL}\rvert^2+\lvert\psi^{\rm h}_{\rm LL}\rvert^2)\)
that lies in the isospin-up components, and the probability of the hole,
the fraction that lies in \(\psi^{\rm h}_{\rm LL}\); both are functions
of the time \(x^0_{\rm LL}\) (in the data files the suffix
\texttt{\_cvode} marks the values of the integration of the lattice
equations by CVODE, \texttt{\_rk4} those of the integration by the
fixed-step Runge-Kutta method of the first version of Test 2, kept in
the folder \texttt{test2/rk4\_reference} for the comparison, and
\texttt{\_exact} those of the exact solution) & new & \texttt{P\_up}
\texttt{P\_up\_cvode} \texttt{P\_up\_rk4} \texttt{P\_up\_exact}
\texttt{P\_hole} \texttt{P\_hole\_cvode} \texttt{P\_hole\_rk4}
\texttt{P\_hole\_exact} \\
\(\chi^{\rm pw}_{\rm LL}\), \(\vartheta^{\rm B}_{\rm LL}\) & the
constant 8-component spinor (Dirac index and isospin) of a plane-wave
solution, a normalized eigenvector of
\(\mathsf h_{\rm LL}(\vec q_{\rm LL})\) or of its lattice form; and the
Bogoliubov angle of a plane-wave mode of the uniform polar superfluid,
\(\cos^2\vartheta^{\rm B}_{\rm LL}=\tfrac12\big(1+(\pm\varepsilon_{\rm LL}-\mu_{\rm LL})/E^\mp_{\rm LL}\big)\)
for an eigenvector of the energy \(\pm\varepsilon_{\rm LL}\), which
divides the mode into its particle part
\(\cos\vartheta^{\rm B}_{\rm LL}\,\chi^{\rm pw}_{\rm LL}\) and its hole
part
\(\sin\vartheta^{\rm B}_{\rm LL}\,\mathsf K^3_{\rm LL}\chi^{\rm pw}_{\rm LL}\)
& new & \texttt{chi\_pw} \texttt{theta\_B} \\
\(\Delta^\infty_{\rm LL}\), \(\ell^{\rm K}_{\rm LL}\) & the value of the
gap far from a domain wall (kink) of the pair field, and the width of
the wall: the kink
\(d^\iota_{\rm LL}=(0,0,\Delta_{\rm LL}(x^1_{\rm LL}))\) with
\(\Delta_{\rm LL}(x^1_{\rm LL})=\Delta^\infty_{\rm LL}\tanh(x^1_{\rm LL}/\ell^{\rm K}_{\rm LL})\)
changes the sign of the pair field at \(x^1_{\rm LL}=0\) & new &
\texttt{Delta\_inf} \texttt{ell\_K} \\
\(\varphi^{\rm K}_{\rm LL}\), \(\chi^{\rm K}_{\rm LL}\) & the profile
and the constant spinor of a bound state of a kink of the pair field:
\(\varphi^{\rm K}_{\rm LL}(x^1_{\rm LL})\) solves
\(d\varphi^{\rm K}_{\rm LL}/dx^1_{\rm LL}=-\Delta_{\rm LL}(x^1_{\rm LL})\varphi^{\rm K}_{\rm LL}\),
so that
\(\varphi^{\rm K}_{\rm LL}=\cosh(x^1_{\rm LL}/\ell^{\rm K}_{\rm LL})^{-\Delta^\infty_{\rm LL}\ell^{\rm K}_{\rm LL}}\)
for the kink of the width \(\ell^{\rm K}_{\rm LL}\);
\(\chi^{\rm K}_{\rm LL}\) is a constant 8-component spinor with
\(\beta_{\rm LL}\chi^{\rm K}_{\rm LL}=\pm\chi^{\rm K}_{\rm LL}\); the
bound state is
\((\psi^{\rm p}_{\rm LL},\psi^{\rm h}_{\rm LL})=(\varphi^{\rm K}_{\rm LL}\chi^{\rm K}_{\rm LL},\,i(\mathsf K^3_{\rm LL})^\dagger\alpha^1_{\rm LL}\varphi^{\rm K}_{\rm LL}\chi^{\rm K}_{\rm LL})\),
with the energy \(\pm m_{\rm LL}\) & new & \texttt{phi\_K}
\texttt{chi\_K} \\
\(\mathsf Q^{\rm K}_{\rm LL}\); \(k^{\rm K}_{\rm LL}\) & the constant
Hermitian \(16\times16\) matrix
\(\mathsf Q^{\rm K}_{\rm LL}=\begin{pmatrix}0&-i\alpha^1_{\rm LL}(\mathsf K^3_{\rm LL})^\dagger\\-i\alpha^1_{\rm LL}\mathsf K^3_{\rm LL}&0\end{pmatrix}\)
(blocks of \(8\times8\)), with \((\mathsf Q^{\rm K}_{\rm LL})^2\) the
unit matrix and the trace 0, through which the square of the
Dirac-Bogoliubov-de Gennes operator of a wall of the pair field depends
on the slope of the gap:
\((\mathcal H^{\rm DB}_{\rm LL})^2=\big(-d^2/d(x^1_{\rm LL})^2+(q^2_{\rm LL})^2+(q^3_{\rm LL})^2+m_{\rm LL}^2+\Delta_{\rm LL}^2\big)+(d\Delta_{\rm LL}/dx^1_{\rm LL})\,\mathsf Q^{\rm K}_{\rm LL}\);
and the number \(k^{\rm K}_{\rm LL}=0,1,2,\dots\) of a bound state of a
kink, below \(\Delta^\infty_{\rm LL}\ell^{\rm K}_{\rm LL}\) & new &
\texttt{Qmat\_K} \texttt{k\_K} \\
\(E^{\rm R}_{\rm LL}\) & the Rayleigh quotient of a trial mode on the
lattice: the expectation value of \(\mathcal H^{\rm lat}_{\rm LL}\) in
the mode divided by the norm of the mode; it approximates the eigenvalue
nearest to it with an error of second order in the residual of the mode
& new & \texttt{E\_R} \\
\(B^\uparrow_{\rm LL}\), \(B^\downarrow_{\rm LL}\);
\(n^{\rm L}_{\rm LL}\) & the magnetic fields that the isospin-up and the
isospin-down components feel in the constant field strengths
\(\mathsf F^{\rm F}_{{\rm LL},12}\) (of the phase group) and
\(\mathsf F^3_{{\rm LL},12}\) (of the third isospin direction),
\(B^\uparrow_{\rm LL}=g^{\rm F}_{\rm LL}\mathsf F^{\rm F}_{{\rm LL},12}+\tfrac12g^{\rm v}_{\rm LL}\mathsf F^3_{{\rm LL},12}\)
and
\(B^\downarrow_{\rm LL}=g^{\rm F}_{\rm LL}\mathsf F^{\rm F}_{{\rm LL},12}-\tfrac12g^{\rm v}_{\rm LL}\mathsf F^3_{{\rm LL},12}\)
(the fields point along \(x^3_{\rm LL}\)); and the number
\(n^{\rm L}_{\rm LL}=0,1,2,\dots\) of a Landau level, of the energies
\(\pm\big(m_{\rm LL}^2+(q^3_{\rm LL})^2+2n^{\rm L}_{\rm LL}\lvert B^\uparrow_{\rm LL}\rvert\big)^{1/2}\)
for isospin up and likewise with \(B^\downarrow_{\rm LL}\) for isospin
down & new & \texttt{B\_up} \texttt{B\_down} \texttt{n\_L} \\
\(n^{\rm act}_{\rm LL}\) & the number of active directions of a lattice
problem, the directions with more than one site (1, 2 or 3); the mode
functions of a problem with \(n^{\rm act}_{\rm LL}\) active directions
are normalized to
\(\sum_{\vec s_{\rm LL}}(\ell^{\rm x}_{\rm LL})^{n^{\rm act}_{\rm LL}}(\lvert\psi^{\rm p}_{\rm LL}\rvert^2+\lvert\psi^{\rm h}_{\rm LL}\rvert^2)=1\)
& new & \texttt{n\_act} \\
\(N^{\rm F}_{\rm LL}\), \(E^{\rm MF}_{\rm LL}\) & the fermion number of
a self-consistent lattice state, the sum over its modes and over the
sites of the squared hole components
\(\lvert\psi^{\rm h}_{\rm LL}\rvert^2\) (conserved by the
self-consistent evolution), and the mean-field energy functional, the
sum over the modes of the expectation values of
\(\mathcal H^{\rm lat}_{\rm LL}\) without the pair field plus
\(2(\ell^{\rm x}_{\rm LL})^{n^{\rm act}_{\rm LL}}/G_{\rm LL}\) times the
sum over the sites of \(\sum_\iota\lvert d^\iota_{\rm LL}\rvert^2\)
(conserved in time-independent gauge fields) & new & \texttt{N\_F}
\texttt{E\_MF} \\
\(\mathsf W^{\rm m}_{\rm LL}\) & the matrix whose columns are the modes
of a self-consistent lattice state (unit vectors with 16 components per
site); the self-consistent evolution moves all columns together & new &
\texttt{W\_modes} \\
\(\mathcal M^{\rm sc}_{\rm LL}\); \(y^{\rm A}_{\rm LL}\),
\(r^{\rm A}_{\rm LL}\), \(Y^{\rm A}_{\rm LL}\), \(R^{\rm A}_{\rm LL}\),
\(c^{\rm A}_{\rm LL}\), \(a^{\rm A}_{\rm LL}\); \(n^{\rm it}_{\rm LL}\)
& the self-consistency map, which takes a pair field to
\(-\tfrac12G_{\rm LL}\) times the pair expectation of the
positive-energy eigenvectors of \(\mathcal H^{\rm lat}_{\rm LL}\) with
that pair field, divided by
\((\ell^{\rm x}_{\rm LL})^{n^{\rm act}_{\rm LL}}\); the quantities of
the Anderson iteration for its fixed point: the iterate (the real and
imaginary parts of all \(d^\iota_{\rm LL}\) on all sites, as one real
vector), its residual (the map applied to the iterate, minus the
iterate), the matrices whose columns are the differences of successive
iterates and of successive residuals, the coefficients of the
least-squares combination, the damping factor, and the number of the
iteration & new & \texttt{M\_sc} \texttt{y\_A} \texttt{r\_A}
\texttt{Y\_A} \texttt{R\_A} \texttt{c\_A} \texttt{a\_A}
\texttt{n\_it} \\
\(\vec x^{\rm c}_{\rm LL}\), \(w^{\rm pk}_{\rm LL}\),
\(\vec q^{\rm pk}_{\rm LL}\), \(\chi^{\rm pk}_{\rm LL}\),
\(s^{\rm br}_{\rm LL}\), \(\mathsf Z^\pm_{\rm LL}(\vec q_{\rm LL})\) &
the Gaussian wave packet of the solvers: its centre, its width (the
standard deviation of its density), its central momentum, its
8-component spinor (the projection of a reference spinor of definite
spin and isospin on the chosen branch at the central momentum), the sign
\(\pm1\) of the chosen branch \(\pm\varepsilon^{\rm lat}_{\rm LL}\), and
the projectors
\(\tfrac12(\mathsf 1_4\pm\mathsf h^{\rm lat}_{\rm LL}(\vec q_{\rm LL})/\varepsilon^{\rm lat}_{\rm LL}(\vec q_{\rm LL}))\)
on the two branches of the free lattice operator at the momentum
\(\vec q_{\rm LL}\) & new & \texttt{x\_c} \texttt{w\_pk} \texttt{q\_pk}
\texttt{chi\_pk} \texttt{s\_br} \texttt{Zproj} \\
\(\mathsf h^{\rm lat}_{\rm LL}(\vec q_{\rm LL})\) & the \(4\times4\)
matrix of the free lattice operator (without isospin) on the plane wave
\(\exp(i\vec q_{\rm LL}\cdot\vec x_{\rm LL})\):
\(\sum_j\alpha^j_{\rm LL}\sin(q^j_{\rm LL}\ell^{\rm x}_{\rm LL})/\ell^{\rm x}_{\rm LL}+\beta_{\rm LL}\big(m_{\rm LL}+(r^{\rm W}_{\rm LL}/\ell^{\rm x}_{\rm LL})\sum_j(1-\cos(q^j_{\rm LL}\ell^{\rm x}_{\rm LL}))\big)\),
the sums over the active directions; its square is
\(\varepsilon^{\rm lat}_{\rm LL}(\vec q_{\rm LL})^2\) times
\(\mathsf 1_4\) & new & \texttt{h\_lat\_q} \\
\(\vec v^{\rm g}_{\rm LL}\), \(\bar v^{\rm g}_{\rm LL}\) & the group
velocity, the gradient of \(\varepsilon^{\rm lat}_{\rm LL}\) with
respect to \(\vec q_{\rm LL}\), and its mean over the distribution of
momenta of a wave packet & new & \texttt{v\_g} \texttt{v\_g\_mean} \\
\(\ell^{\rm B}_{\rm LL}\) & the magnetic length
\(\lvert B^\uparrow_{\rm LL}\rvert^{-1/2}\) of the isospin-up component
in a constant magnetic field & new & \texttt{l\_B} \\
\(G'_{\rm LL}\); \(\bar\Delta_{\rm LL}\); \(\omega^{\rm H}_{\rm LL}\) &
the coupling after a quench (the coupling is changed from \(G_{\rm LL}\)
to \(G'_{\rm LL}\) at \(x^0_{\rm LL}=0\)); the mean over the time of the
gap \(\lvert\vec d_{\rm LL}\rvert\) after the quench; and the angular
frequency of the amplitude oscillation of the pair field (the amplitude
mode, the analogue of the Higgs boson) & new & \texttt{G\_quench}
\texttt{Delta\_mean} \texttt{omega\_H} \\
\(u^{\rm Bl}_{\rm LL}\), \(v^{\rm Bl}_{\rm LL}\);
\(\vec b^{\rm Bl}_{\rm LL}\), \(\vec h^{\rm Bl}_{\rm LL}\);
\(\vec b^{\rm fl}_{\rm LL}\), \(\Delta^{\rm fl}_{\rm LL}\) & the two
complex coefficients of a state of the plane (3.12) (of the particle
vector and of the hole vector that span it); the Bloch vector of that
state, the expectation values of \(\sigma^{\rm LL}_1\),
\(\sigma^{\rm LL}_2\), \(\sigma^{\rm LL}_3\) in it, and the vector of
the coefficients of the \(2\times2\) matrix (3.12) in the Pauli
matrices, about which it precesses; their small deviations from the
ground state in the linearized equations: the deviation of the Bloch
vector and of the gap & new & \texttt{u\_Bl} \texttt{v\_Bl}
\texttt{b\_Bl} \texttt{h\_Bl} \texttt{b\_fl} \texttt{Delta\_fl} \\
\(\mathsf K^{\rm g}_{\rm LL}\), \(\mathsf K^{\rm g'}_{\rm LL}\) & two
general pair matrices: complex \(8\times8\) matrices with which
\(\Psi^T_{\rm LL}\mathsf K^{\rm g}_{\rm LL}\Psi_{\rm LL}=\sum_{k,k'}(\mathsf K^{\rm g}_{\rm LL})_{kk'}\Psi_{{\rm LL},k}\Psi_{{\rm LL},k'}\)
annihilates two spinor charges at one point; only the antisymmetric part
of a pair matrix contributes, because the components of
\(\Psi_{\rm LL}\) anticommute & new & \texttt{Kgen}
\texttt{Kgen\textquotesingle{}} \\
\(\mathsf V^\pm_{\rm LL}\); \(\mathsf B^{\rm W}_{\rm LL}\) & the isospin
factors of the two Weyl blocks of a Lorentz-invariant pair matrix, which
is
\(\mathrm{diag}(\mathsf E_{\rm LL}\otimes\mathsf V^-_{\rm LL},\mathsf E_{\rm LL}\otimes\mathsf V^+_{\rm LL})\)
with two symmetric complex \(2\times2\) matrices
\(\mathsf V^\pm_{\rm LL}\) on the isospin index; and a complex
\(2\times2\) matrix on the Weyl index, the Weyl factor of a term of a
\(4\times4\) block of a pair matrix (in the proof that the
Lorentz-invariant pairs have this form) & new & \texttt{Vblock}
\texttt{BW} \\
\(\mathsf K^{(+1)}_{\rm LL}\), \(\mathsf K^{(0)}_{\rm LL}\),
\(\mathsf K^{(-1)}_{\rm LL}\); \(\Pi^{(+1)}_{\rm LL}\),
\(\Pi^{(0)}_{\rm LL}\), \(\Pi^{(-1)}_{\rm LL}\) & the scalar pair
matrices of definite electric charge,
\(\mathsf K^{(\pm1)}_{\rm LL}=(\mathsf K^1_{\rm LL}\mp i\mathsf K^2_{\rm LL})/\sqrt2\)
and \(\mathsf K^{(0)}_{\rm LL}=\mathsf K^3_{\rm LL}\), and the pair
operators \(\Pi^{(+1)}_{\rm LL}\), \(\Pi^{(0)}_{\rm LL}\),
\(\Pi^{(-1)}_{\rm LL}\) made with them, which annihilate a pair of the
electric charge \(+1\) (two isospin-up charges), \(0\) (one isospin-up
and one isospin-down charge) and \(-1\) (two isospin-down charges) & new
& \texttt{K(+1)} \texttt{K(0)} \texttt{K(-1)} \texttt{Pi(+1)}
\texttt{Pi(0)} \texttt{Pi(-1)} \\
\(S^{\rm sp}_{\rm LL}\), \(I^{\rm v}_{\rm LL}\) & the spin and the
isospin of a state of two spinor charges: the sums of the squares of the
three total spin generators and of the three total isospin generators
have the eigenvalues \(S^{\rm sp}_{\rm LL}(S^{\rm sp}_{\rm LL}+1)\) and
\(I^{\rm v}_{\rm LL}(I^{\rm v}_{\rm LL}+1)\); the superfluid boson
singlet has \(S^{\rm sp}_{\rm LL}=0\) (spin singlet) and
\(I^{\rm v}_{\rm LL}=1\) (isospin triplet) & new & \texttt{S\_sp}
\texttt{I\_v} \\
\(\hat c_{{\rm LL},k}(\vec s_{\rm LL})\);
\(\hat\Pi^\iota_{\rm LL}(\vec s_{\rm LL})\),
\(\hat n^{\rm F}_{\rm LL}(\vec s_{\rm LL})\),
\(\hat n^\iota_{\rm LL}(\vec s_{\rm LL})\) & operators of one lattice
site \(\vec s_{\rm LL}\) (a cell of the volume
\((\ell^{\rm x}_{\rm LL})^3\)): the annihilation operators
\(\hat c_{{\rm LL},k}(\vec s_{\rm LL})=(\ell^{\rm x}_{\rm LL})^{3/2}\Psi_{{\rm LL},k}(\vec s_{\rm LL})\)
of the eight components, which obey the canonical anticommutation
relations; the pair operator of the site
\(\hat\Pi^\iota_{\rm LL}=\sum_{k,k'}(\mathsf K^\iota_{\rm LL})_{kk'}\hat c_{{\rm LL},k}\hat c_{{\rm LL},k'}\);
the number of spinor charges on the site
\(\hat n^{\rm F}_{\rm LL}=\sum_k\hat c^\dagger_{{\rm LL},k}\hat c_{{\rm LL},k}\);
and the isospin of the site
\(\hat n^\iota_{\rm LL}=\sum_{k,k'}\hat c^\dagger_{{\rm LL},k}(T^\iota_{\rm LL})_{kk'}\hat c_{{\rm LL},k'}\)
& new & \texttt{c\_site} \texttt{Pi\_site} \texttt{n\_F\_site}
\texttt{n\_iota\_site} \\
\(\hat H^{\rm q}_{\rm LL}\); \(\lvert\Phi^{\rm SF}_{\rm LL}\rangle\),
\(\lvert\Phi^{\rm N}_{\rm LL}\rangle\) & the part of the mean-field
Hamiltonian of the uniform polar state that belongs to one pair of
momenta \(\vec q_{\rm LL}\), \(-\vec q_{\rm LL}\) (an operator on 16
fermion modes); the superfluid boson singlet state, its ground state, a
coherent state of pairs; and the normal state from which it is built,
the Dirac sea of these modes (every state of energy
\(-\varepsilon_{\rm LL}\) filled, every state of energy
\(+\varepsilon_{\rm LL}\) empty) & new & \texttt{H\_q}
\texttt{SF\_state} \texttt{Dirac\_sea} \\
\(\hat B^{\rm p}_{\rm LL}(\vec q_{\rm LL})\),
\(\hat B^{\rm a}_{\rm LL}(\vec q_{\rm LL})\); \(t^{\rm p}_{\rm LL}\),
\(t^{\rm a}_{\rm LL}\) & the pair operators of the momenta
\(\vec q_{\rm LL}\) and \(-\vec q_{\rm LL}\) in the scalar isospin
channel of \(\mathsf K^3_{\rm LL}\):
\(\hat B^{\rm p}_{\rm LL}(\vec q_{\rm LL})^\dagger\) creates a pair of
spinor charges of positive energy (particles),
\(\hat B^{\rm a}_{\rm LL}(\vec q_{\rm LL})^\dagger\) removes a pair from
the states of negative energy (creates a pair of antiparticles); and
their weights
\(t^{\rm p}_{\rm LL}=\Delta_{\rm LL}/(E^-_{\rm LL}+\varepsilon_{\rm LL}-\mu_{\rm LL})\)
and
\(t^{\rm a}_{\rm LL}=\Delta_{\rm LL}/(E^+_{\rm LL}+\varepsilon_{\rm LL}+\mu_{\rm LL})\)
in the superfluid boson singlet state & new & \texttt{B\_p}
\texttt{B\_a} \texttt{t\_p} \texttt{t\_a} \\
\(\mathsf Y^\pm_{\rm LL}(\vec q_{\rm LL})\); \(\hat f_{\rm LL}\),
\(\hat f'_{\rm LL}\) & the projectors
\(\tfrac12(\mathsf 1_8\pm\mathsf h_{\rm LL}(\vec q_{\rm LL})/\varepsilon_{\rm LL})\)
on the states of energy \(\pm\varepsilon_{\rm LL}\) of the plane-wave
Hamiltonian (rank 4 each); and the annihilation operators of a spinor
charge in an eigenstate of \(\mathsf h_{\rm LL}(\vec q_{\rm LL})\) and
in its partner state of the momentum \(-\vec q_{\rm LL}\) (the complex
conjugate of \(\mathsf K^3_{\rm LL}\) times the eigenvector) & new &
\texttt{Yproj} \texttt{f\_state} \texttt{f\_partner} \\
\(\mathcal L^{\rm HS}_{\rm LL}\) & the Lagrangian density of Test 2
after the Hubbard-Stratonovich identity: the Lagrangian
\(\mathcal L^\nabla_{\rm LL}\) with its pairing term replaced by
\(-\sum_\iota\lvert d^\iota_{\rm LL}\rvert^2/G_{\rm LL}-\tfrac12\sum_\iota\big((d^\iota_{\rm LL})^*\Pi^\iota_{\rm LL}+d^\iota_{\rm LL}(\Pi^\iota_{\rm LL})^\dagger\big)\),
in which \(\vec d_{\rm LL}\) is an independent complex isovector field &
new & \texttt{Lagr\_HS} \\
\(\nabla^{\rm B}_{{\rm LL},p}\) & the gauge-covariant derivative of the
SFBS field,
\(\nabla^{\rm B}_{{\rm LL},p}\vec d_{\rm LL}=\partial_p\vec d_{\rm LL}-2i\,g^{\rm F}_{\rm LL}A^{\rm F}_{{\rm LL},p}\vec d_{\rm LL}+g^{\rm v}_{\rm LL}\mathcal A^{\rm v}_{{\rm LL},p}\vec d_{\rm LL}\)
(fermion number 2, isospin vector, Lorentz scalar) & new &
\texttt{nablaB} \texttt{nablaB\_p} \\
\(\mathcal A^{\rm v}_{{\rm LL},p}\), \(\vec A^{\rm v}_{{\rm LL},p}\) &
the isospin potential of the direction \(p\) in the vector
representation: the vector
\(\vec A^{\rm v}_{{\rm LL},p}=(A^1_{{\rm LL},p},A^2_{{\rm LL},p},A^3_{{\rm LL},p})\)
and the real antisymmetric \(3\times3\) matrix with the entries
\((\mathcal A^{\rm v}_{{\rm LL},p})_{\iota\iota''}=\sum_{\iota'}\epsilon^{\rm LL}_{\iota\iota'\iota''}A^{\iota'}_{{\rm LL},p}\),
so that
\(\mathcal A^{\rm v}_{{\rm LL},p}\vec d_{\rm LL}=\vec A^{\rm v}_{{\rm LL},p}\times\vec d_{\rm LL}\)
& new & \texttt{Avec\_matrix} \texttt{Avec\_p} \\
\(\Psi^{\rm a}_{\rm LL}\), \(\Psi^{\rm b}_{\rm LL}\) & two arbitrary
eight-component fields that transform like \(\Psi_{\rm LL}\); used in
the Leibniz rule of the covariant derivative of the SFBS & new &
\texttt{Psi\_a} \texttt{Psi\_b} \\
\(\mathcal S^{\rm B}_{\rm LL}\) & the effective action of the SFBS:
\(\exp(i\mathcal S^{\rm B}_{\rm LL})\) is the functional integral over
the spinor field of
\(\exp(i\int d^4x_{\rm LL}\,\mathcal L^{\rm HS}_{\rm LL})\) at given
\(\vec d_{\rm LL}\) and given gauge fields & new & \texttt{S\_B} \\
\(\mathcal L^{\rm B}_{\rm LL}\); \(Z^{\rm t}_{\rm LL}\),
\(Z^{\rm s}_{\rm LL}\) & the effective Lagrangian density of the SFBS to
second order in the derivatives,
\(Z^{\rm t}_{\rm LL}\sum_\iota\lvert(\nabla^{\rm B}_{{\rm LL},0}\vec d_{\rm LL})^\iota\rvert^2-Z^{\rm s}_{\rm LL}\sum_j\sum_\iota\lvert(\nabla^{\rm B}_{{\rm LL},j}\vec d_{\rm LL})^\iota\rvert^2-\Omega_{\rm LL}(\vec d_{\rm LL})\),
and its stiffness coefficients of the time derivatives and of the space
derivatives; they depend on the vacuum about which the expansion is
made, written as the argument: \(Z^{\rm t}_{\rm LL}(0)\),
\(Z^{\rm s}_{\rm LL}(0)\) in the normal vacuum (plain text
\texttt{Z\_t0}, \texttt{Z\_s0}) & new & \texttt{Lagr\_B} \texttt{Z\_t}
\texttt{Z\_s} \texttt{Z\_t0} \texttt{Z\_s0} \\
\(\chi^{\rm pr}_{\rm LL}(\omega^{\rm B}_{\rm LL},k^{\rm B}_{\rm LL})\) &
the pair susceptibility of the vacuum: the coefficient of
\(\lvert\vec d_{\rm LL}\rvert^2\) that the spinor field contributes to
the effective action, per duration and volume, of a plane wave of the
SFBS field with the frequency \(\omega^{\rm B}_{\rm LL}\) and the wave
number \(k^{\rm B}_{\rm LL}\) (plain text \texttt{chi\_pr\_GL}: its
derivative expansion to second order) & new & \texttt{chi\_pr}
\texttt{chi\_pr\_GL} \\
\(\omega^{\rm B}_{\rm LL}\), \(\vec k^{\rm B}_{\rm LL}\),
\(k^{\rm B}_{\rm LL}\) & the frequency, the wave vector and the wave
number of a plane wave of the SFBS field (plain text
\texttt{omega\_B\_KG}: the frequency of the Klein-Gordon form of the
effective field equation) & new & \texttt{omega\_B} \texttt{k\_B}
\texttt{omega\_B\_KG} \\
\(M^{\rm pr}_{\rm LL}(\vec q_{\rm LL},\vec k^{\rm B}_{\rm LL})\);
\(\varepsilon_{\rm LL}(+)\), \(\varepsilon_{\rm LL}(-)\) & the pair
matrix element
\({\rm tr}\big(\mathsf Y^+_{\rm LL}(\vec q_{\rm LL}+\vec k^{\rm B}_{\rm LL}/2)\mathsf Y^+_{\rm LL}(\vec q_{\rm LL}-\vec k^{\rm B}_{\rm LL}/2)\big)\)
of two spinor charges with the relative momentum \(\vec q_{\rm LL}\) and
the total momentum \(\vec k^{\rm B}_{\rm LL}\), and the energies
\(\varepsilon_{\rm LL}(\pm)=(\lvert\vec q_{\rm LL}\pm\vec k^{\rm B}_{\rm LL}/2\rvert^2+m_{\rm LL}^2)^{1/2}\)
of the two charges & new & \texttt{M\_pr} \texttt{eps(+)}
\texttt{eps(-)} \\
\(\mathsf H^{\rm mod}_{\rm LL}(\vec q_{\rm LL})\) & the \(16\times16\)
Bogoliubov matrix of the static pair field
\((0,0,d^3_{\rm LL}\exp(i\vec k^{\rm B}_{\rm LL}\cdot\vec x_{\rm LL}))\)
for the relative momentum \(\vec q_{\rm LL}\), acting on the Nambu pair
\((\Psi_{\rm LL}(\vec q_{\rm LL}+\vec k^{\rm B}_{\rm LL}/2),\Psi_{\rm LL}(\vec k^{\rm B}_{\rm LL}/2-\vec q_{\rm LL})^*)\)
& new & \texttt{H\_mod} \\
\(G^{\rm bd}_{\rm LL}\) & the binding coupling,
\(1/G^{\rm bd}_{\rm LL}=\chi^{\rm pr}_{\rm LL}(2m_{\rm LL},0)\): for
\(G^{\rm bd}_{\rm LL}<G_{\rm LL}<G^{\rm crit}_{\rm LL}\) the SFBS of the
normal vacuum is a bound state of two spinor charges & new &
\texttt{G\_bind} \\
\(M^{\rm B}_{\rm LL}\) & the mass of the SFBS boson of the normal
vacuum: the frequency \(\omega^{\rm B}_{\rm LL}\) below the threshold
\(2m_{\rm LL}\) with
\(\chi^{\rm pr}_{\rm LL}(\omega^{\rm B}_{\rm LL},0)=1/G_{\rm LL}\)
(plain text \texttt{M\_B\_GL}: the Ginzburg-Landau mass
\((a^{\rm GL}_{\rm LL}/Z^{\rm t}_{\rm LL}(0))^{1/2}\);
\texttt{line\_MB}: a fitted oscillation with this frequency) & new &
\texttt{M\_B} \texttt{M\_B\_GL} \texttt{line\_MB} \\
\(a^{\rm GL}_{\rm LL}\) & the Ginzburg-Landau coefficient of
\(\lvert\vec d_{\rm LL}\rvert^2\) in the effective potential of the
normal vacuum, \(1/G_{\rm LL}-1/G^{\rm crit}_{\rm LL}\) & new &
\texttt{a\_GL} \\
\(\vec P^{\rm B}_{\rm LL}\) & the canonical momentum of the SFBS field,
\(\vec P^{\rm B}_{\rm LL}=Z^{\rm t}_{\rm LL}\nabla^{\rm B}_{{\rm LL},0}\vec d_{\rm LL}\),
the derivative of \(\mathcal L^{\rm B}_{\rm LL}\) with respect to
\(\partial_0\vec d^{\,*}_{\rm LL}\) & new & \texttt{P\_B} \\
\(\mathcal J^{\iota,p}_{{\rm B},{\rm LL}}\),
\(\mathcal J^{{\rm F},p}_{{\rm B},{\rm LL}}\) & the currents of the
SFBS: the derivatives of \(\mathcal L^{\rm B}_{\rm LL}\) with respect to
\(A^\iota_{{\rm LL},p}\) and \(A^{\rm F}_{{\rm LL},p}\), the sources of
the gauge fields in the effective theory & new & \texttt{Jv\_B}
\texttt{JF\_B} \\
\(u^{\rm pl}_{{\rm LL},\pm}\), \(v^{\rm pl}_{{\rm LL},\pm}\) & the two
components of a mode of the uniform polar mean-field dynamics in the
plane spanned by \((\chi^{\rm pw}_{\rm LL},0)\) and
\((0,\mathsf K^3_{\rm LL}\chi^{\rm pw}_{\rm LL})\), for an eigenvector
\(\chi^{\rm pw}_{\rm LL}\) of the energy \(\pm\varepsilon_{\rm LL}\);
functions of the time, with
\(\lvert u^{\rm pl}_{{\rm LL},\pm}\rvert^2+\lvert v^{\rm pl}_{{\rm LL},\pm}\rvert^2=1\)
& new & \texttt{u\_pl} \texttt{v\_pl} \\
\(W^{\rm g}_{\rm LL}(\Delta_{\rm LL}^2)\) & the condensation function of
one sector,
\(W^{\rm g}_{\rm LL}(\Delta_{\rm LL}^2)=\int\!d^3q_{\rm LL}(2\pi)^{-3}\big(E^+_{\rm LL}+E^-_{\rm LL}-\lvert\varepsilon_{\rm LL}+\mu_{\rm LL}\rvert-\lvert\varepsilon_{\rm LL}-\mu_{\rm LL}\rvert\big)\)
with the gap \(\Delta_{\rm LL}\) in \(E^\pm_{\rm LL}\); increasing and
strictly concave in \(\Delta_{\rm LL}^2\); the effective potential is
\(\lvert\vec d_{\rm LL}\rvert^2/G_{\rm LL}-W^{\rm g}_{\rm LL}((\Delta^+_{\rm LL})^2)-W^{\rm g}_{\rm LL}((\Delta^-_{\rm LL})^2)\)
& new & \texttt{W\_g} \\
\(\Xi_{\rm LL}((M^{\rm q}_{\rm LL})^2)\); \(\Xi^{(1)}_{\rm LL}\),
\(\Xi^{(2)}_{\rm LL}\), \(\Xi^{(3)}_{\rm LL}\) & the function
\(\Xi_{\rm LL}((M^{\rm q}_{\rm LL})^2)=2\int_{q_{\rm LL}<\Lambda_{\rm LL}}d^3q_{\rm LL}(2\pi)^{-3}(q_{\rm LL}^2+(M^{\rm q}_{\rm LL})^2)^{1/2}\),
which gives the condensation function at \(\mu_{\rm LL}=0\) in closed
form, and its first, second and third derivatives with respect to its
argument & new & \texttt{Xi}
\texttt{Xi\textquotesingle{}\textquotesingle{}} \texttt{Xi\_1}
\texttt{Xi\_2} \texttt{Xi\_3} \\
\(M^{\rm q}_{\rm LL}\) & the quasi-particle mass of the condensed vacuum
at \(\mu_{\rm LL}=0\), \((m_{\rm LL}^2+\Delta_{\rm LL}^2)^{1/2}\): the
smallest quasi-particle energy, at zero momentum & new &
\texttt{M\_q} \\
\(\Omega^{\rm GL}_{\rm LL}(\vec d_{\rm LL})\); \(b^{\rm GL}_{\rm LL}\) &
the Ginzburg-Landau form of the effective potential,
\(a^{\rm GL}_{\rm LL}\lvert\vec d_{\rm LL}\rvert^2+b^{\rm GL}_{\rm LL}(2\lvert\vec d_{\rm LL}\rvert^4-\lvert\vec d_{\rm LL}\cdot\vec d_{\rm LL}\rvert^2)\),
and its quartic coefficient, at \(\mu_{\rm LL}=0\) equal to
\(Z^{\rm t}_{\rm LL}(0)\) & new & \texttt{Omega\_GL} \texttt{b\_GL} \\
\(\Delta^{\rm vac}_{\rm LL}\) & the gap of the condensed vacuum: the
radius of the circle of minima of the effective potential at
\(\mu_{\rm LL}=0\) for \(G_{\rm LL}>G^{\rm crit}_{\rm LL}\), the
solution of the gap equation (plain text \texttt{Delta\_vac\_GL}: its
Ginzburg-Landau value
\((-a^{\rm GL}_{\rm LL}/(2b^{\rm GL}_{\rm LL}))^{1/2}\)) & new &
\texttt{Delta\_vac} \texttt{Delta\_vac\_GL} \\
\(\Omega^{\rm s}_{\rm LL}(\Delta_{\rm LL})\);
\(\Omega^{\rm s(2)}_{\rm LL}\), \(\Omega^{\rm s(3)}_{\rm LL}\),
\(\Omega^{\rm s(4)}_{\rm LL}\) & the potential of one sector,
\(\Omega^{\rm s}_{\rm LL}(\Delta_{\rm LL})=\Delta_{\rm LL}^2/(2G_{\rm LL})-W^{\rm g}_{\rm LL}(\Delta_{\rm LL}^2)\),
an even function of its real argument, with
\(\Omega_{\rm LL}=\Omega^{\rm s}_{\rm LL}(\Delta^+_{\rm LL})+\Omega^{\rm s}_{\rm LL}(\Delta^-_{\rm LL})\);
and its derivatives of the orders 2, 3 and 4 at
\(\Delta^{\rm vac}_{\rm LL}\), the coefficients of the expansion of the
potential about the vacuum (plain text:
\texttt{Omega\_s\textquotesingle{}} and \texttt{Omega\_s\^{}(1)} the
first derivative, \texttt{Omega\_s2} to \texttt{Omega\_s4} the parameter
names of the data files) & new & \texttt{Omega\_s}
\texttt{Omega\_s\textquotesingle{}} \texttt{Omega\_s\^{}(1)}
\texttt{Omega\_s\^{}(2)} \texttt{Omega\_s\^{}(3)}
\texttt{Omega\_s\^{}(4)} \texttt{Omega\_s2} \texttt{Omega\_s3}
\texttt{Omega\_s4} \\
\(\Omega^{\rm min}_{\rm LL}\) & the smallest value of the effective
potential, \(2\,\Omega^{\rm s}_{\rm LL}(\Delta^{\rm vac}_{\rm LL})\),
taken on the vacuum manifold (plain text \texttt{Omega\_min}, listed
with \(\Omega_{\rm LL}\)) & new & \\
\(\Xi^{(4)}_{\rm LL}\) & the fourth derivative of \(\Xi_{\rm LL}\) with
respect to its argument \((M^{\rm q}_{\rm LL})^2\),
\(-\tfrac{15}8\int_{q_{\rm LL}<\Lambda_{\rm LL}}d^3q_{\rm LL}(2\pi)^{-3}(q_{\rm LL}^2+(M^{\rm q}_{\rm LL})^2)^{-7/2}\)
& new & \texttt{Xi\_4} \\
\(y^+_{\rm LL}\), \(y^-_{\rm LL}\);
\(\vec d^{\rm nf}_{\rm LL}(y^+_{\rm LL},y^-_{\rm LL})\) & the two real
coordinates of the plane of the normal form and its fields,
\(\vec d^{\rm nf}_{\rm LL}(y^+_{\rm LL},y^-_{\rm LL})=(0,\tfrac i2(y^+_{\rm LL}-y^-_{\rm LL}),\tfrac12(y^+_{\rm LL}+y^-_{\rm LL}))\),
on which
\(\Omega_{\rm LL}=\Omega^{\rm s}_{\rm LL}(y^+_{\rm LL})+\Omega^{\rm s}_{\rm LL}(y^-_{\rm LL})\);
for \(y^+_{\rm LL}\ge y^-_{\rm LL}\ge0\) the normal form with the gaps
\(y^\pm_{\rm LL}\) & new & \texttt{y\_+} \texttt{y\_-} \texttt{d\_nf} \\
\(t^{\rm fl}_{\rm LL}\) & the parameter of the gradient flow of the
effective potential (a fictitious time): along the flow the six real
coordinates of \(\vec d_{\rm LL}\) change with the rates minus the
derivatives of \(\Omega_{\rm LL}\) with respect to them & new & \\
\(\kappa^{\rm st}_{\rm LL}\) & a real parameter along a straight line of
fields, used in the proof of Theorem 11.2 of the document of Test 2 &
new & \\
\(\mathcal V^{\rm vac}_{\rm LL}\) & the vacuum manifold: the set of the
minima of the effective potential, the polar fields of modulus
\(\Delta^{\rm vac}_{\rm LL}\),
\(\exp(2i\theta^{\rm F}_{\rm LL})\mathsf R^{\rm v}_{\rm LL}(\mathsf U_{\rm LL})(0,0,\Delta^{\rm vac}_{\rm LL})\);
a manifold of dimension 3 & new & \texttt{V\_vac} \\
\(\theta^\iota_{\rm LL}\) & the three angles of an isospin rotation,
\(\mathsf U_{\rm LL}=\exp(i\sum_\iota\theta^\iota_{\rm LL}\tau^{\rm LL}_\iota/2)\);
\(\theta^3_{\rm LL}\) alone gives the rotations
\(\exp(i\theta^3_{\rm LL}Q_{\rm LL})\) of the electric charge (plain
text \texttt{al} in the report of check SB-14) & new & \texttt{th\^{}1}
\texttt{th\^{}2} \texttt{th\^{}3} \texttt{th\^{}iota} \texttt{al} \\
\(\vec d^{\rm tr}_{\rm LL}\) & the first-order change of the field
\(\vec d_{\rm LL}\) under the transformation (9.4) with the small angles
\(\theta^{\rm F}_{\rm LL}\), \(\theta^\iota_{\rm LL}\):
\(2i\theta^{\rm F}_{\rm LL}\vec d_{\rm LL}-(\theta^1_{\rm LL},\theta^2_{\rm LL},\theta^3_{\rm LL})\times\vec d_{\rm LL}\)
& new & \\
\(H^{\rm vac}_{\rm LL}\); \(U(1)^{\rm Q}_{\rm LL}\) & the stabilizer of
the vacuum \((0,0,\Delta^{\rm vac}_{\rm LL})\): the transformations
(9.4) that leave it unchanged, a group \(O(2)\) formed by the group
\(U(1)^{\rm Q}_{\rm LL}\) of the electric charge (the rotations
\(\exp(i\theta^3_{\rm LL}Q_{\rm LL})\) about the third isospin axis) and
the reflection \(\theta^{\rm F}_{\rm LL}=\pi/2\),
\(\mathsf U_{\rm LL}=i\tau^{\rm LL}_1\), which reverses the charge & new
& \texttt{H\_vac} \texttt{U(1)\^{}Q} \\
\(SO(3)\), \(O(2)\) & the group of the rotations of three-dimensional
space (here: of the isospin space; on the field \(\vec d_{\rm LL}\) the
transformations (9.4) act as the phases times \(SO(3)\)), and the group
of the rotations and reflections of a plane & new & \texttt{SO(3)}
\texttt{O(2)} \\
\(\vartheta^{\rm G}_{{\rm LL},0}\), \(\vartheta^{\rm G}_{{\rm LL},1}\),
\(\vartheta^{\rm G}_{{\rm LL},2}\); \(\vartheta^{\rm G}_{\rm LL}\),
\(\varphi^{\rm G}_{\rm LL}\) & the three Goldstone fields of SFBSv: the
phase \(\vartheta^{\rm G}_{{\rm LL},0}\) (the angle of the broken
fermion number) and the angles \(\vartheta^{\rm G}_{{\rm LL},1}\),
\(\vartheta^{\rm G}_{{\rm LL},2}\) of the broken isospin rotations about
the first two axes;
\(\vartheta^{\rm G}_{{\rm LL},1}+i\vartheta^{\rm G}_{{\rm LL},2}=\vartheta^{\rm G}_{\rm LL}\exp(i\varphi^{\rm G}_{\rm LL})\)
with \(\vartheta^{\rm G}_{\rm LL}\ge0\) & new & \texttt{thG0}
\texttt{thG1} \texttt{thG2} \texttt{thG} \texttt{phiG} \\
\(\mathsf U^{\rm G}_{\rm LL}\) & the isospin rotation of the Goldstone
fields of SFBSv,
\(\exp\big(\tfrac i2(\vartheta^{\rm G}_{{\rm LL},1}\tau^{\rm LL}_1+\vartheta^{\rm G}_{{\rm LL},2}\tau^{\rm LL}_2)\big)\),
a matrix of \(SU(2)\) whose rotation turns the third isospin axis by the
angle \(\vartheta^{\rm G}_{\rm LL}\) & new & \texttt{UG}
\texttt{U(thG1,thG2)} \\
\(h^{\rm a}_{\rm LL}\) & the amplitude field of SFBSv: the change of the
length of the real part of the field from \(\Delta^{\rm vac}_{\rm LL}\),
the analogue of the Higgs field; electrically neutral & new &
\texttt{h\_a} \\
\(c^{\rm n}_{{\rm LL},1}\), \(c^{\rm n}_{{\rm LL},2}\);
\(\varrho^{\rm n}_{\rm LL}\) & the two non-unitary fields of SFBSv, the
components of an imaginary part orthogonal to the real part of the field
(\(c^{\rm n}_{{\rm LL},1}\mp i\,c^{\rm n}_{{\rm LL},2}\) carry the
electric charge \(\pm1\)), and their modulus
\(\varrho^{\rm n}_{\rm LL}=((c^{\rm n}_{{\rm LL},1})^2+(c^{\rm n}_{{\rm LL},2})^2)^{1/2}\)
& new & \texttt{c\_n1} \texttt{c\_n2} \texttt{rho\_n} \\
\(\Omega^{\rm H}_{\rm LL}(\vec d_{\rm LL})\),
\(\lambda^{\rm H}_{\rm LL}\) & the potential of the Higgs form with the
minimum on \(\mathcal V^{\rm vac}_{\rm LL}\),
\(\lambda^{\rm H}_{\rm LL}(\lvert\vec d_{\rm LL}\rvert^2-(\Delta^{\rm vac}_{\rm LL})^2)^2\),
used for comparison; its coupling
\(\lambda^{\rm H}_{\rm LL}=\Omega^{\rm s(2)}_{\rm LL}/(4(\Delta^{\rm vac}_{\rm LL})^2)\)
gives it the curvature of the computed potential along
\(h^{\rm a}_{\rm LL}\) & new & \texttt{Omega\_H} \texttt{lambda\_H} \\
\(\mathcal L^{\rm kin}_{\rm LL}\) & the kinetic Lagrangian of the SFBS,
the part of \(\mathcal L^{\rm B}_{\rm LL}\) that contains the
derivatives,
\(Z^{\rm t}_{\rm LL}\sum_\iota\lvert(\nabla^{\rm B}_{{\rm LL},0}\vec d_{\rm LL})^\iota\rvert^2-Z^{\rm s}_{\rm LL}\sum_j\sum_\iota\lvert(\nabla^{\rm B}_{{\rm LL},j}\vec d_{\rm LL})^\iota\rvert^2\),
so that
\(\mathcal L^{\rm B}_{\rm LL}=\mathcal L^{\rm kin}_{\rm LL}-\Omega_{\rm LL}(\vec d_{\rm LL})\)
& new & \texttt{Lagr\_kin} \texttt{L\_kin} \\
\(\mathsf N^{\rm g}_{\rm LL}(\vec d_{\rm LL})\),
\(\mathsf N^{\rm v}_{\rm LL}(\vec d_{\rm LL})\);
\(\vec w^{\rm n}_{\rm LL}\) & the contact matrix: the real symmetric
\(4\times4\) matrix of the terms of
\(\sum_\iota\lvert(\nabla^{\rm B}_{{\rm LL},p}\vec d_{\rm LL})^\iota\rvert^2\)
that are quadratic in the potentials
\((A^{\rm F}_{{\rm LL},p},A^1_{{\rm LL},p},A^2_{{\rm LL},p},A^3_{{\rm LL},p})\),
with the corner entry
\(4(g^{\rm F}_{\rm LL})^2\lvert\vec d_{\rm LL}\rvert^2\), the first row
and column
\(4g^{\rm F}_{\rm LL}g^{\rm v}_{\rm LL}(w^{\rm n}_{\rm LL})^\iota\) and
the isospin block \((g^{\rm v}_{\rm LL})^2\mathsf N^{\rm v}_{\rm LL}\),
whose entries are
\(\lvert\vec d_{\rm LL}\rvert^2-\lvert d^\iota_{\rm LL}\rvert^2\) on the
diagonal and \(-{\rm Re}(d^\iota_{\rm LL}(d^{\iota'}_{\rm LL})^*)\) off
it; and the real vector
\(\vec w^{\rm n}_{\rm LL}={\rm Re}\,\vec d_{\rm LL}\times{\rm Im}\,\vec d_{\rm LL}\),
zero for the polar fields & new & \texttt{N\_g} \texttt{N\_v}
\texttt{w\_n} \\
\(E^{\rm c}_{\rm LL}\) & the quasi-particle energy of the condensed
vacuum \((0,0,\Delta_{\rm LL})\) at \(\mu_{\rm LL}=0\),
\(E^{\rm c}_{\rm LL}=(q_{\rm LL}^2+(M^{\rm q}_{\rm LL})^2)^{1/2}\), an
eigenvalue of \(\mathsf H_{\rm LL}(\vec q_{\rm LL})\) eight times with
each sign; \(E^{\rm c}_{\rm LL}(\Lambda_{\rm LL})\) is its value at
\(q_{\rm LL}=\Lambda_{\rm LL}\) (plain text \texttt{E\_Lambda}) & new &
\texttt{E\_c} \texttt{E\_Lambda} \\
\(\vec a^{\rm F}_{\rm LL}\) & the fermion-number coupling times the
space components of a constant fermion-number potential,
\(\vec a^{\rm F}_{\rm LL}=g^{\rm F}_{\rm LL}(A^{\rm F}_{{\rm LL},1},A^{\rm F}_{{\rm LL},2},A^{\rm F}_{{\rm LL},3})\):
the shift of the kinetic momentum of both members of a pair & new &
\texttt{a\_F} \\
\(\vec u^{\rm f}_{\rm LL}\), \(\vec u^{\rm f'}_{\rm LL}\);
\(\delta^{\rm f}_{\rm LL}\) & a fixed complex unit vector, the direction
of a deviation
\(\delta^{\rm f}_{\rm LL}(x_{\rm LL})\,\vec u^{\rm f}_{\rm LL}\) of the
field from the vacuum \((0,0,\Delta_{\rm LL})\), one of the six
directions \((1,0,0)\), \((0,1,0)\), \((0,0,i)\) (Goldstone) and
\((i,0,0)\), \((0,i,0)\), \((0,0,1)\) (massive); a second such
direction; and the small real amplitude of the deviation, a function of
space and time & new & \texttt{u\_f} \texttt{delta\_f} \\
\(\mathsf H^{\rm f}_{\rm LL}\), \(\mathsf H^{\rm f'}_{\rm LL}\) & the
change of the \(16\times16\) Bogoliubov matrix per unit of
\(\delta^{\rm f}_{\rm LL}\),
\(\begin{pmatrix}0&\mathsf P_{\rm LL}(\vec u^{\rm f}_{\rm LL})\\\mathsf P_{\rm LL}(\vec u^{\rm f}_{\rm LL})^\dagger&0\end{pmatrix}\),
where \(\mathsf P_{\rm LL}(\vec u^{\rm f}_{\rm LL})\) is the pairing
block with \(\vec u^{\rm f}_{\rm LL}\) in place of \(\vec d_{\rm LL}\);
and the same for \(\vec u^{\rm f'}_{\rm LL}\) & new & \texttt{H\_f} \\
\(\chi^{\rm f}_{\rm LL}(\omega^{\rm B}_{\rm LL})\),
\(\Gamma^{\rm f}_{\rm LL}(\omega^{\rm B}_{\rm LL})\);
\(\Gamma^{\rm G}_{\rm LL}\), \(\Gamma^{\rm m}_{\rm LL}\) & the pair
susceptibility of a direction \(\vec u^{\rm f}_{\rm LL}\) of the
condensed vacuum at the frequency \(\omega^{\rm B}_{\rm LL}\) and zero
wave number, and the kernel of the direction,
\(\Gamma^{\rm f}_{\rm LL}=\lvert\vec u^{\rm f}_{\rm LL}\rvert^2/G_{\rm LL}-\chi^{\rm f}_{\rm LL}\):
minus the coefficient of the time average of
\((\delta^{\rm f}_{\rm LL})^2\) in the effective action per duration and
volume (plain text \texttt{Gamma\_bdg}: the kernel computed from the
Bogoliubov matrices); the kernels of the three Goldstone directions and
of the three massive directions,
\(\Gamma^{\rm G}_{\rm LL}=-(\omega^{\rm B}_{\rm LL})^2\Upsilon_{\rm LL}\)
and
\(\Gamma^{\rm m}_{\rm LL}=(4\Delta_{\rm LL}^2-(\omega^{\rm B}_{\rm LL})^2)\Upsilon_{\rm LL}\)
& new & \texttt{chi\_f} \texttt{Gamma\_f} \texttt{Gamma\_bdg}
\texttt{Gamma\_G} \texttt{Gamma\_m} \\
\(\Upsilon_{\rm LL}(\omega^{\rm B}_{\rm LL})\),
\(\Upsilon^{(1)}_{\rm LL}\); \(c^{\rm w}_{\rm LL}\) & the common factor
of the kernels of the condensed vacuum,
\(\Upsilon_{\rm LL}=2\int_{q_{\rm LL}<\Lambda_{\rm LL}}d^3q_{\rm LL}(2\pi)^{-3}\big(E^{\rm c}_{\rm LL}(4(E^{\rm c}_{\rm LL})^2-(\omega^{\rm B}_{\rm LL})^2)\big)^{-1}\),
with \(\Upsilon_{\rm LL}(0)=Z^{\rm t}_{\rm LL}(\Delta_{\rm LL})\); its
derivative with respect to \((\omega^{\rm B}_{\rm LL})^2\) at zero; and
the positive number
\(c^{\rm w}_{\rm LL}=(M^{\rm q}_{\rm LL})^2-(\omega^{\rm B}_{\rm LL})^2/4\)
of the proof of Theorem 12.2 of the document of Test 2 & new &
\texttt{Upsilon} \texttt{Upsilon\_1} \\
\(\vec d^{\rm u}_{\rm LL}\) & the SFBS field in the unitary gauge, the
column of the SFBSv parametrization,
\((i\,c^{\rm n}_{{\rm LL},1},\,i\,c^{\rm n}_{{\rm LL},2},\,\Delta^{\rm vac}_{\rm LL}+h^{\rm a}_{\rm LL})\):
the field after the gauge transformation that removes the three
Goldstone fields & new & \\
\(V^{\rm F}_{{\rm LL},p}\), \(V^\iota_{{\rm LL},p}\) & the gauge
potentials in the unitary gauge, the vector fields of definite mass:
\(V^{\rm F}_{{\rm LL},p}=A^{\rm F}_{{\rm LL},p}-\partial_p\vartheta^{\rm G}_{{\rm LL},0}/g^{\rm F}_{\rm LL}\)
and
\(\sum_\iota V^\iota_{{\rm LL},p}\tau^{\rm LL}_\iota/2=(\mathsf U^{\rm G}_{\rm LL})^\dagger\mathsf a_{{\rm LL},p}\mathsf U^{\rm G}_{\rm LL}-(i/g^{\rm v}_{\rm LL})(\partial_p(\mathsf U^{\rm G}_{\rm LL})^\dagger)\mathsf U^{\rm G}_{\rm LL}\);
to first order
\(V^\iota_{{\rm LL},p}=A^\iota_{{\rm LL},p}-\partial_p\vartheta^{\rm G}_{{\rm LL},\iota}/g^{\rm v}_{\rm LL}\)
(\(\iota=1,2\)) and \(V^3_{{\rm LL},p}=A^3_{{\rm LL},p}\), the photon &
new & \texttt{V} \texttt{V\^{}iota} \texttt{VF} \\
\(W^\pm_{{\rm LL},p}\) & the charged massive vector fields,
\((V^1_{{\rm LL},p}\mp iV^2_{{\rm LL},p})/\sqrt2\), of the electric
charges \(\pm1\) & new & \texttt{W\_+} \texttt{W\_-} \\
\(\mathsf M^{\rm t}_{\rm LL}\), \(\mathsf M^{\rm s}_{\rm LL}\) & the
mass matrices of the time and of the space components of the four vector
fields
\((V^{\rm F}_{{\rm LL},p},V^1_{{\rm LL},p},V^2_{{\rm LL},p},V^3_{{\rm LL},p})\)
in the vacuum: \(2Z^{\rm t}_{\rm LL}\) and \(2Z^{\rm s}_{\rm LL}\) times
the contact matrix of the vacuum,
\({\rm diag}(4(g^{\rm F}_{\rm LL})^2,(g^{\rm v}_{\rm LL})^2,(g^{\rm v}_{\rm LL})^2,0)(\Delta^{\rm vac}_{\rm LL})^2\);
the quadratic Lagrangian contains one half of the first and minus one
half of the second, each between the row and the column of the
components & new & \\
\(M^{\rm W}_{\rm LL}\), \(M^{\rm F}_{\rm LL}\) & the rest masses of the
charged vector fields \(W^\pm_{{\rm LL},p}\) and of the fermion-number
vector field \(V^{\rm F}_{{\rm LL},p}\),
\((M^{\rm W}_{\rm LL})^2=2Z^{\rm s}_{\rm LL}(g^{\rm v}_{\rm LL})^2(\Delta^{\rm vac}_{\rm LL})^2\)
and
\((M^{\rm F}_{\rm LL})^2=8Z^{\rm s}_{\rm LL}(g^{\rm F}_{\rm LL})^2(\Delta^{\rm vac}_{\rm LL})^2\)
(plain text \texttt{M\_W\_time}, \texttt{M\_F\_time}: the masses of the
time components, with \(Z^{\rm t}_{\rm LL}\) in place of
\(Z^{\rm s}_{\rm LL}\)) & new & \texttt{M\_W} \texttt{M\_F}
\texttt{M\_W\_time} \texttt{M\_F\_time} \\
\(\theta^{\rm mix}_{\rm LL}\) & the mixing angle: the angle between the
massless gauge field (the null vector of the mass matrix) and the axis
of the potential \(A^3_{{\rm LL},p}\) in the plane of
\(A^{\rm F}_{{\rm LL},p}\) and \(A^3_{{\rm LL},p}\); zero in the vacuum,
\(\arctan(g^{\rm v}_{\rm LL}/(2g^{\rm F}_{\rm LL}))\) at the non-unitary
saddle point (plain text: the values found at the vacuum and at the
saddle point, and the prediction) & new & \texttt{theta\_mix}
\texttt{theta\_mix\_vacuum} \texttt{theta\_mix\_saddle}
\texttt{theta\_mix\_pred} \\
\(\mathcal L^{(2)}_{\rm LL}\) & the part of the second order, in the
fields of SFBSv and the gauge potentials, of the Lagrangian (13.1) of
Test 2 (the effective Lagrangian of the SFBS with the Yang-Mills and
Maxwell terms) & new & \\
\(w^{\rm lat}_{\rm LL}\), \(w^{\rm cmp}_{\rm LL}\);
\(n^{\rm cpx}_{\rm LL}\) & a Nambu mode on the lattice written as one
complex vector, the 16 components of \(\psi^{\rm p}_{\rm LL}\) and
\(\psi^{\rm h}_{\rm LL}\) on every site, the sites one after the other,
so that the lattice equations read
\(i\,\partial_0w^{\rm lat}_{\rm LL}=\mathcal H^{\rm lat}_{\rm LL}w^{\rm lat}_{\rm LL}\)
(for the self-consistent evolution the columns of
\(\mathsf W^{\rm m}_{\rm LL}\) one after the other); one complex
component of it; and the number of its complex components, \(16\) times
the number of sites (times the number of modes) & new & \texttt{w\_lat}
\texttt{w\_cmp} \texttt{n\_complex} \\
\(r^{\rm tol}_{\rm LL}\), \(a^{\rm tol}_{\rm LL}\) & the relative and
the absolute tolerance of the time integration by CVODE: the estimated
local error of every internal step is accepted only if its weighted
root-mean-square norm is at most 1, with the weight
\(1/(r^{\rm tol}_{\rm LL}\lvert w^{\rm cmp}_{\rm LL}\rvert+a^{\rm tol}_{\rm LL})\)
for the real and for the imaginary part of every complex component
\(w^{\rm cmp}_{\rm LL}\) of the state & new & \texttt{rtol}
\texttt{atol} \\
\(h^{\rm CV}_{\rm LL}\), \(q^{\rm CV}_{\rm LL}\);
\(N^{\rm st}_{\rm LL}\), \(N^{\rm ev}_{\rm LL}\) & the size and the
order of an internal step of CVODE (both chosen by CVODE, the order from
1 to 12 for the Adams-Moulton formulas); the number of internal steps of
an integration and the number of evaluations of its right side
\(-i\mathcal H^{\rm lat}_{\rm LL}w^{\rm lat}_{\rm LL}\) (the cost; one
step of the classical Runge-Kutta method of fourth order needs 4
evaluations) & new & \texttt{h\_CV} \texttt{q\_CV} \texttt{N\_steps}
\texttt{N\_evals} \\
\(B^{\rm err}_{\rm LL}\); \(\delta^{\rm rel}_{\rm LL}\),
\(w^{\rm ex}_{\rm LL}\) & the accumulated local error that the error
test of CVODE permits,
\(B^{\rm err}_{\rm LL}=N^{\rm st}_{\rm LL}\sqrt2\,(r^{\rm tol}_{\rm LL}\lVert w^{\rm lat}_{\rm LL}\rVert+a^{\rm tol}_{\rm LL}\sqrt{n^{\rm cpx}_{\rm LL}})\)
for a state of the norm \(\lVert w^{\rm lat}_{\rm LL}\rVert\), the
expected size of the global error of a unitary evolution; and the
relative deviation of a computed state from the exact solution,
\(\lVert w^{\rm lat}_{\rm LL}-w^{\rm ex}_{\rm LL}\rVert/\lVert w^{\rm ex}_{\rm LL}\rVert\),
with the exact solution \(w^{\rm ex}_{\rm LL}\) of the same equations &
new & \texttt{budget} \texttt{max\_deviation} \texttt{w\_exact} \\
\end{longtable}

\end{document}
